%ID VD_3_1_2_01
\textbf{Bước 1. Tách li độ và vận tốc.} Với $a=1$ ở (i), (iii), (iv), và $a=2$ ở (ii), dùng
\[u=\tfrac12[f(x-at)+f(x+at)]+\frac{H(x+at)-H(x-at)}{2a},\qquad H'=g.\]
\textbf{Bước 2. Tính từng trường hợp.}

(i) $g=0$, nên $H=0$. Công thức $\sin(x-t)+\sin(x+t)=2\sin x\cos t$ cho $u=\sin x\cos t$.

(ii) Chỉ đổi độ dịch thành $2t$: $u=\sin x\cos2t$. Hai chu kỳ cơ bản lần lượt là $2\pi$ và $\pi$.

(iii) $f=0$, $g(s)=s|s|$. Trên $s>0$, nguyên hàm là $s^3/3$; trên $s<0$, là $-s^3/3$. Nối tại $0$ được $H(s)=|s|^3/3$, do đó
\[u=\frac{|x+t|^3-|x-t|^3}{6}.\]
Hệ số $1/2$ ở nguyên hàm trong bản scan không đúng: $(|s|^3/3)'=s|s|$. Hàm $g$ thuộc $C^1$, nên nghiệm thuộc $C^2$, kể cả trên $x\pm t=0$; đạo hàm cấp ba có thể nhảy.

(iv) $g(s)=|s|$ có nguyên hàm $H=s|s|/2$, nên
\[u=\frac{(x+t)|x+t|-(x-t)|x-t|}{4}.\]
Vì $H\in C^1$ nhưng không $C^2$ tại $0$, nghiệm $C^1$ và không $C^2$ trên hai đặc trưng $x=\pm t$.

\textbf{Bước 3. Kiểm tra.} Tại $t=0$, hai hiệu nguyên hàm bằng không; đạo hàm theo $t$ cho $g(x)$. Mỗi hạng chỉ phụ thuộc $x\pm at$ nên thỏa PDE nơi đạo hàm tồn tại. Hai nghiệm (iii), (iv) tăng theo $t$ ở những điểm thích hợp, nên không tuần hoàn.

%ID VD_3_1_2_02
\textbf{Bước 1. Giải phần có nguồn với dữ liệu không.} Nguồn $F(x,t)=xt$, $a=1$. Tích phân trong của Duhamel là
\[\int_{x-(t-s)}^{x+(t-s)}ys\,dy=2xs(t-s).\]
Suy $u_F=x\int_0^t s(t-s)ds=xt^3/6$.

\textbf{Bước 2. Thêm phần thuần nhất.} Ở (i), $f=g=0$, nên $u=xt^3/6$. Ở (ii), $f=\sin x$, $g=x|x|$, nguyên hàm đúng là $|x|^3/3$:
\[u=\sin x\cos t+\frac{|x+t|^3-|x-t|^3}{6}+\frac{xt^3}{6}.\]
\textbf{Bước 3. Kiểm tra.} $(u_F)_{tt}=xt$, $(u_F)_{xx}=0$, $u_F(x,0)=(u_F)_t(x,0)=0$. Phần còn lại thỏa phương trình thuần nhất và cho đúng $f,g$. Dữ liệu $g\in C^1$ bảo đảm nghiệm $C^2$; không cần loại hai đường $x=\pm t$ ở mức đạo hàm cấp hai.

\textbf{Chi tiết hai tích phân nguồn.} Đặt $h=t-s$. Tích phân trong bằng
\[s\left[\frac{y^2}{2}\right]_{x-h}^{x+h}
=\frac s2[(x+h)^2-(x-h)^2]=2xsh.\]
Nhân $1/2$ của Duhamel, tích phân ngoài là
\[x\int_0^t(ts-s^2)ds=x\left[\frac{ts^2}{2}-\frac{s^3}{3}\right]_0^t
=x(t^3/2-t^3/3)=xt^3/6.\]

%ID VD_3_1_2_03
\textbf{Bước 1. Đặt tích phân trong.} Viết $V=(2a)^{-1}\int_0^t G(x,t,s)ds$, với
\[G=\int_{x-a(t-s)}^{x+a(t-s)}F(y,s)dy.\]
Ở đây $s$ là tham số khi lấy đạo hàm theo $t$.

\textbf{Bước 2. Đạo hàm các cận.} Với $F_\pm=F(x\pm a(t-s),s)$,
\[G_t=a(F_++F_-),\quad G_x=F_+-F_-,\quad G_{tt}=a^2(F_{x,+}-F_{x,-})=a^2G_{xx}.\]
Khi $s=t$, $G=0$ và $G_t=2aF(x,t)$.

\textbf{Bước 3. Đạo hàm tích phân ngoài.} Quy tắc Leibniz cho
\[V_t=\frac1{2a}\int_0^tG_tds,\qquad V_{tt}=F(x,t)+\frac1{2a}\int_0^tG_{tt}ds.\]
Do đó $V_{tt}-a^2V_{xx}=F$. Tại $t=0$ cả $V,V_t$ bằng không. Các đạo hàm trên hợp lệ nếu $F$ có các đạo hàm liên tục cần dùng (chẳng hạn $F\in C^1$ theo không gian và liên tục theo thời gian). Thêm nghiệm d'Alembert của dữ liệu $f,g$ được nghiệm đầy đủ.

%ID VD_3_1_2_04
\textbf{Bước 1. Xác định hai miền.} Với biên $u(0,t)=h(t)$, ở $x\ge at$ dùng dữ liệu đầu; ở $x<at$ thêm tín hiệu biên tại thời điểm trễ $t-x/a$ và lấy hiệu hai giá trị $f$.

\textbf{Bước 2. Tính cụ thể.}

(i) $a=1$, $f=g=0$, $h=t^3$. Trước khi sóng biên đến, $u=0$; sau đó $u=(t-x)^3$. Gộp lại $u=(t-x)_+^3$. Hàm này $C^2$ trên đường $x=t$.

(ii) $a=2$, $f=\sin x$, $g=1$, $h=0$. Ngoài mặt sóng, phần vận tốc là $t$. Trong mặt sóng, đoạn tích phân $[2t-x,2t+x]$ dài $2x$, cho $x/2$. Phần li độ ở cả hai miền bằng $\sin x\cos2t$. Vì vậy
\[u=\sin x\cos2t+\min(t,x/2).\]
Nghiệm liên tục nhưng không $C^1$ trên $x=2t$: vận tốc đầu $g(0)=1$ trái với $h'(0)=0$.

(iii) $a=1$, $f=x^2$, $g=6x$, $h=t^2$. Ngoài mặt sóng, half-sum bằng $x^2+t^2$, tích phân vận tốc bằng $6xt$. Trong mặt sóng,
\[(t-x)^2+\tfrac12[(x+t)^2-(t-x)^2]+\tfrac12\int_{t-x}^{t+x}6sds=x^2+t^2+6xt.\]
\textbf{Bước 3. Kiểm tra.} Nghiệm (iii) có $u_{tt}=u_{xx}=2$, $u(x,0)=x^2$, $u_t(x,0)=6x$, $u(0,t)=t^2$. Ba điều kiện tương thích tại góc đều đúng.

\textbf{Chi tiết rút gọn câu (iii).} Ngoài mặt sóng, half-sum là $[(x-t)^2+(x+t)^2]/2=x^2+t^2$, còn tích phân là $3[(x+t)^2-(x-t)^2]/2=6xt$. Trong mặt sóng, half-difference là $[(x+t)^2-(t-x)^2]/2=2xt$, tích phân vận tốc là $3[(t+x)^2-(t-x)^2]/2=6xt$. Cộng thêm biên $(t-x)^2=x^2-2xt+t^2$ cho cùng $x^2+t^2+6xt$.

%ID VD_3_1_2_05
\textbf{Bước 1. Điều kiện cần ở góc.} Nếu $u\in C^2$ đến $(0,0)$ thì $u(0,0)=f(0)=h(0)$, $u_t(0,0)=g(0)=h'(0)$. Từ PDE ở góc, $h''(0)=u_{tt}(0,0)=a^2f''(0)$.

\textbf{Bước 2. So sánh hai công thức trên $x=at$.} Đặt $s=x+at$, $r=at-x$. Công thức trong là $h(t-x/a)+(f(s)-f(r))/2+(2a)^{-1}\int_r^s g$. Ngoài là half-sum $f(x-at),f(s)$ cộng tích phân $g$. Tại $r=0$, chênh lệch vết trong trừ ngoài của $u$ là $h(0)-f(0)$.

\textbf{Bước 3. Đạo hàm cấp một.} Trong miền phản xạ,
\[u_x=-h'(t-x/a)/a+\tfrac12[f'(s)+f'(r)]+\frac{g(s)+g(r)}{2a}.\]
Ngoài miền phản xạ, $u_x=(f'(s)+f'(x-at))/2+(g(s)-g(x-at))/(2a)$. Hiệu vết là $(g(0)-h'(0))/a$; tương tự hiệu vết $u_t$ là $h'(0)-g(0)$.

\textbf{Bước 4. Đạo hàm cấp hai.} Hiệu vết $u_{xx}$ bằng $h''(0)/a^2-f''(0)$, vì các hạng $g'$ trùng nhau. Hiệu $u_{xt}$ bằng $af''(0)-h''(0)/a$; $u_{tt}$ bằng $h''(0)-a^2f''(0)$. Như vậy ba điều kiện đã nêu đủ để ghép thành $C^2$, nếu $f\in C^2,g\in C^1,h\in C^2$. Nghiệm thỏa dữ liệu và PDE theo công thức; tính duy nhất suy từ hai sóng chạy hoặc năng lượng của hiệu hai nghiệm.

\textbf{Chi tiết so các vết cấp hai.} Trong miền phản xạ,
\[\begin{aligned}
u_{xx}&=h''(\tau)/a^2+(f''(s)-f''(r))/2+(g'(s)-g'(r))/(2a),\\
u_{xt}&=-h''(\tau)/a+a(f''(s)+f''(r))/2+(g'(s)+g'(r))/2,\\
u_{tt}&=h''(\tau)+a^2(f''(s)-f''(r))/2+a(g'(s)-g'(r))/2.
\end{aligned}\]
Ngoài miền phản xạ, viết $\xi=x-at$, $\eta=x+at$:
\[\begin{aligned}
u_{xx}&=(f''(\xi)+f''(\eta))/2+(g'(\eta)-g'(\xi))/(2a),\\
u_{xt}&=a(f''(\eta)-f''(\xi))/2+(g'(\eta)+g'(\xi))/2,\\
u_{tt}&=a^2(f''(\xi)+f''(\eta))/2+a(g'(\eta)-g'(\xi))/2.
\end{aligned}\]
Trên $x=at$, $r=\xi=\tau=0$, $s=\eta=2at$. Trừ từng cặp: các hạng tại $2at$ và mọi hạng $g'$ mất, còn lần lượt $h''(0)/a^2-f''(0)$, $af''(0)-h''(0)/a$, $h''(0)-a^2f''(0)$.

%ID BT_3_1_3_01
\textbf{Bước 1. Miền hỗ trợ.} Nếu $|x|>M+at$, hai cận $x\pm at$ và toàn đoạn giữa chúng cùng nằm ngoài $[-M,M]$. D'Alembert cho $u=0$; các đạo hàm cũng bằng không ngoài miền ấy khi dữ liệu đủ trơn.

\textbf{Bước 2. Đạo hàm năng lượng.} Đặt $I(t)=\int_{\mathbb R}(a^2u_x^2+u_t^2)dx$. Trên mỗi khoảng thời gian hữu hạn có thể tích phân trên một đoạn cố định đủ lớn, vì sóng có hỗ trợ hữu hạn. Dùng PDE:
\[I'=2a^2\int(u_xu_{xt}+u_tu_{xx})dx=2a^2[u_tu_x]_{-\infty}^{+\infty}=0.\]
\textbf{Bước 3. Giá trị đầu.} Tại $t=0$, $u_x=f'$, $u_t=g$, nên
\[I(t)=I(0)=\int_{\mathbb R}(a^2(f')^2+g^2)dx.\]
Điều kiện hỗ trợ bảo đảm hạng biên bằng không, không phải chỉ vì ký hiệu cận là vô hạn. Với dữ liệu năng lượng, kết quả tiếp tục đúng bằng xấp xỉ dữ liệu trơn.

\textbf{Chi tiết tích phân từng phần trong đạo hàm năng lượng.} Đặt $E=\frac12\int_{-R}^R(u_t^2+a^2u_x^2+\kappa u^2)dx$. Với $u_{tt}-a^2u_{xx}+\gamma u_t+\kappa u=F$,
\[\begin{aligned}
E'&=\int_{-R}^R(u_tu_{tt}+a^2u_xu_{xt}+\kappa uu_t)dx\\
&=\int_{-R}^Ru_t(a^2u_{xx}-\gamma u_t-\kappa u+F)dx
+a^2\int_{-R}^Ru_xu_{xt}dx+\kappa\int_{-R}^Ruu_tdx\\
&=a^2[u_tu_x]_{-R}^R-a^2\int_{-R}^Ru_{tx}u_xdx
+a^2\int_{-R}^Ru_xu_{xt}dx-\gamma\int_{-R}^Ru_t^2dx+\int_{-R}^RFu_tdx\\
&=a^2[u_tu_x]_{-R}^R-\gamma\int_{-R}^Ru_t^2dx+\int_{-R}^RFu_tdx.
\end{aligned}\]
Hai hạng $\kappa uu_t$ triệt tiêu ngay khi thế PDE; hai tích phân $u_xu_{xt}$ triệt tiêu sau tích phân từng phần. Chọn $R$ lớn hơn miền hỗ trợ trên khoảng thời gian xét, nên $u_t=u_x=0$ ở $x=\pm R$. Với $I=2E$, nhân đẳng thức đạo hàm với $2$ để có $I'=0$.

%ID BT_3_1_3_02
\textbf{Bước 1. Tính sóng bằng nguyên hàm.} Dùng $H'=g$ và half-sum của $f$. Công thức tổng quát trong phần kiến thức cho cách kiểm tra $u_t(x,0)=g(x)$; dưới đây ghi rõ nguyên hàm và kết quả từng câu.

(i) $a=3$, $H=-\cos2s/2$. Hiệu hai cos bằng $2\sin2x\sin6t$, nên
\[u=\cos x\cos3t+\tfrac16\sin2x\sin6t,\qquad T=2\pi/3.\]
(ii) $a=5$, $H=s^2/2$. Hiệu bình phương bằng $20xt$, cho $u=\sin3x\cos15t+xt$, không tuần hoàn.

(iii) $a=2$, $H=\sin s$. Half-sum của bình phương bằng $x^2+4t^2$, nên $u=x^2+4t^2+\cos x\sin2t/2$, không tuần hoàn.

(iv) $a=6$, $H=s/2+\sin2s/4$ vì $\cos^2s=(1+\cos2s)/2$:
\[u=\tfrac12(e^{-|x-6t|}+e^{-|x+6t|})+t/2+\cos2x\sin12t/24.\]
Hạng $t/2$ ngăn tuần hoàn; li độ có góc tại $0$ nên không $C^1$ trên $x=\pm6t$.

(v) $a=1$, $g=(3\sin s-\sin3s)/4$, $H=-3\cos s/4+\cos3s/12$. Suy
\[u=\sin x\cos t-\tfrac12-\tfrac12\cos2x\cos2t+\tfrac34\sin x\sin t-\tfrac1{12}\sin3x\sin3t.\]
Mode tần số $1$ khác không, nên chu kỳ cơ bản là $2\pi$.

(vi) $a=7$, $f=1$. Nguyên hàm của xung tam giác được nối liên tục:
\[H(s)=\begin{cases}0&s\le-1,\\(s+1)^2/2&-1<s\le0,\\1-(1-s)^2/2&0<s<1,\\1&s\ge1.\end{cases}\quad
u=1+\frac{H(x+7t)-H(x-7t)}{14}.\]
Tại mỗi $x$ cố định, sau khi cả hai cận ra ngoài hỗ trợ, $u=15/14$, khác giá trị đầu $1$; vì thế không tuần hoàn. Nghiệm $C^1$, các đạo hàm cấp hai gãy trên $x\pm7t=-1,0,1$.

\textbf{Bước 2. Dựng đồ thị.} Sóng phải $A=f/2-H/(2a)$ được dịch sang phải $at$, sóng trái $B=f/2+H/(2a)$ sang trái $at$. Các hình kèm theo gồm lát cắt thời gian; những vị trí gãy được suy từ các cận chạm điểm gãy của dữ liệu. Với các câu trơn, đạo hàm từng mode xác nhận PDE; với câu (iv), (vi), chỉ khẳng định cổ điển ngoài các đường gãy đã chỉ ra.

\textbf{Chi tiết phép thế nguyên hàm và lượng giác.} Ở (i),
\[\frac{H(x+3t)-H(x-3t)}6
=\frac{\cos(2x-6t)-\cos(2x+6t)}{12}
=\frac{\sin2x\sin6t}{6}.\]
Ở (ii), $[(x+5t)^2-(x-5t)^2]/20=xt$. Ở (iii), $[\sin(x+2t)-\sin(x-2t)]/4=\cos x\sin2t/2$. Ở (iv), phần tuyến tính của $H$ cho $(12t/2)/12=t/2$; phần $\sin2s/4$ cho $2\cos2x\sin12t/48=\cos2x\sin12t/24$.

Ở (v), $f(s)=\sin s-\cos^2s=\sin s-1/2-\cos2s/2$, nên half-sum là $\sin x\cos t-1/2-\cos2x\cos2t/2$. Từ $\sin^3s=(3\sin s-\sin3s)/4$, nguyên hàm đã ghi cho
\[\frac{H(x+t)-H(x-t)}2=\frac34\sin x\sin t-\frac1{12}\sin3x\sin3t.\]
Ở (vi), trên $-1<s<0$, $H=\int_{-1}^s(r+1)dr=(s+1)^2/2$; trên $0<s<1$, $H=H(0)+\int_0^s(1-r)dr=1/2+s-s^2/2=1-(1-s)^2/2$. Sau $1$ giữ $H(1)=1$, bảo đảm không có bước nhảy của nguyên hàm.

%ID BT_3_1_3_03
\textbf{Bước 1. Tách sóng li độ và vận tốc.} $a=5$, $f=10\chi_{[-1,1]}$, $g=\chi_{[-1,1]}$. Tại $x=15$, cận phải $15+5t$ luôn ngoài hỗ trợ. Cận trái $15-5t$ nằm trong đoạn khi $14/5\le t\le16/5$. Đặt $\ell(t)=|[15-5t,15+5t]\cap[-1,1]|$.

\textbf{Bước 2. Viết áp suất theo từng giai đoạn.}
\[u(15,t)=\begin{cases}0&0\le t<14/5,\\5+(5t-14)/10&14/5\le t\le16/5,\\1/5&t>16/5.\end{cases}\]
Hạng $5$ là một nửa của li độ $10$; hạng thứ hai là $\ell/(2a)$. Ở hai đầu mút, giá trị phụ thuộc quy ước chỉ thị đóng; các vết một phía vẫn như trên.

\textbf{Bước 3. So sánh ngưỡng.} Trong khoảng giữa, hàm tăng, lớn nhất $5+2/10=26/5=5.2$, thấp hơn ngưỡng $8$. Sau xung chỉ còn $1/5$. Do đó áp suất không vượt ngưỡng. Không thể chỉ xét tích phân vận tốc mà bỏ sóng li độ gây bước nhảy $5$.

%ID BT_3_1_3_04
\textbf{Bước 1. Xung của nguồn tại thời điểm $T$.} Gọi $w(x,\tau;T)$ là nghiệm thuần nhất với $w(x,0;T)=0$, $w_\tau(x,0;T)=F(x,T)$. D'Alembert cho
\[w(x,\tau;T)=\frac1{2a}\int_{x-a\tau}^{x+a\tau}F(y,T)dy.\]
\textbf{Bước 2. Cộng các xung.} Với $V=\int_0^t w(x,t-T;T)dT$, Leibniz cho $V_{tt}=F(x,t)+\int_0^t w_{\tau\tau}dT$, còn $a^2V_{xx}=\int_0^ta^2w_{xx}dT$. Suy $V_{tt}-a^2V_{xx}=F$ và dữ liệu đầu bằng không.

\textbf{Bước 3. Khôi phục dữ liệu đầu.} Nếu đề có $f,g\ne0$ thì $\int w$ chỉ là phần nguồn; phải thêm $u_0$ của d'Alembert. Miền tích phân là tam giác ngược $0\le T\le t$, $|y-x|\le a(t-T)$; đổi thứ tự tích phân trên miền này cho dạng tích phân Green tương đương.

\textbf{Bước 4. Dạng Green trực tiếp.} Trên tam giác đó, đặt $p=u_T$, $q=-a^2u_y$, nên $p_T+q_y=F$. Tích phân và lấy thông lượng hai cạnh đặc trưng, trên đó $du/dT=u_T\mp au_y$, cho
\[\int_0^t\int_{x-a(t-T)}^{x+a(t-T)}F\,dy\,dT
=2au(x,t)-a[f(x-at)+f(x+at)]-\int_{x-at}^{x+at}g(y)dy.\]
Chia $2a$ thu đúng công thức. Công thức này đồng thời chỉ ra dấu và hệ số của từng dữ liệu.

\textbf{Chi tiết thông lượng trên từng cạnh tam giác Green.} Đặt $\ell(T)=x-a(t-T)$, $r(T)=x+a(t-T)$ và $M(T)=\int_{\ell(T)}^{r(T)}u_T(y,T)dy$. Quy tắc Leibniz cho
\[M'=\int_\ell^r u_{TT}dy-a u_T(r,T)-a u_T(\ell,T).\]
Tích phân $u_{TT}-a^2u_{yy}=F$ theo $y$, rồi thay dòng trên:
\[\int_\ell^r Fdy=M'+a(u_T+a u_y)(\ell,T)+a(u_T-a u_y)(r,T).\]
Cạnh trái có $\ell'=a$, nên $(u_T+a u_y)(\ell,T)=d[u(\ell(T),T)]/dT$; cạnh phải có $r'=-a$, nên $(u_T-a u_y)(r,T)=d[u(r(T),T)]/dT$. Tích phân mỗi cạnh từ $0$ đến $t$ cho
\[a[u(x,t)-f(x-at)],\qquad a[u(x,t)-f(x+at)].\]
Mặt đáy cho $M(t)-M(0)=0-\int_{x-at}^{x+at}g(y)dy$, vì ở đỉnh đoạn tích phân co về độ dài $0$. Cộng ba đóng góp được đẳng thức Green ở bước 4, rồi chia $2a$ thu half-sum, tích phân vận tốc và tích phân nguồn.

%ID BT_3_1_3_05
\textbf{Bước 1. Giải phần thuần nhất.} Thay $f,g$ vào d'Alembert. Nguyên hàm cần dùng lần lượt là $-\cos2s/2$, $\arctan s$, $-\cos s$, $e^s$, $\sin s$ và $H$ của xung tam giác. Sau đó tìm phần nguồn có hai dữ liệu đầu bằng không.

\textbf{Bước 2. Tính từng nguồn.}

(i) $a=2$, nguồn $x\cos t$ có nghiệm riêng $x(1-\cos t)$ vì đạo hàm không gian cấp hai bằng không. Suy
\[u=e^{-x}\cosh2t+\tfrac14\sin2x\sin4t+x(1-\cos t).\]
(ii) $a=3$. Nguồn $e^x$ dùng $e^xq$, $q''-9q=1$, $q(0)=q'(0)=0$, cho $q=(\cosh3t-1)/9$. Nguồn $\sin t$ độc lập $x$ cho $t-\sin t$:
\[u=\frac{\arctan(x+3t)-\arctan(x-3t)}6+\frac{e^x(\cosh3t-1)}9+t-\sin t.\]
(iii) $a=4$, nguồn $x+t$ cho $xt^2/2+t^3/6$:
\[u=x^2+16t^2+\tfrac14\sin x\sin4t+xt^2/2+t^3/6.\]
(iv) $a=7$, phần vận tốc $e^x$ là $e^x\sinh7t/7$, phần nguồn $xt$ là $xt^3/6$.

(v) $a=1$, $F=\sin(x+t)=\sin x\cos t+\cos x\sin t$. Hai ODE là $q''+q=\cos t$ và $r''+r=\sin t$ với dữ liệu không, ở cộng hưởng. Suy $q=t\sin t/2$, $r=(\sin t-t\cos t)/2$:
\[u=x+\cos x\sin t+\tfrac12\{\sin x\,t\sin t+\cos x(\sin t-t\cos t)\}.\]
(vi) $a=2$, $f$ là hình thang trên $[0,3]$ với nhánh $s,1,3-s$ tại các đoạn $[0,1],[1,2],[2,3]$. $H$ là nguyên hàm tam giác ở Bài 2(vi). Nguồn $-e^x$ cho
\[u=\tfrac12[f(x-2t)+f(x+2t)]+\frac{H(x+2t)-H(x-2t)}4-\frac{e^x(\cosh2t-1)}4.\]
\textbf{Bước 3. Kiểm tra.} Mỗi phần nguồn và đạo hàm thời gian của nó bằng không ở $t=0$. Các ODE vừa lập chứng minh PDE, không chỉ đoán dạng đáp án. Câu (vi) có góc của $f,g$, nên nghiệm được hiểu từng miền/ theo nghĩa yếu trên các đường $x\pm2t=0,1,2,3$ và $x\pm2t=-1,0,1$.

\textbf{Bước 4. Vẽ đặc trưng và các lát cắt.} Với mỗi tốc độ $a$ của từng câu, các đường xuất phát từ $(b,0)$ là $x=b+at$ và $x=b-at$. Nguồn thay đổi giá trị trên miền ảnh hưởng nhưng không đổi hai họ đặc trưng này. Hình kèm lấy $t=0,0.15,0.35,0.65$ và thay trực tiếp vào sáu công thức đã tính; hình thứ hai cho mặt phẳng $x,t$ với các đường đặc trưng mẫu.

\textbf{Chi tiết giải các ODE nguồn.} (i) Đặt $u_F=xq$, có $q''=\cos t$, $q'(0)=q(0)=0$. Tích phân $q'=\sin t$, $q=1-\cos t$. (ii) Với $u_F=e^xq$, $q''-9q=1$: thử $q_p=-1/9$, phần thuần nhất $A\cosh3t+B\sinh3t$. Dữ liệu cho $A=1/9,B=0$. Nguồn $\sin t$ cho $r''=\sin t$, $r'=1-\cos t$, $r=t-\sin t$.

(iii) $q''=x+t$, tích phân theo thời gian với $x$ cố định được $q'=xt+t^2/2$, $q=xt^2/2+t^3/6$. (iv) $q''=xt$ cho $q'=xt^2/2$, $q=xt^3/6$.

(v) Tính trực tiếp $q=t\sin t/2$: $q'=(\sin t+t\cos t)/2$, $q''=\cos t-t\sin t/2$, nên $q''+q=\cos t$. Với $r=(\sin t-t\cos t)/2$, $r'=t\sin t/2$, $r''=(\sin t+t\cos t)/2$, nên $r''+r=\sin t$. Cả $q,r$ và đạo hàm đầu bằng $0$. (vi) $u_F=e^xq$ cho $q''-4q=-1$, $q_p=1/4$, $q=A\cosh2t+B\sinh2t+1/4$; dữ liệu cho $A=-1/4,B=0$.

%ID BT_3_1_3_06
\textbf{Bước 1. Đạo hàm công thức nửa trục.} Trong miền $x<at$, đặt $s=x+at$, $r=at-x$, $\tau=t-x/a$. Từ $u=h(\tau)+(f(s)-f(r))/2+(2a)^{-1}\int_r^s g$ thu
\[u_x=-h'(\tau)/a+\tfrac12(f'(s)+f'(r))+\frac{g(s)+g(r)}{2a}.\]
Đây là công thức (3.1.13); dấu cộng của $g(r)$ do cận dưới có đạo hàm $r_x=-1$.

\textbf{Bước 2. Đạo hàm thời gian hai lần.}
\[u_{tt}=h''(\tau)+\tfrac{a^2}2(f''(s)-f''(r))+\tfrac a2(g'(s)-g'(r)).\]
Ngoài mặt sóng, thay dấu trừ giữa hai $f''$ thành dấu cộng, và giữ cùng hiệu $g'$. Tại $x=at$, hiệu hai vết là $h''(0)-a^2f''(0)$, cho công thức ghép (3.1.15).

\textbf{Bước 3. Kết luận.} Cần $h(0)=f(0)$ để liên tục, $h'(0)=g(0)$ để ghép cấp một, và $h''(0)=a^2f''(0)$ để ghép cấp hai. Những đạo hàm tính ở đây cũng cho $u_{tt}=a^2u_{xx}$ trong miền phản xạ và đúng biên $u(0,t)=h(t)$.

\textbf{Chi tiết lần đạo hàm thứ nhất theo thời gian.} Với $s_t=r_t=a,\tau_t=1$,
\[u_t=h'(\tau)+\frac a2[f'(s)-f'(r)]+\frac12[g(s)-g(r)].\]
Lấy tiếp đạo hàm: $h'$ cho $h''$, hai hạng $f'$ cho $a^2(f''(s)-f''(r))/2$, hai hạng $g$ cho $a(g'(s)-g'(r))/2$. Tương tự, lấy đạo hàm $x$ của công thức $u_x$ cho
\[u_{xx}=h''(\tau)/a^2+(f''(s)-f''(r))/2+(g'(s)-g'(r))/(2a),\]
chứng minh trực tiếp $u_{tt}=a^2u_{xx}$.

%ID BT_3_1_3_07
\textbf{Bước 1. Mở rộng lẻ.} Đặt $\widetilde f(-s)=-f(s)$, $\widetilde g(-s)=-g(s)$. So vết tại $0$: $f(0)=0$ nối giá trị, $f'$ tự nối, $f''(0)=0$ nối cấp hai; $g(0)=0$ nối giá trị, $g'$ tự nối. Do đó $\widetilde f\in C^2$, $\widetilde g\in C^1$ dưới các điều kiện tương thích.

\textbf{Bước 2. Giải trên toàn trục.} Đặt $H(s)=\int_0^s\widetilde g(r)dr$; $H$ chẵn. D'Alembert cho $u=(\widetilde f(x-at)+\widetilde f(x+at))/2+(H(x+at)-H(x-at))/(2a)$. Tại $x=0$, hai giá trị $\widetilde f$ đối nhau và hai giá trị $H$ bằng nhau, nên $u(0,t)=0$.

\textbf{Bước 3. Viết lại bằng dữ liệu dương.} Khi $x\ge at$ giữ nguyên công thức. Khi $0<x<at$, $\widetilde f(x-at)=-f(at-x)$, còn $H(at+x)-H(at-x)=\int_{at-x}^{at+x}g$. Thu công thức phản xạ trong phần kiến thức với $h=0$. Ranh giới đúng là $x=at$; $x=t$ chỉ đúng nếu $a=1$. Dữ liệu đầu và tính $C^2$ suy trực tiếp từ mở rộng.

%ID BT_3_1_3_08
\textbf{Bước 1. Dùng công thức Dirichlet với thời điểm trễ.} Ở $x\ge at$ sóng chưa chạm biên; ở $x<at$ dùng $h(t-x/a)$ và phản xạ lẻ. Dưới đây giữ hai miền khi biên làm thay đổi vận tốc.

(i) $a=2$, $f=\sin x$, $g=\cos2x$, $h=0$. Mở rộng $g$ lẻ, $H(s)=\sin(2|s|)/2$. Do đó
\[u=\sin x\cos2t+\begin{cases}\cos2x\sin4t/4&x\ge2t,\\\sin2x\cos4t/4&0<x<2t.\end{cases}\]
$g(0)=1\ne h'(0)$ nên không $C^1$ trên mặt sóng. Sau khi mặt sóng đi qua mỗi $x$, công thức tuần hoàn chu kỳ $\pi$; toàn nghiệm không có chu kỳ chung vì hai nhánh khác nhau.

(ii) $a=3$, $f=g=0$, $h=t/(1+t)$. Đặt $\tau=t-x/3$: $u=0$ khi $\tau\le0$, $u=\tau/(1+\tau)$ khi $\tau>0$. Trên tia $x=ct$, giới hạn là $1$ nếu $0\le c<3$, và $0$ nếu $c\ge3$. Nghiệm không tuần hoàn: tại mỗi $x$ cố định, sau khi mặt sóng đến, $u$ tăng nghiêm ngặt về $1$.

(iii) $a=1$, mở rộng lẻ $f$ là $s|s|$, còn $g=\sin s$ đã lẻ:
\[u=\tfrac12[(x+t)|x+t|+(x-t)|x-t|]+\sin x\sin t.\]
Nghiệm $C^1$, không $C^2$ trên $x=t$ vì $f''(0)=2\ne0$, và không tuần hoàn.

(iv) $a=5$, $f=0$, $g=\sin x-\sin3x/3$ có $H=-\cos s+\cos3s/9$. Suy $u=\sin x\sin5t/5-\sin3x\sin15t/45$, trơn, chu kỳ $2\pi/5$.

(v) $a=2$, $f=\cos x-\sin x$, $g=0$, $h=\cos2t-\sin2t$. Ngoài mặt sóng, half-sum của $f$ là $(\cos x-\sin x)\cos2t$. Trong mặt sóng:
\[h(t-x/2)+\tfrac12[f(x+2t)-f(2t-x)]=\cos x(\cos2t-\sin2t).\]
Vậy
\[u=\begin{cases}(\cos x-\sin x)\cos2t&x\ge2t,\\\cos x(\cos2t-\sin2t)&0<x<2t.\end{cases}\]
Công thức cũ $\cos(x-2t)-\sin(x+2t)$ không thỏa $u_t(x,0)=0$. Nghiệm đúng liên tục, không $C^1$ vì $h'(0)=-2\ne0$. Chỉ nhánh sau mặt sóng có chu kỳ $\pi$ ở từng điểm.

(vi) $a=1$, $g=\chi_{[1,2]}-\chi_{[3,4]}$. Mở rộng lẻ có nguyên hàm chẵn $H(s)=[|s|-1]_0^1-[|s|-3]_0^1$, nên $u=(H(x+t)-H(x-t))/2$. Tại $x$ cố định, hai giá trị $H$ đều bằng không sau $t>x+4$, nên $u=0$; đây là xung quá độ, không tuần hoàn. Nghiệm liên tục nhưng có gãy cấp một khi các cận gặp $\pm1,\pm2,\pm3,\pm4$.

\textbf{Bước 2. Kiểm tra chung.} Thế $x=0$ vào nhánh trong cho đúng $h$; thế $t=0$ vào nhánh ngoài cho $f,g$. So hai vết trên $x=at$ theo ba điều kiện tương thích, không dùng tính tuần hoàn của nhánh trong để kết luận tuần hoàn toàn miền.

\textbf{Chi tiết hai nhánh vận tốc ở (i).} Ngoài mặt sóng,
\[\frac14\int_{x-2t}^{x+2t}\cos2s\,ds
=\frac{\sin(2x+4t)-\sin(2x-4t)}8=\frac14\cos2x\sin4t.\]
Trong mặt sóng, cận dưới là $2t-x$, nên hiệu sine bằng $\sin(4t+2x)-\sin(4t-2x)=2\cos4t\sin2x$; chia $8$ cho nhánh đã ghi.

\textbf{Chi tiết rút gọn (v).} Trong miền phản xạ, biên là $\cos(2t-x)-\sin(2t-x)$. Half-difference của $f$ bằng
\[-\sin2t\sin x-\cos2t\sin x.\]
Khai triển biên thành $\cos2t\cos x+\sin2t\sin x-\sin2t\cos x+\cos2t\sin x$ rồi cộng, hai hạng chứa $\sin x$ triệt tiêu, còn $\cos x(\cos2t-\sin2t)$.

%ID BT_3_1_3_09
\textbf{Bước 1. Tìm các hàm sóng ở hai miền.} Biên nghiêng là $t=\alpha x$, không phải $x=0$. Viết $u=p+A(x+at)+B(x-at)$, trong đó $p$ là nghiệm riêng. Dữ liệu ở $t=0$ xác định $A,B$ trên đối số dương; dữ liệu biên xác định $B$ ở đối số âm.

(i) $a=1$, $\alpha=2$, $F=0$, $f=x$, $g=0$, vết biên $3x$. Trên $s>0$, $A(s)=B(s)=s/2$ (bỏ hai hằng số triệt tiêu). Tại $t=2x$, $A(3x)+B(-x)=3x$, nên $B(-x)=3x/2$, tức $B(s)=-3s/2$ khi $s<0$. Suy
\[u=\begin{cases}x&0\le t\le x,\\2t-x&x<t<2x.\end{cases}\]
Hai vết giá trị bằng nhau, nhưng $u_t$ đổi từ $0$ sang $2$; không có nghiệm $C^1$ trên toàn tam giác thỏa mọi dữ liệu này.

(ii) $a=2$, $\alpha=3$, $F=1$, $f=g=0$, biên bằng không. Chọn $p=-x^2/8$, vì $p_{tt}-4p_{xx}=1$. Dữ liệu đầu cho $A=B=s^2/16$ trên $s>0$. Tại $t=3x$, $A(7x)+B(-5x)-x^2/8=0$, suy $B(s)=-47s^2/400$ khi $s<0$:
\[u=\begin{cases}t^2/2&0\le t\le x/2,\\-x^2/8+(x+2t)^2/16-47(x-2t)^2/400&x/2<t<3x.\end{cases}\]
\textbf{Bước 2. Kiểm tra.} Tại đường $x=2t$, $B$ và $B'$ nối bằng không nên nghiệm $C^1$, nhưng $B''$ đổi từ $1/8$ sang $-47/200$, nên không $C^2$. Mỗi nhánh thỏa PDE; tại $t=3x$ đa thức triệt tiêu đúng. Hai kết quả phải hiểu từng miền khi đề yêu cầu mức trơn cao hơn mức dữ liệu cho phép.

%ID BT_3_1_3_10
\textbf{Bước 1. Điều kiện Neumann.} $u_x(0,t)=0$ kéo theo $f'(0)=0$, $g'(0)=0$ nếu nghiệm $C^2$. Mở rộng chẵn của $f,g$ có vết $f'$ và $g'$ đổi dấu, nên hai điều kiện này bảo đảm các vết cần thiết nối nhau.

\textbf{Bước 2. Áp dụng d'Alembert.} Với mở rộng chẵn $\widetilde f,\widetilde g$, nguyên hàm $H=\int_0^s\widetilde g$ là lẻ. Nghiệm là half-sum cộng hiệu $H/(2a)$. Nó chẵn theo $x$, do đó đạo hàm không gian ở $0$ bằng không. Trong miền $x<at$ có thể viết trực tiếp
\[u=\frac{f(at-x)+f(at+x)}2+\frac1a\int_0^{at-x}g(s)ds+\frac1{2a}\int_{at-x}^{at+x}g(s)ds.\]
Hạng tích phân đầu xuất hiện vì đoạn $[x-at,0]$ cũng góp vận tốc sau phản xạ chẵn.

\textbf{Bước 3. Ví dụ.} Với $a=2$, $f=x^2$, $g=\cos x$, cả hai đã chẵn và tương thích. Half-sum cho $x^2+4t^2$, nguyên hàm $\sin s$ cho $\cos x\sin2t/2$. Vậy $u=x^2+4t^2+\cos x\sin2t/2$. Kiểm tra $u_x(0,t)=0$, $u_t(x,0)=\cos x$ và $u_{tt}=4u_{xx}$ trực tiếp.

%ID BT_3_1_3_11
\textbf{Bước 1. Chuyển đạo hàm thành bài Dirichlet.} Đặt $v=u_x$. Khi dữ liệu đủ trơn, $v_{tt}=a^2v_{xx}$, $v(0,t)=h(t)$, $v(x,0)=f'(x)$, $v_t(x,0)=g'(x)$. Tương thích cấp hai cho bài này là $h(0)=f'(0)$, $h'(0)=g'(0)$, $h''(0)=a^2f'''(0)$.

\textbf{Bước 2. Tích phân lại và giữ hằng số theo thời gian.} $u(x,t)=C(t)+\int_0^xv(s,t)ds$. Thế vào PDE cho $C''(t)=a^2v_x(0,t)$, $C(0)=f(0)$, $C'(0)=g(0)$. Không được bỏ $C$ chỉ vì đã biết $u_x$.

\textbf{Bước 3. Công thức tiện dụng tương đương.} Gọi $u_e$ là nghiệm d'Alembert của các mở rộng chẵn $f,g$. Đặt
\[u=u_e-a\int_0^{(t-x/a)_+}h(s)ds.\]
Ở $x<at$, đạo hàm không gian của hạng cuối là $h(t-x/a)$, cho đúng $u_x(0,t)=h(t)$; hạng đó là sóng chạy nên thỏa PDE. Nếu $h(0)\ne0$, góc của $u_e$ do $f'(0)$ được hạng biên bù khi tương thích đúng.

\textbf{Bước 4. Viết hai nhánh.} Ngoài mặt sóng dùng d'Alembert với dữ liệu dương. Trong mặt sóng dùng công thức Bài 10 rồi trừ $a\int_0^{t-x/a}h$. Hai dữ liệu đầu giữ nguyên vì hạng bổ sung chưa đến $x>0$ tại $t=0$. Nếu chỉ cần $u\in C^2$, hai tương thích đầu đủ với mức trơn phù hợp; điều kiện thứ ba ở trên phục vụ $v\in C^2$, tức mức trơn cao hơn của $u$.

%ID BT_3_1_3_12
\textbf{Bước 1. Tách hai sóng trên dữ liệu gốc.} Đặt $H(s)=\int_0^sg$, $A=f/2-H/(2a)$, $B=f/2+H/(2a)$ trên $0\le s\le L$. Hai biên là $u(0,t)=\alpha(t)$ và $u(L,t)=\beta(t)$. Chỉ xét $t\le L/a$, nên mỗi chân đặc trưng phản xạ tối đa một lần.

\textbf{Bước 2. Viết đủ bốn miền.} Đặt $\xi=x-at$, $\eta=x+at$:
\[u=\begin{cases}
A(\xi)+B(\eta)&\xi\ge0,\ \eta\le L,\\
\alpha(t-x/a)-B(at-x)+B(x+at)&\xi<0,\ \eta\le L,\\
A(x-at)+\beta(t-(L-x)/a)-A(2L-x-at)&\xi\ge0,\ \eta>L,\\
\alpha(t-x/a)+\beta(t-(L-x)/a)-B(at-x)-A(2L-x-at)&\xi<0,\ \eta>L.
\end{cases}\]
Hai công thức phản xạ suy từ $A(-s)=\alpha(s/a)-B(s)$ và $B(L+s)=\beta(s/a)-A(L-s)$; mọi đối số của $A,B$ còn lại đều trong $[0,L]$.

\textbf{Bước 3. Trường hợp (i).} $a=2$, $L=\pi$, $f=\sin x$, $g=1$, $\alpha=\beta=0$. $A=\sin s/2-s/4$, $B=\sin s/2+s/4$. Trong bốn miền theo đúng thứ tự trên:
\[u=\sin x\cos2t+\begin{cases}t&\text{miền I},\\x/2&\text{miền II},\\(\pi-x)/2&\text{miền III},\\\pi/2-t&\text{miền IV}.\end{cases}\]
Phần vận tốc liên tục, nhưng có góc trên hai mặt sóng vì $g$ không tương thích với biên không.

\textbf{Bước 4. Trường hợp (ii).} $a=L=1$, $f=g=0$, $\alpha(t)=\sin\pi t$, $\beta(t)=t$. Vì $A=B=0$, bốn nhánh là $0$, $\sin\pi(t-x)$, $t-1+x$, và $\sin\pi(t-x)+t-1+x$. Thế $x=0,1$ trong khoảng $0\le t\le1$ cho đúng hai biên. Vận tốc góc không tương thích ($\alpha'(0)=\pi$, $\beta'(0)=1$), nên không có nghiệm $C^1$ trên toàn miền; các nhánh và đường ghép đã mô tả chính xác sự mất trơn.

\textbf{Chi tiết thế vào miền cả hai mặt sóng đã phản xạ ở (i).} Đặt $p=2t-x$, $q=2\pi-x-2t$. Khi $\xi<0,\eta>\pi$, $p,q\in[0,\pi]$ và $u=-B(p)-A(q)$:
\[u=-\frac{\sin(2t-x)}2-\frac{\sin(2\pi-x-2t)}2+\frac{q-p}{4}
=\frac{\sin(x+2t)+\sin(x-2t)}2+\frac{2\pi-4t}{4}.\]
Tổng sine bằng $2\sin x\cos2t$, hạng tuyến tính bằng $\pi/2-t$. Ở miền chỉ phản xạ trái, $-B(2t-x)+B(x+2t)$ cho hạng tuyến tính $[(x+2t)-(2t-x)]/4=x/2$. Ở miền chỉ phản xạ phải, $A(x-2t)-A(2\pi-x-2t)$ cho $[-(x-2t)+(2\pi-x-2t)]/4=(\pi-x)/2$.

%ID VD_3_2_2_01
\textbf{Bước 1. Chia dữ liệu và hai sóng.} $a=1$, $f(s)=s(s+1)$ trên $[-1,0]$, bằng $0$ ngoài đoạn, $g=0$. Vì vậy $A=B=f/2$, không cần nguyên hàm vận tốc.

\textbf{Bước 2. Dịch từng đoạn.} Sóng phải $A(x-t)$ khác không khi $t-1<x<t$, sóng trái $B(x+t)$ khác không khi $-t-1<x<-t$. Nghiệm chính xác là
\[u(x,t)=\tfrac12(x-t)(x-t+1)\chi_{[-1,0]}(x-t)+\tfrac12(x+t)(x+t+1)\chi_{[-1,0]}(x+t).\]
Khi hai đoạn chồng nhau, cộng hai đa thức; khi chỉ một đoạn hiện diện lấy một hạng; ngoài cả hai đoạn lấy $0$. Đây cũng là cách dựng các hình trong sách.

\textbf{Bước 3. Xét tính trơn.} $f$ bằng $0$ tại hai đầu nên $u$ liên tục. $f'$ trong đoạn bằng $2s+1$, có bước nhảy tại $-1,0$, nên các đường gãy cấp một là $x\pm t=-1,0$. Ở giao điểm $(-1/2,1/2)$, lát cắt $u(x,1/2)$ có đạo hàm trái theo $x$ bằng $1/2$, đạo hàm phải bằng $-1/2$; vì vậy điểm ấy vẫn không $C^1$. Tại giao của các mặt sóng phải so trực tiếp các đạo hàm một phía, không mặc nhiên giả sử chúng triệt tiêu.

\textbf{Bước 4. Kiểm tra.} Thế $t=0$ cho $u=f$, còn đạo hàm theo $t$ của hai sóng đối nhau nên vận tốc đầu bằng $0$ nơi $f'$ tồn tại. Mỗi đa thức của $x\pm t$ thỏa PDE; toàn nghiệm đúng theo nghĩa yếu trên các đường gãy.

\textbf{Chi tiết đạo hàm tại giao điểm hai mặt sóng.} Với $t=1/2$, bên trái $x=-1/2$, chỉ $x+t=x+1/2\in[-1,0]$ góp:
\[u(x,1/2)=\tfrac12(x+1/2)(x+3/2),\quad u_x=x+1\to1/2.\]
Bên phải, chỉ $x-t=x-1/2\in[-1,0]$ góp:
\[u(x,1/2)=\tfrac12(x-1/2)(x+1/2),\quad u_x=x\to-1/2.\]
Hai vết giá trị đều $0$, nhưng hai đạo hàm không bằng nhau, nên không thể gọi điểm giao là điểm trơn.

%ID VD_3_2_2_02
\textbf{Bước 1. Tích phân vận tốc.} $f=0$, $g=\chi_{[-1,0]}$, $a=1$. Chọn
\[H(s)=\begin{cases}0&s\le-1,\\s+1&-1<s<0,\\1&s\ge0.\end{cases}\]
Hai hằng số đuôi đã nối liên tục với nhánh giữa.

\textbf{Bước 2. Tách và ghép sóng.} $A=-H/2$, $B=H/2$, nên $u=(H(x+t)-H(x-t))/2$. Tương đương
\[u=\tfrac12\big[\min(x+t,0)-\max(x-t,-1)\big]_+.\]
Biểu thức trong ngoặc là độ dài giao của khoảng ảnh hưởng với hỗ trợ vận tốc. Trên mỗi miền mà các min/max chọn cùng nhánh, $u$ là affine theo $x,t$.

\textbf{Bước 3. Miền gãy và giá trị lâu dài.} $H$ liên tục nên $u$ liên tục. $H'$ nhảy tại $-1,0$, phát sinh các đường $x\pm t=-1,0$ mất trơn cấp một. Với mỗi $x$ cố định, khi $t>\max(|x|,|x+1|)$, toàn đoạn $[-1,0]$ nằm trong khoảng ảnh hưởng, nên $u=1/2$. Kiểm tra tại $t=0$ có hiệu $H=0$ và $u_t=g$ tại các điểm liên tục của dữ liệu.

%ID VD_3_2_2_03
\textbf{Bước 1. Phản xạ Dirichlet.} $a=2$, $f=0$, $g=\chi_{[1,2]}$ trên nửa trục, biên $u(0,t)=0$. Mở rộng vận tốc lẻ: bằng $1$ trên $(1,2)$, $-1$ trên $(-2,-1)$, bằng $0$ còn lại.

\textbf{Bước 2. Nguyên hàm chẵn.} Với $H(0)=0$, trên đối số dương $H=0$ trước $1$, $s-1$ trên $(1,2)$, $1$ sau $2$; phản xạ được $H(s)=[|s|-1]_0^1$. Vì thế
\[A=-H/4,\quad B=H/4,\qquad u=\frac{[|x+2t|-1]_0^1-[|x-2t|-1]_0^1}{4}.\]
\textbf{Bước 3. Kiểm tra biên và mặt sóng.} Tại $x=0$, hai cận có cùng trị tuyệt đối nên $u=0$. $u_t(x,0)=H'(x)=g(x)$ tại điểm liên tục. Nghiệm liên tục nhưng không $C^1$ trên các nhánh $x\pm2t=\pm1,\pm2$ nằm trong $x,t>0$. Chúng chia đồ thị thành các đoạn thẳng. Tại mỗi $x$ cố định, sau $t>(x+2)/2$, hai đuôi bằng nhau và $u=0$.

%ID VD_3_2_2_04
\textbf{Bước 1. Phản xạ Neumann.} Dữ liệu như Ví dụ 3, nhưng $u_x(0,t)=0$. Mở rộng vận tốc chẵn thành $\chi_{[1,2]}(|s|)$, nên nguyên hàm phải lẻ, không dùng nguyên hàm chẵn của ví dụ trước.

\textbf{Bước 2. Tính hai sóng.}
\[H(s)=\operatorname{sgn}(s)[|s|-1]_0^1,\quad A=-H/4,\quad B=H/4,\quad u=\frac{H(x+2t)-H(x-2t)}4.\]
Trước $x=2t$ vẫn có thể dùng các cận dương; sau đó giá trị $H(x-2t)$ âm, tạo xung phản xạ cùng dấu.

\textbf{Bước 3. Kiểm tra.} $u_x(0,t)=(H'(2t)-H'(-2t))/4=0$. $u$ liên tục, gãy cấp một trên $x\pm2t=\pm1,\pm2$. Độ dịch đúng là $2t$ vì PDE có hệ số $4$; những chỗ ghi $x\pm t$ trong hình/công thức cần hiểu theo tốc độ thật. Sau $t>(x+2)/2$, hai đuôi $H$ là $1,-1$, nên $u=1/2$, khác phản xạ Dirichlet trở về $0$.

%ID VD_3_2_2_05
\textbf{Bước 1. Phản xạ hai biên cố định.} $a=L=1$, $f=\chi_{[1/4,1/3]}$, $g=4\chi_{[2/3,3/4]}$. Mở rộng lẻ tại $0,1$ cho chu kỳ $2$. Đặt $r=\operatorname{mod}(s,2)$: $f_e(s)=f(r)$ nếu $r\le1$, $-f(2-r)$ nếu $r>1$; vận tốc mở rộng tương tự.

\textbf{Bước 2. Nguyên hàm trong một chu kỳ.} Trên $[0,1]$, $H(s)=4[s-2/3]_0^{1/12}$; nó tăng từ $0$ lên $1/3$. Trên $[1,2]$, $H(s)=H(2-s)$, sau đó mở rộng chu kỳ $2$. Sự giảm trên nửa sau do vận tốc phản xạ âm, nên $H(2)=0$.

\textbf{Bước 3. Ghép nghiệm.}
\[u=\tfrac12[f_e(x-t)+f_e(x+t)]+\tfrac12[H(x+t)-H(x-t)].\]
Hai sóng chạy là $A=f_e/2-H/2$, $B=f_e/2+H/2$. Tính từng hình bằng giảm các đối số về $[0,2]$, rồi chọn các nhánh trên; không cần viết một danh sách dài công thức thời gian.

\textbf{Bước 4. Chu kỳ và tính trơn.} Cả $f_e,H$ có chu kỳ $2$, nên $u(t+2)=u(t)$. Li độ bước nhảy gây gián đoạn trên các đường $x\pm t=\pm1/4+2k,\pm1/3+2k$; vận tốc gây gãy cấp một trên $x\pm t=\pm2/3+2k,\pm3/4+2k$, $k\in\mathbb Z$. Ngoài các đường ấy, nghiệm affine cộng hằng. Biên $u(0,t)=u(1,t)=0$ theo tính lẻ tại hai đầu; dữ liệu ban đầu đúng trong miền theo nghĩa yếu.

%ID VD_3_2_2_06
\textbf{Bước 1. Phản xạ biên hỗn hợp.} Dữ liệu giống Ví dụ 5, nhưng $u(0,t)=0$, $u_x(1,t)=0$. Trên $[0,2]$, mở rộng $f_e(s)=f(s)$ nếu $s\le1$, $f(2-s)$ nếu $s\ge1$; sau đó $f_e(s+2)=-f_e(s)$ và chu kỳ $4$. Làm tương tự với $g_e$.

\textbf{Bước 2. Tính nguyên hàm bốn nhánh.} Đặt $h(s)=4[s-2/3]_0^{1/12}$ trên $[0,1]$. Trên $[0,4]$:
\[H(s)=\begin{cases}h(s)&0\le s\le1,\\2/3-h(2-s)&1\le s\le2,\\2/3-H(s-2)&2\le s\le4.\end{cases}\]
Hạng $2/3$ là tích phân vận tốc trên $[0,2]$; toàn chu kỳ $4$ có trung bình bằng $0$, nên mở rộng $H$ chu kỳ $4$.

\textbf{Bước 3. Nghiệm và phản chu kỳ.} Vẫn dùng $u=(f_e(x-t)+f_e(x+t))/2+(H(x+t)-H(x-t))/2$. Ta có $H(s+2)=2/3-H(s)$, nên
\[u(x,t+2)=-u(x,t),\qquad u(x,t+4)=u(x,t).\]
Vì dữ liệu khác không, không có chu kỳ $2$; câu yêu cầu chu kỳ $2$ trong sách cần sửa thành phản chu kỳ $2$ và chu kỳ $4$. Mode đầu khác không (tích phân $f\sin(\pi x/2)>0$), xác nhận chu kỳ cơ bản $4$.

\textbf{Bước 4. Kiểm tra biên.} $f_e,g_e$ lẻ tại $0$, chẵn tại $1$, nên nghiệm có đúng vết Dirichlet và Neumann nơi vết đạo hàm tồn tại. Các bước nhảy dữ liệu lan dọc $x\pm t=b$, trong đó $b$ là các điểm biên hỗ trợ cùng các ảnh phản xạ đã xây dựng. Mức trơn là nghiệm từng miền/ nghiệm yếu, không $C^2$ trên toàn hình chữ nhật.

\textbf{Chi tiết nối nguyên hàm và phản chu kỳ.} Xung dương có diện tích $4(3/4-2/3)=1/3$. Phản xạ chẵn tại $1$ tạo thêm diện tích $1/3$ trên $(1,2)$, nên $H(2)=2/3$. Vì $g_e(s+2)=-g_e(s)$,
\[H(s+2)=H(2)+\int_2^{s+2}g_e(r)dr=2/3-\int_0^sg_e(r)dr=2/3-H(s).\]
Khi dịch $t$ thêm $2$, các đối số $x+t$ tăng $2$, $x-t$ giảm $2$; cả hai giá trị $H$ thành $2/3-H$ của giá trị cũ. Vì thế hiệu $H$ đổi dấu. Các giá trị $f_e$ cũng đổi dấu, cho $u(t+2)=-u(t)$.

%ID BT_3_2_3_LT_01
\textbf{Bước 1. Tính liên tục tuyệt đối.} Với $F(x)=\int_0^xf(s)ds$, trên đoạn hữu hạn $[a,b]$ hàm Riemann khả tích $f$ bị chặn bởi $M$. Với các khoảng rời nhau $(a_j,b_j)$,
\[\sum_j|F(b_j)-F(a_j)|\le M\sum_j(b_j-a_j).\]
Chọn tổng độ dài nhỏ hơn $\varepsilon/M$ chứng minh liên tục tuyệt đối (trường hợp $M=0$ hiển nhiên), và $F'=f$ tại mọi điểm liên tục của $f$.

\textbf{Bước 2. Tính chẵn/lẻ.} Đổi biến $s=-r$ cho $F(-x)=-\int_0^xf(-r)dr$. Nếu $f$ lẻ, kết quả là $F(x)$; nếu $f$ chẵn, kết quả là $-F(x)$. Vì vậy nguyên hàm lấy gốc tại $0$ đổi tính chẵn/lẻ.

\textbf{Bước 3. Điều kiện tuần hoàn.} Nếu $f$ có chu kỳ $L$ thì
\[F(x+L)-F(x)=\int_x^{x+L}f(s)ds=\int_0^Lf(s)ds.\]
Do đó $F$ có chu kỳ $L$ khi và chỉ khi trung bình của $f$ bằng $0$. Nếu $f$ lẻ và tuần hoàn, tích phân trên chu kỳ đối xứng $[-L/2,L/2]$ bằng $0$, nên điều kiện đúng. Ví dụ $f=1$ tuần hoàn nhưng $F=x$ không tuần hoàn, cho thấy không được bỏ điều kiện trung bình.

%ID BT_3_2_3_LT_02
\textbf{Bước 1. Dùng đối xứng của dữ liệu.} Với biên cố định hai đầu và $f(L-x)=f(x)$, $g(L-x)=g(x)$, hệ số sine $A_n=(2/L)\int_0^Lf(x)\sin(n\pi x/L)dx$. Đổi $y=L-x$ cho $A_n=(-1)^{n+1}A_n$; do đó $A_n=0$ khi $n$ chẵn. Cùng lập luận cho $B_n$.

\textbf{Bước 2. Dịch thời gian từng mode còn lại.} Với $n$ lẻ, $\omega_n=n\pi c/L$ và $\omega_nL/c=n\pi$, nên cả cos và sin đổi dấu khi dịch $L/c$. Vì vậy $u(x,t+L/c)=-u(x,t)$ và chu kỳ $2L/c$.

\textbf{Bước 3. Chứng minh bằng phản xạ nếu dữ liệu chưa trơn.} Mở rộng lẻ chu kỳ $2L$ của một hàm đối xứng giữa đoạn thỏa $f_e(s+L)=-f_e(s)$; tương tự $g_e$. Với $H'=g_e$, $H(s+L)=C-H(s)$. Thế vào d'Alembert, hằng số $C$ triệt tiêu, thu cùng phản chu kỳ. Lập luận này áp dụng trực tiếp cho dữ liệu chia đoạn, không cần đạo hàm chuỗi Fourier.

%ID BT_3_2_3_LT_03
\textbf{Bước 1. Viết hàm đường thẳng dưới dạng affine.} $f(x)=m_1x+b_1$, $g(x)=m_2x+b_2$ trên đoạn xét. Với hai hệ số $a,b$,
\[af+bg=(am_1+bm_2)x+ab_1+bb_2.\]
Đây vẫn là hàm affine, có thể là hằng nếu hệ số góc bằng $0$.

\textbf{Bước 2. Co giãn và tịnh tiến đối số.} $f(ax)=m_1ax+b_1$, $f(x-h)=m_1x+b_1-m_1h$. Phải giới hạn $x$ để $ax$ hoặc $x-h$ nằm trong đoạn định nghĩa ban đầu; nếu $a<0$, chiều của đoạn đảo; nếu $a=0$, cần $0$ thuộc miền của $f$ và kết quả là hằng.

\textbf{Bước 3. Hiệu của hai dịch chuyển bậc hai.} Với $f(s)=As^2+Bs+C$, khai triển hai bình phương và triệt tiêu $Aa^2x^2,C$:
\[f(ax+h_1)-f(ax+h_2)=2Aa(h_1-h_2)x+A(h_1^2-h_2^2)+B(h_1-h_2).\]
Kết quả affine trên giao hai miền hợp lệ. Vì vậy tích phân một vận tốc affine cho nguyên hàm bậc hai, nhưng hiệu hai cận dịch cùng tốc độ có thể trở lại affine; đây là cơ sở của các đoạn thẳng trong đồ thị sóng ở mục 3.2.

%ID BT_3_2_3_01
\textbf{Bước 1. Đọc dữ liệu và chọn phản xạ.} Toàn trục, $a=2$, $f=0$, $g=s$ trên $[-1,0]$, bằng không bên ngoài.

\textbf{Bước 2. Tính từng nhánh và nối nguyên hàm.} Tích phân từ $-1$ đến $s$ cho $(s^2-1)/2$; trước $-1$ lấy $0$, sau $0$ giữ $-1/2$. Do đó $H$ liên tục ở cả hai đầu.
\[f_e(s)=0,\qquad H(s)=\tfrac12(\operatorname{clip}(s,-1,0)^2-1).\]
Ở đây $\operatorname{clip}(s,p,q)=\min(q,\max(p,s))$; $\operatorname{mod}(s,L)\in[0,L)$ kể cả với $s<0$.

\textbf{Bước 3. Xác định hai sóng chạy và ghép nghiệm.} $A=f_e/2-H/(2a)$, $B=f_e/2+H/(2a)$. Thay hai cận:
\[u(x,t)=\frac{f_e(x-2t)+f_e(x+2t)}2+\frac{H(x+2t)-H(x-2t)}{4}.\]

\textbf{Bước 4. Trả lời các yêu cầu và kiểm tra.} Tại $t=1$, cận $x+2=-1$ cho $x=-3$, cận $x-2=-1$ cho $x=1$. Bước nhảy của $g$ ở $-1$ khác không nên hai điểm này không có $u_x$. Ở $0$, $g$ nối về $0$ nên $u_x$ vẫn tồn tại; chỉ $u_{xx}$ gãy.

Tại $t=0$, hiệu nguyên hàm bằng $0$ và đạo hàm thời gian cho $H'(x)=g(x)$ trong miền dữ liệu; half-sum cho đúng li độ. Mỗi hạng là sóng của $x\pm at$, nên thỏa PDE ngoài các mặt gãy và theo nghĩa phân phối trên toàn miền. Đồ thị hai sóng ban đầu và các lát cắt thời gian được tính trực tiếp từ các nhánh trên.

%ID BT_3_2_3_02
\textbf{Bước 1. Đọc dữ liệu và chọn phản xạ.} Toàn trục, $a=2$, $f=(1-|s|)_+$, $g=\chi_{[1,2]}$.

\textbf{Bước 2. Tính từng nhánh và nối nguyên hàm.} $H=0$ trước $1$, $s-1$ trên $[1,2]$, $1$ sau $2$. Các nhánh của $f$ có hệ số góc $0,1,-1,0$ tại $-1,0,1$.
\[f_e(s)=(1-|s|)_+,\qquad H(s)=\operatorname{clip}(s-1,0,1).\]
Ở đây $\operatorname{clip}(s,p,q)=\min(q,\max(p,s))$; $\operatorname{mod}(s,L)\in[0,L)$ kể cả với $s<0$.

\textbf{Bước 3. Xác định hai sóng chạy và ghép nghiệm.} $A=f_e/2-H/(2a)$, $B=f_e/2+H/(2a)$. Thay hai cận:
\[u(x,t)=\frac{f_e(x-2t)+f_e(x+2t)}2+\frac{H(x+2t)-H(x-2t)}{4}.\]

\textbf{Bước 4. Trả lời các yêu cầu và kiểm tra.} $f,H$ liên tục nên $u$ liên tục ở mọi $(x,t)$, kể cả $t=1$. Ở $t=1$, vị trí ứng viên của $u_x$ là $-3,-2,-1,0,1,2,3,4$; cộng các bước nhảy theo công thức đạo hàm để xét từng lát cắt. Trong hai biến, các đường xuất phát từ $-1,0,1$ của $f$ và $1,2$ của $g$ là các mặt sóng cấp một.

Tại $t=0$, hiệu nguyên hàm bằng $0$ và đạo hàm thời gian cho $H'(x)=g(x)$ trong miền dữ liệu; half-sum cho đúng li độ. Mỗi hạng là sóng của $x\pm at$, nên thỏa PDE ngoài các mặt gãy và theo nghĩa phân phối trên toàn miền. Đồ thị hai sóng ban đầu và các lát cắt thời gian được tính trực tiếp từ các nhánh trên.

%ID BT_3_2_3_03
\textbf{Bước 1. Đọc dữ liệu và chọn phản xạ.} Toàn trục, $a=3$, $f=\chi_{[-1,1]}$, $g=s\chi_{[1,2]}$.

\textbf{Bước 2. Tính từng nhánh và nối nguyên hàm.} Trên $[1,2]$, $\int_1^s r,dr=(s^2-1)/2$; sau $2$ giữ $3/2$. Phần li độ là hai xung biên độ $1/2$, chạy ngược chiều.
\[f_e(s)=\chi_{[-1,1]}(s),\qquad H(s)=\tfrac12(\operatorname{clip}(s,1,2)^2-1).\]
Ở đây $\operatorname{clip}(s,p,q)=\min(q,\max(p,s))$; $\operatorname{mod}(s,L)\in[0,L)$ kể cả với $s<0$.

\textbf{Bước 3. Xác định hai sóng chạy và ghép nghiệm.} $A=f_e/2-H/(2a)$, $B=f_e/2+H/(2a)$. Thay hai cận:
\[u(x,t)=\frac{f_e(x-3t)+f_e(x+3t)}2+\frac{H(x+3t)-H(x-3t)}{6}.\]

\textbf{Bước 4. Trả lời các yêu cầu và kiểm tra.} Tại $t=1$, hai đoạn li độ là $[-4,-2]$ và $[2,4]$, có bước nhảy tại $-4,-2,2,4$. Phần $H$ liên tục nên không thể khử các bước nhảy ấy. Trong hai biến, ở giao điểm hai đường, cần so vết theo nhiều hướng; triệt tiêu trên một lát cắt không xóa mất kỳ dị hai chiều.

Tại $t=0$, hiệu nguyên hàm bằng $0$ và đạo hàm thời gian cho $H'(x)=g(x)$ trong miền dữ liệu; half-sum cho đúng li độ. Mỗi hạng là sóng của $x\pm at$, nên thỏa PDE ngoài các mặt gãy và theo nghĩa phân phối trên toàn miền. Đồ thị hai sóng ban đầu và các lát cắt thời gian được tính trực tiếp từ các nhánh trên.

%ID BT_3_2_3_04
\textbf{Bước 1. Đọc dữ liệu và chọn phản xạ.} Toàn trục, $a=2$, $f=\chi_{[-2,-1]}$, $g=s\chi_{[0,1]}$.

\textbf{Bước 2. Tính từng nhánh và nối nguyên hàm.} $H=0$ trước $0$, $s^2/2$ trên $[0,1]$, $1/2$ sau $1$. $g$ liên tục tại $0$, có bước nhảy từ $1$ về $0$ tại $1$.
\[f_e(s)=\chi_{[-2,-1]}(s),\qquad H(s)=\tfrac12\operatorname{clip}(s,0,1)^2.\]
Ở đây $\operatorname{clip}(s,p,q)=\min(q,\max(p,s))$; $\operatorname{mod}(s,L)\in[0,L)$ kể cả với $s<0$.

\textbf{Bước 3. Xác định hai sóng chạy và ghép nghiệm.} $A=f_e/2-H/(2a)$, $B=f_e/2+H/(2a)$. Thay hai cận:
\[u(x,t)=\frac{f_e(x-2t)+f_e(x+2t)}2+\frac{H(x+2t)-H(x-2t)}{4}.\]

\textbf{Bước 4. Trả lời các yêu cầu và kiểm tra.} $u$ gián đoạn trên các nhánh $x\pm2t=-2,-1$. Vận tốc chỉ làm gãy $u_x$ trên $x\pm2t=1$; tại $x\pm2t=0$ chỉ gãy đạo hàm cấp hai. Ở $t=1$, các bước nhảy li độ nằm tại $-4,-3,0,1$.

Tại $t=0$, hiệu nguyên hàm bằng $0$ và đạo hàm thời gian cho $H'(x)=g(x)$ trong miền dữ liệu; half-sum cho đúng li độ. Mỗi hạng là sóng của $x\pm at$, nên thỏa PDE ngoài các mặt gãy và theo nghĩa phân phối trên toàn miền. Đồ thị hai sóng ban đầu và các lát cắt thời gian được tính trực tiếp từ các nhánh trên.

%ID BT_3_2_3_05
\textbf{Bước 1. Đọc dữ liệu và chọn phản xạ.} Toàn trục, $a=4$, $f=10\chi_{[-1,1]}$, $g=\chi_{[-1,1]}$.

\textbf{Bước 2. Tính từng nhánh và nối nguyên hàm.} $H(s)=0$ khi $s\le-1$, $s+1$ khi $-1<s<1$, $2$ khi $s\ge1$. Tại $x=10$, sóng phải đến lúc $10-4t=1$ và rời lúc $10-4t=-1$.
\[f_e(s)=10\chi_{[-1,1]}(s),\qquad H(s)=\operatorname{clip}(s+1,0,2).\]
Ở đây $\operatorname{clip}(s,p,q)=\min(q,\max(p,s))$; $\operatorname{mod}(s,L)\in[0,L)$ kể cả với $s<0$.

\textbf{Bước 3. Xác định hai sóng chạy và ghép nghiệm.} $A=f_e/2-H/(2a)$, $B=f_e/2+H/(2a)$. Thay hai cận:
\[u(x,t)=\frac{f_e(x-4t)+f_e(x+4t)}2+\frac{H(x+4t)-H(x-4t)}{8}.\]

\textbf{Bước 4. Trả lời các yêu cầu và kiểm tra.} Đó là $t=9/4$ và $11/4$. Trong khoảng giữa, $u(10,t)=5+(4t-9)/8$; trước đó $0$, sau đó $1/4$. Cực đại $21/4<6$, kể cả vết một phía tại bước nhảy. Không đạt ngưỡng phá hủy.

Tại $t=0$, hiệu nguyên hàm bằng $0$ và đạo hàm thời gian cho $H'(x)=g(x)$ trong miền dữ liệu; half-sum cho đúng li độ. Mỗi hạng là sóng của $x\pm at$, nên thỏa PDE ngoài các mặt gãy và theo nghĩa phân phối trên toàn miền. Đồ thị hai sóng ban đầu và các lát cắt thời gian được tính trực tiếp từ các nhánh trên.

%ID BT_3_2_3_06
\textbf{Bước 1. Đọc dữ liệu và chọn phản xạ.} Nửa trục Dirichlet, $a=3$, $f=s(s-1)$ trên $[0,1]$, $g=6$.

\textbf{Bước 2. Tính từng nhánh và nối nguyên hàm.} Mở rộng lẻ của $f$ là $\operatorname{sgn}(s)|s|(|s|-1)\chi_{[0,1]}(|s|)$. Mở rộng lẻ của $g$ là $6\operatorname{sgn}(s)$, nên $H=6|s|$.
\[f_e(s)=\operatorname{sgn}(s)|s|(|s|-1)\chi_{[0,1]}(|s|),\qquad H(s)=6|s|.\]
Ở đây $\operatorname{clip}(s,p,q)=\min(q,\max(p,s))$; $\operatorname{mod}(s,L)\in[0,L)$ kể cả với $s<0$.

\textbf{Bước 3. Xác định hai sóng chạy và ghép nghiệm.} $A=f_e/2-H/(2a)$, $B=f_e/2+H/(2a)$. Thay hai cận:
\[u(x,t)=\frac{f_e(x-3t)+f_e(x+3t)}2+\frac{H(x+3t)-H(x-3t)}{6}.\]

\textbf{Bước 4. Trả lời các yêu cầu và kiểm tra.} Ở $t=1/3$, $x-1=0$ cho bước nhảy của $g$, tức $x=1$. Góc của $f$ tại $s=1$ cho $x=2$; các góc khác nằm ngoài $x>0$. $f$ nối đạo hàm ở $0$, nên không khử bước nhảy vận tốc tại $x=1$.

Tại $t=0$, hiệu nguyên hàm bằng $0$ và đạo hàm thời gian cho $H'(x)=g(x)$ trong miền dữ liệu; half-sum cho đúng li độ. Mỗi hạng là sóng của $x\pm at$, nên thỏa PDE ngoài các mặt gãy và theo nghĩa phân phối trên toàn miền. Đồ thị hai sóng ban đầu và các lát cắt thời gian được tính trực tiếp từ các nhánh trên.

%ID BT_3_2_3_07
\textbf{Bước 1. Đọc dữ liệu và chọn phản xạ.} Nửa trục Neumann, $a=1$, $f=0$, $g=s-1$ trên $[0,1]$, bằng không ngoài đoạn.

\textbf{Bước 2. Tính từng nhánh và nối nguyên hàm.} Mở rộng chẵn $\widetilde g=(|s|-1)\chi_{[0,1]}(|s|)$. Tích phân từ $0$ cho $s^2/2-s$ ở $s\in[0,1]$, sau đó giữ $-1/2$; phản xạ nguyên hàm lẻ.
\[f_e(s)=0,\qquad H(s)=\operatorname{sgn}(s)(m^2/2-m),\quad m=\min(|s|,1).\]
Ở đây $\operatorname{clip}(s,p,q)=\min(q,\max(p,s))$; $\operatorname{mod}(s,L)\in[0,L)$ kể cả với $s<0$.

\textbf{Bước 3. Xác định hai sóng chạy và ghép nghiệm.} $A=f_e/2-H/(2a)$, $B=f_e/2+H/(2a)$. Thay hai cận:
\[u(x,t)=\frac{f_e(x-1t)+f_e(x+1t)}2+\frac{H(x+1t)-H(x-1t)}{2}.\]

\textbf{Bước 4. Trả lời các yêu cầu và kiểm tra.} Vì $\widetilde g$ liên tục cả ở $0,\pm1$, $u_x$ tồn tại ở mọi $x$ tại $t=1$. Khi $t\ge2$ tại $x=1$, hai cận đã ở hai phía ngoài hỗ trợ, $H(1+t)=-1/2$, $H(1-t)=1/2$, nên $u=-1/2$.

Tại $t=0$, hiệu nguyên hàm bằng $0$ và đạo hàm thời gian cho $H'(x)=g(x)$ trong miền dữ liệu; half-sum cho đúng li độ. Mỗi hạng là sóng của $x\pm at$, nên thỏa PDE ngoài các mặt gãy và theo nghĩa phân phối trên toàn miền. Đồ thị hai sóng ban đầu và các lát cắt thời gian được tính trực tiếp từ các nhánh trên.

%ID BT_3_2_3_08
\textbf{Bước 1. Đọc dữ liệu và chọn phản xạ.} Nửa trục Neumann, $a=3$, $f=s(2-s)$ trên $[0,2]$, $g=6$.

\textbf{Bước 2. Tính từng nhánh và nối nguyên hàm.} Mở rộng chẵn $f=|s|(2-|s|)\chi_{[0,2]}(|s|)$ và $g=6$, nên $H=6s$. Hiệu $H$ góp đúng $6t$.
\[f_e(s)=|s|(2-|s|)\chi_{[0,2]}(|s|),\qquad H(s)=6s.\]
Ở đây $\operatorname{clip}(s,p,q)=\min(q,\max(p,s))$; $\operatorname{mod}(s,L)\in[0,L)$ kể cả với $s<0$.

\textbf{Bước 3. Xác định hai sóng chạy và ghép nghiệm.} $A=f_e/2-H/(2a)$, $B=f_e/2+H/(2a)$. Thay hai cận:
\[u(x,t)=\frac{f_e(x-3t)+f_e(x+3t)}2+\frac{H(x+3t)-H(x-3t)}{6}.\]

\textbf{Bước 4. Trả lời các yêu cầu và kiểm tra.} Ở $t=1/3$, $x-1=0$ là góc của mở rộng chẵn $f$, cho $x=1$; $x-1=2$ cho $x=3$. Các góc ở cận phải không cho điểm dương khác. Sau $t>1$ tại $x=1$, phần li độ hết, $u=6t\to\infty$.

Tại $t=0$, hiệu nguyên hàm bằng $0$ và đạo hàm thời gian cho $H'(x)=g(x)$ trong miền dữ liệu; half-sum cho đúng li độ. Mỗi hạng là sóng của $x\pm at$, nên thỏa PDE ngoài các mặt gãy và theo nghĩa phân phối trên toàn miền. Đồ thị hai sóng ban đầu và các lát cắt thời gian được tính trực tiếp từ các nhánh trên.

%ID BT_3_2_3_09
\textbf{Bước 1. Đọc dữ liệu và chọn phản xạ.} Nửa trục Dirichlet, $a=1$, $f=0$, $g=2s-2$ trên $[0,2]$.

\textbf{Bước 2. Tính từng nhánh và nối nguyên hàm.} Mở rộng lẻ là $\operatorname{sgn}(s)(2|s|-2)\chi_{[0,2]}(|s|)$. Tích phân dương cho $s^2-2s$; tại $2$ bằng $0$, sau đó giữ $0$. Nguyên hàm chẵn.
\[f_e(s)=0,\qquad H(s)=m^2-2m,\quad m=\min(|s|,2).\]
Ở đây $\operatorname{clip}(s,p,q)=\min(q,\max(p,s))$; $\operatorname{mod}(s,L)\in[0,L)$ kể cả với $s<0$.

\textbf{Bước 3. Xác định hai sóng chạy và ghép nghiệm.} $A=f_e/2-H/(2a)$, $B=f_e/2+H/(2a)$. Thay hai cận:
\[u(x,t)=\frac{f_e(x-1t)+f_e(x+1t)}2+\frac{H(x+1t)-H(x-1t)}{2}.\]

\textbf{Bước 4. Trả lời các yêu cầu và kiểm tra.} Tại $t=1$, các cận gặp bước nhảy ở $0,2$ tại $x=1,3$; các bước nhảy khác không cho vị trí dương mới. Sau $t\ge3$ ở $x=1$, cả hai giá trị nguyên hàm đều $0$, nên giới hạn bằng $0$.

Tại $t=0$, hiệu nguyên hàm bằng $0$ và đạo hàm thời gian cho $H'(x)=g(x)$ trong miền dữ liệu; half-sum cho đúng li độ. Mỗi hạng là sóng của $x\pm at$, nên thỏa PDE ngoài các mặt gãy và theo nghĩa phân phối trên toàn miền. Đồ thị hai sóng ban đầu và các lát cắt thời gian được tính trực tiếp từ các nhánh trên.

%ID BT_3_2_3_10
\textbf{Bước 1. Đọc dữ liệu và chọn phản xạ.} Nửa trục Neumann, $a=1$, $f=0$, $g=2-s$ trên $[0,2]$.

\textbf{Bước 2. Tính từng nhánh và nối nguyên hàm.} Mở rộng chẵn là $(2-|s|)\chi_{[0,2]}(|s|)$. Tích phân dương cho $2s-s^2/2$, tại $2$ bằng $2$; nguyên hàm lẻ có hai đuôi $\pm2$.
\[f_e(s)=0,\qquad H(s)=\operatorname{sgn}(s)(2m-m^2/2),\quad m=\min(|s|,2).\]
Ở đây $\operatorname{clip}(s,p,q)=\min(q,\max(p,s))$; $\operatorname{mod}(s,L)\in[0,L)$ kể cả với $s<0$.

\textbf{Bước 3. Xác định hai sóng chạy và ghép nghiệm.} $A=f_e/2-H/(2a)$, $B=f_e/2+H/(2a)$. Thay hai cận:
\[u(x,t)=\frac{f_e(x-1t)+f_e(x+1t)}2+\frac{H(x+1t)-H(x-1t)}{2}.\]

\textbf{Bước 4. Trả lời các yêu cầu và kiểm tra.} $\widetilde g$ liên tục nên tại $t=1$, $u_x$ tồn tại mọi $x>0$. Ở $x=1/2$, khi $t\ge5/2$, hiệu hai đuôi là $4$, chia $2$ cho $u=2$. Góc của $\widetilde g$ tại $0$ ảnh hưởng cấp hai, không ảnh hưởng tồn tại $u_x$.

Tại $t=0$, hiệu nguyên hàm bằng $0$ và đạo hàm thời gian cho $H'(x)=g(x)$ trong miền dữ liệu; half-sum cho đúng li độ. Mỗi hạng là sóng của $x\pm at$, nên thỏa PDE ngoài các mặt gãy và theo nghĩa phân phối trên toàn miền. Đồ thị hai sóng ban đầu và các lát cắt thời gian được tính trực tiếp từ các nhánh trên.

%ID BT_3_2_3_11
\textbf{Bước 1. Đọc dữ liệu và chọn phản xạ.} Nửa trục Neumann $u_x(0,t)=0$, $a=1$, $f=0$, $g=\chi_{[1,2]}$.

\textbf{Bước 2. Tính từng nhánh và nối nguyên hàm.} Biên trong bản scan là Neumann: phải mở rộng chẵn $\widetilde g=\chi_{[1,2]}(|s|)$. Tích phân dương cho $[s-1]_0^1$; nguyên hàm lẻ là $\operatorname{sgn}(s)[|s|-1]_0^1$. Công thức lẻ của vận tốc trong lời giải cũ là nhầm biên.
\[f_e(s)=0,\qquad H(s)=\operatorname{sgn}(s)\operatorname{clip}(|s|-1,0,1).\]
Ở đây $\operatorname{clip}(s,p,q)=\min(q,\max(p,s))$; $\operatorname{mod}(s,L)\in[0,L)$ kể cả với $s<0$.

\textbf{Bước 3. Xác định hai sóng chạy và ghép nghiệm.} $A=f_e/2-H/(2a)$, $B=f_e/2+H/(2a)$. Thay hai cận:
\[u(x,t)=\frac{f_e(x-1t)+f_e(x+1t)}2+\frac{H(x+1t)-H(x-1t)}{2}.\]

\textbf{Bước 4. Trả lời các yêu cầu và kiểm tra.} Vì $H$ lẻ, $H'(-s)=H'(s)$, nên $u_x(0,t)=(\widetilde g(t)-\widetilde g(-t))/2=0$. Nghiệm liên tục; đạo hàm cấp một gãy trên $x\pm t=\pm1,\pm2$ ở phần nằm trong nửa trục.

Tại $t=0$, hiệu nguyên hàm bằng $0$ và đạo hàm thời gian cho $H'(x)=g(x)$ trong miền dữ liệu; half-sum cho đúng li độ. Mỗi hạng là sóng của $x\pm at$, nên thỏa PDE ngoài các mặt gãy và theo nghĩa phân phối trên toàn miền. Đồ thị hai sóng ban đầu và các lát cắt thời gian được tính trực tiếp từ các nhánh trên.

%ID BT_3_2_3_12
\textbf{Bước 1. Đọc dữ liệu và chọn phản xạ.} Nửa trục Neumann, $a=2$, $f=0$, $g=\chi_{[1,2]}$.

\textbf{Bước 2. Tính từng nhánh và nối nguyên hàm.} Mở rộng chẵn và nguyên hàm giống Bài 11. Chỉ tốc độ đổi từ $1$ thành $2$, nên độ dịch là $2t$ và hệ số tích phân là $1/4$.
\[f_e(s)=0,\qquad H(s)=\operatorname{sgn}(s)\operatorname{clip}(|s|-1,0,1).\]
Ở đây $\operatorname{clip}(s,p,q)=\min(q,\max(p,s))$; $\operatorname{mod}(s,L)\in[0,L)$ kể cả với $s<0$.

\textbf{Bước 3. Xác định hai sóng chạy và ghép nghiệm.} $A=f_e/2-H/(2a)$, $B=f_e/2+H/(2a)$. Thay hai cận:
\[u(x,t)=\frac{f_e(x-2t)+f_e(x+2t)}2+\frac{H(x+2t)-H(x-2t)}{4}.\]

\textbf{Bước 4. Trả lời các yêu cầu và kiểm tra.} $u_x(0,t)=0$ bởi tính chẵn của vận tốc mở rộng. $H$ liên tục nên $u$ liên tục; $H'$ nhảy ở $\pm1,\pm2$ nên các đường mất trơn cấp một là $x\pm2t=\pm1,\pm2$. Khi $t>(x+2)/2$, $u=1/2$ tại mỗi điểm cố định.

Tại $t=0$, hiệu nguyên hàm bằng $0$ và đạo hàm thời gian cho $H'(x)=g(x)$ trong miền dữ liệu; half-sum cho đúng li độ. Mỗi hạng là sóng của $x\pm at$, nên thỏa PDE ngoài các mặt gãy và theo nghĩa phân phối trên toàn miền. Đồ thị hai sóng ban đầu và các lát cắt thời gian được tính trực tiếp từ các nhánh trên.

%ID BT_3_2_3_13
\textbf{Bước 1. Đọc dữ liệu và chọn phản xạ.} Nửa trục Dirichlet, $a=2$, $f=0$, $g=s\chi_{[1,2]}$.

\textbf{Bước 2. Tính từng nhánh và nối nguyên hàm.} Mở rộng lẻ; trên đối số dương, $H=0$ trước $1$, $(s^2-1)/2$ trên $[1,2]$, $3/2$ sau $2$. Phản xạ $H$ chẵn, nên dạng gộp dùng $|s|$.
\[f_e(s)=0,\qquad H(s)=\tfrac12(\operatorname{clip}(|s|,1,2)^2-1).\]
Ở đây $\operatorname{clip}(s,p,q)=\min(q,\max(p,s))$; $\operatorname{mod}(s,L)\in[0,L)$ kể cả với $s<0$.

\textbf{Bước 3. Xác định hai sóng chạy và ghép nghiệm.} $A=f_e/2-H/(2a)$, $B=f_e/2+H/(2a)$. Thay hai cận:
\[u(x,t)=\frac{f_e(x-2t)+f_e(x+2t)}2+\frac{H(x+2t)-H(x-2t)}{4}.\]

\textbf{Bước 4. Trả lời các yêu cầu và kiểm tra.} Tại $x=0$, hai giá trị $H$ bằng nhau nên $u=0$. Nghiệm liên tục nhưng các bước nhảy của $g$ ở $\pm1,\pm2$ phát sinh các mặt gãy cấp một. Sau $t>(x+2)/2$, cả hai đuôi của $H$ bằng $3/2$, nên nghiệm trở về $0$.

Tại $t=0$, hiệu nguyên hàm bằng $0$ và đạo hàm thời gian cho $H'(x)=g(x)$ trong miền dữ liệu; half-sum cho đúng li độ. Mỗi hạng là sóng của $x\pm at$, nên thỏa PDE ngoài các mặt gãy và theo nghĩa phân phối trên toàn miền. Đồ thị hai sóng ban đầu và các lát cắt thời gian được tính trực tiếp từ các nhánh trên.

%ID BT_3_2_3_14
\textbf{Bước 1. Đọc dữ liệu và chọn phản xạ.} Nửa trục Neumann, $a=3$, $f=0$, $g=s\chi_{[2,3]}$.

\textbf{Bước 2. Tính từng nhánh và nối nguyên hàm.} Mở rộng chẵn. Trên $s>0$, nguyên hàm là $0$ trước $2$, $(s^2-4)/2$ trên $[2,3]$, $5/2$ sau $3$; phản xạ nguyên hàm lẻ.
\[f_e(s)=0,\qquad H(s)=\tfrac12\operatorname{sgn}(s)(\operatorname{clip}(|s|,2,3)^2-4).\]
Ở đây $\operatorname{clip}(s,p,q)=\min(q,\max(p,s))$; $\operatorname{mod}(s,L)\in[0,L)$ kể cả với $s<0$.

\textbf{Bước 3. Xác định hai sóng chạy và ghép nghiệm.} $A=f_e/2-H/(2a)$, $B=f_e/2+H/(2a)$. Thay hai cận:
\[u(x,t)=\frac{f_e(x-3t)+f_e(x+3t)}2+\frac{H(x+3t)-H(x-3t)}{6}.\]

\textbf{Bước 4. Trả lời các yêu cầu và kiểm tra.} Nghiệm liên tục, $u_x(0,t)=0$, các đường gãy cấp một là $x\pm3t=\pm2,\pm3$. Sau $t>(x+3)/3$, hiệu hai đuôi là $5$, nên $u=5/6$. Điều này khác trường hợp Dirichlet có hai đuôi nguyên hàm bằng nhau.

Tại $t=0$, hiệu nguyên hàm bằng $0$ và đạo hàm thời gian cho $H'(x)=g(x)$ trong miền dữ liệu; half-sum cho đúng li độ. Mỗi hạng là sóng của $x\pm at$, nên thỏa PDE ngoài các mặt gãy và theo nghĩa phân phối trên toàn miền. Đồ thị hai sóng ban đầu và các lát cắt thời gian được tính trực tiếp từ các nhánh trên.

%ID BT_3_2_3_15
\textbf{Bước 1. Đọc dữ liệu và chọn phản xạ.} Đoạn $(0,1)$, Dirichlet hai đầu, $a=1$, $f=0$, $g=1$.

\textbf{Bước 2. Tính từng nhánh và nối nguyên hàm.} Mở rộng lẻ tại $0,1$ cho vận tốc bằng $1$ trên $(0,1)$, $-1$ trên $(1,2)$ rồi chu kỳ $2$. Tích phân cho $H=s$ ở $[0,1]$, $2-s$ ở $[1,2]$; trung bình vận tốc một chu kỳ bằng $0$.
\[f_e(s)=0,\qquad H(s)=1-|\operatorname{mod}(s,2)-1|.\]
Ở đây $\operatorname{clip}(s,p,q)=\min(q,\max(p,s))$; $\operatorname{mod}(s,L)\in[0,L)$ kể cả với $s<0$.

\textbf{Bước 3. Xác định hai sóng chạy và ghép nghiệm.} $A=f_e/2-H/(2a)$, $B=f_e/2+H/(2a)$. Thay hai cận:
\[u(x,t)=\frac{f_e(x-1t)+f_e(x+1t)}2+\frac{H(x+1t)-H(x-1t)}{2}.\]

\textbf{Bước 4. Trả lời các yêu cầu và kiểm tra.} $H(s+2)=H(s)$ nên nghiệm có chu kỳ thời gian $2$. Chiếu $g$ lên sine: $B_n=2\int_0^1\sin(n\pi x)dx=2(1-(-1)^n)/(n\pi)$. Chia $\omega_n=n\pi$ cho chuỗi ở dưới. $g$ khác không ở góc biên nên nghiệm không $C^1$ qua mặt sóng từ các góc.

Tách biến cho $u=\sum_{n\text{ lẻ}}4\sin(n\pi t)\sin(n\pi x)/(n^2\pi^2)$. Dữ liệu $g=1$ tại hai đầu không tương thích nên đây là nghiệm yếu; chu kỳ thời gian 2.

Tại $t=0$, hiệu nguyên hàm bằng $0$ và đạo hàm thời gian cho $H'(x)=g(x)$ trong miền dữ liệu; half-sum cho đúng li độ. Mỗi hạng là sóng của $x\pm at$, nên thỏa PDE ngoài các mặt gãy và theo nghĩa phân phối trên toàn miền. Đồ thị hai sóng ban đầu và các lát cắt thời gian được tính trực tiếp từ các nhánh trên.

\textbf{Chi tiết tích phân của dữ liệu hằng.} Với $k_n=n\pi/L$,
\[\int_0^L\sin k_nx\,dx=\left[-\frac{\cos k_nx}{k_n}\right]_0^L=\frac{1-(-1)^n}{k_n}.\]
Nhân $2/L$ cho hệ số của hằng $1$ là $2(1-(-1)^n)/(n\pi)$; hằng $-1$ đổi dấu. Khi chiếu vận tốc, còn chia $\omega_n$ trong hệ số $\sin\omega_nt$; khi chiếu li độ, hệ số $\cos\omega_nt$ giữ nguyên.

%ID BT_3_2_3_16
\textbf{Bước 1. Đọc dữ liệu và chọn phản xạ.} Đoạn $(0,\pi)$, Neumann hai đầu, $a=2$, $f=x$, $g=0$.

\textbf{Bước 2. Tính từng nhánh và nối nguyên hàm.} Mở rộng chẵn tại $0,\pi$: trên $[0,2\pi]$, $f_e=s$ trước $\pi$ và $2\pi-s$ sau $\pi$. Đó là hàm tam giác chu kỳ $2\pi$; $H=0$.
\[f_e(s)=\pi-|\operatorname{mod}(s,2\pi)-\pi|,\qquad H(s)=0.\]
Ở đây $\operatorname{clip}(s,p,q)=\min(q,\max(p,s))$; $\operatorname{mod}(s,L)\in[0,L)$ kể cả với $s<0$.

\textbf{Bước 3. Xác định hai sóng chạy và ghép nghiệm.} $A=f_e/2-H/(2a)$, $B=f_e/2+H/(2a)$. Thay hai cận:
\[u(x,t)=\frac{f_e(x-2t)+f_e(x+2t)}2+\frac{H(x+2t)-H(x-2t)}{4}.\]

\textbf{Bước 4. Trả lời các yêu cầu và kiểm tra.} $A_0=\pi/2$. Tích phân từng phần cho $A_n=(2/\pi)((-1)^n-1)/n^2$, chỉ còn $n$ lẻ. Ghép với $\omega_n=2n$ cho chuỗi bên dưới và chu kỳ $\pi$. $f'(0)=f'(\pi)=1$ không thỏa biên Neumann, nên có góc phản xạ.

Tách biến cho $u=\pi/2-(4/\pi)\sum_{n\text{ lẻ}}\cos(2nt)\cos(nx)/n^2$. Chu kỳ $\pi$. Dữ liệu $f=x$ không có đạo hàm biên 0, nên không có nghiệm cổ điển tại các góc.

Tại $t=0$, hiệu nguyên hàm bằng $0$ và đạo hàm thời gian cho $H'(x)=g(x)$ trong miền dữ liệu; half-sum cho đúng li độ. Mỗi hạng là sóng của $x\pm at$, nên thỏa PDE ngoài các mặt gãy và theo nghĩa phân phối trên toàn miền. Đồ thị hai sóng ban đầu và các lát cắt thời gian được tính trực tiếp từ các nhánh trên.

\textbf{Chi tiết hệ số cos của dữ liệu $x$.} Trên $(0,\pi)$, trung bình là $A_0=\pi^{-1}\int_0^\pi x\,dx=\pi/2$. Với $n\ge1$,
\[\int_0^\pi x\cos nx\,dx=\left[\frac{x\sin nx}{n}\right]_0^\pi-\frac1n\int_0^\pi\sin nx\,dx
=\frac{(-1)^n-1}{n^2}.\]
Nhân $2/\pi$ cho $2((-1)^n-1)/(\pi n^2)$. Nếu dữ liệu này là li độ thì giữ hệ số với $\cos2nt$; nếu là vận tốc thì chia $2n$ và nhân $\sin2nt$. Vì $n$ lẻ có $(-1)^n-1=-2$, hệ số lần lượt là $-4/(\pi n^2)$ và $-2/(\pi n^3)$.

%ID BT_3_2_3_17
\textbf{Bước 1. Đọc dữ liệu và chọn phản xạ.} Đoạn $(0,1)$, Neumann tại $0$, Dirichlet tại $1$, $a=2$, $f=1-x$, $g=0$.

\textbf{Bước 2. Tính từng nhánh và nối nguyên hàm.} Phản xạ chẵn tại $0$, lẻ tại $1$ cho trên $[0,4]$ hai nhánh $1-s$ ở $[0,2]$, $s-3$ ở $[2,4]$. Vì thế có công thức chính xác $f_e(s)=|\operatorname{mod}(s,4)-2|-1$.
\[f_e(s)=|\operatorname{mod}(s,4)-2|-1,\qquad H(s)=0.\]
Ở đây $\operatorname{clip}(s,p,q)=\min(q,\max(p,s))$; $\operatorname{mod}(s,L)\in[0,L)$ kể cả với $s<0$.

\textbf{Bước 3. Xác định hai sóng chạy và ghép nghiệm.} $A=f_e/2-H/(2a)$, $B=f_e/2+H/(2a)$. Thay hai cận:
\[u(x,t)=\frac{f_e(x-2t)+f_e(x+2t)}2+\frac{H(x+2t)-H(x-2t)}{4}.\]

\textbf{Bước 4. Trả lời các yêu cầu và kiểm tra.} $f_e(s+2)=-f_e(s)$ nên $u(t+1)=-u(t)$, chu kỳ $2$. Với $k_n=(n+1/2)\pi$, tích phân $2\int_0^1(1-x)\cos(k_nx)dx=2/k_n^2$ cho chuỗi. Góc tại phản xạ đầu $0$ xuất phát từ $f'(0)=-1$ không tương thích.

Mô hình $u_{tt}=4u_{xx}$, $u_x(0,t)=0,u(1,t)=0$, $f=1-x,g=0$. Với $k_n=(n+1/2)\pi$, nghiệm $u=\sum_{n\ge0}2\cos(2k_nt)\cos(k_nx)/k_n^2$. Chu kỳ 2; dữ liệu không tương thích đạo hàm tại đầu tự do.

Tại $t=0$, hiệu nguyên hàm bằng $0$ và đạo hàm thời gian cho $H'(x)=g(x)$ trong miền dữ liệu; half-sum cho đúng li độ. Mỗi hạng là sóng của $x\pm at$, nên thỏa PDE ngoài các mặt gãy và theo nghĩa phân phối trên toàn miền. Đồ thị hai sóng ban đầu và các lát cắt thời gian được tính trực tiếp từ các nhánh trên.

\textbf{Chi tiết chiếu li độ lên hệ cos nửa nguyên.} $k_n=(n+1/2)\pi$, $\cos k_n=0$, chuẩn $\int_0^1\cos^2k_nx\,dx=1/2$. Tích phân từng phần:
\[\int_0^1(1-x)\cos k_nx\,dx=\left[\frac{(1-x)\sin k_nx}{k_n}\right]_0^1+\frac1{k_n}\int_0^1\sin k_nx\,dx
=\frac{1-\cos k_n}{k_n^2}=\frac1{k_n^2}.\]
Chia chuẩn $1/2$ cho hệ số $2/k_n^2$; vận tốc đầu bằng $0$ nên không có phần sine theo thời gian.

%ID BT_3_2_3_18
\textbf{Bước 1. Đọc dữ liệu và chọn phản xạ.} Đoạn $(0,1)$, $a=2$, hai biên $u=1$, dữ liệu đầu $f=0$, $g=1$.

\textbf{Bước 2. Tính từng nhánh và nối nguyên hàm.} Lấy hàm nâng $b=1$, $w=u-1$ có biên không, $w_0=-1$, $w_{t,0}=1$. Mở rộng lẻ chu kỳ $2$ của $-1$ và nguyên hàm tam giác của vận tốc được dùng trong công thức; cuối cùng phải cộng lại $1$.
\[f_e(s)=\widetilde f=-1\text{ trên }(0,1),\text{ thác triển lẻ chu kỳ }2,\qquad H(s)=1-|\operatorname{mod}(s,2)-1|.\]
Chọn giá trị mở rộng li độ bằng $0$ tại các số nguyên (trung bình hai vết), để công thức có đúng giá trị biên tại cả thời điểm hai mặt sóng gặp nhau.
Ở đây $\operatorname{clip}(s,p,q)=\min(q,\max(p,s))$; $\operatorname{mod}(s,L)\in[0,L)$ kể cả với $s<0$.

\textbf{Bước 3. Xác định hai sóng chạy và ghép nghiệm.} $A=f_e/2-H/(2a)$, $B=f_e/2+H/(2a)$. Thay hai cận:
\[u(x,t)=\frac{f_e(x-2t)+f_e(x+2t)}2+\frac{H(x+2t)-H(x-2t)}{4}+1.\]

\textbf{Bước 4. Trả lời các yêu cầu và kiểm tra.} Hệ số li độ $A_n=-2(1-(-1)^n)/(n\pi)$, vận tốc $B_n=2(1-(-1)^n)/(n\pi)$, tần số $2n\pi$. Dữ liệu li độ và biên mâu thuẫn ngay ở góc nên không có nghiệm cổ điển trên bao đóng; chuỗi và sóng chạy cho nghiệm trong miền với các bước nhảy.

Chọn $v=1$, $w=u-1$ có biên 0, $w(x,0)=-1,w_t(x,0)=1$. Nghiệm $u=1+\sum_{n\text{ lẻ}}[-4\cos(2n\pi t)/(n\pi)+2\sin(2n\pi t)/(n^2\pi^2)]\sin(n\pi x)$. Góc dữ liệu $u(x,0)=0$ mâu thuẫn biên 1: chỉ có nghiệm yếu với mặt sóng.

Tại $t=0$, hiệu nguyên hàm bằng $0$ và đạo hàm thời gian cho $H'(x)=g(x)$ trong miền dữ liệu; half-sum cho đúng li độ sau khi cộng hàm nâng biên. Mỗi hạng là sóng của $x\pm at$, nên thỏa PDE ngoài các mặt gãy và theo nghĩa phân phối trên toàn miền. Đồ thị hai sóng ban đầu và các lát cắt thời gian được tính trực tiếp từ các nhánh trên.

\textbf{Chi tiết tích phân của dữ liệu hằng.} Với $k_n=n\pi/L$,
\[\int_0^L\sin k_nx\,dx=\left[-\frac{\cos k_nx}{k_n}\right]_0^L=\frac{1-(-1)^n}{k_n}.\]
Nhân $2/L$ cho hệ số của hằng $1$ là $2(1-(-1)^n)/(n\pi)$; hằng $-1$ đổi dấu. Khi chiếu vận tốc, còn chia $\omega_n$ trong hệ số $\sin\omega_nt$; khi chiếu li độ, hệ số $\cos\omega_nt$ giữ nguyên.

%ID BT_3_2_3_19
\textbf{Bước 1. Đọc dữ liệu và chọn phản xạ.} Đoạn $(0,\pi)$, $a=2$, $u_x(0,t)=u_x(\pi,t)=1$, $f=x$, $g=x$.

\textbf{Bước 2. Tính từng nhánh và nối nguyên hàm.} Lấy $b=x$. Phần $w=u-x$ có li độ $0$, vận tốc $x$ và biên Neumann không. Mở rộng chẵn vận tốc thành tam giác $q(s)=\pi-|\operatorname{mod}(s,2\pi)-\pi|$. Tích phân từng nhánh cho $H$ ở dưới, mỗi nửa chu kỳ góp $\pi^2/2$.
\[f_e(s)=0,\qquad H(s)=k\pi^2/2+J(r),\ r=s-k\pi,\ k=\lfloor s/\pi\rfloor;\ J(r)=r^2/2\ (k\text{ chẵn}),\ \pi r-r^2/2\ (k\text{ lẻ}).\]
Ở đây $\operatorname{clip}(s,p,q)=\min(q,\max(p,s))$; $\operatorname{mod}(s,L)\in[0,L)$ kể cả với $s<0$.

\textbf{Bước 3. Xác định hai sóng chạy và ghép nghiệm.} $A=f_e/2-H/(2a)$, $B=f_e/2+H/(2a)$. Thay hai cận:
\[u(x,t)=\frac{f_e(x-2t)+f_e(x+2t)}2+\frac{H(x+2t)-H(x-2t)}{4}+x.\]

\textbf{Bước 4. Trả lời các yêu cầu và kiểm tra.} Trung bình vận tốc là $\pi/2$, sinh hạng $\pi t/2$. Các mode còn lại có $B_n=2(( -1)^n-1)/(\pi n^2)$, chia $2n$ cho hệ số thời gian. Nghiệm không tuần hoàn; tương thích cấp hai ở góc không đúng vì $g'=1$ nhưng đạo hàm thời gian của biên bằng $0$.

Chọn $v=x$. $w=u-x$ có biên Neumann 0, dữ liệu $w=0,w_t=x$. Nghiệm $u=x+\pi t/2-(2/\pi)\sum_{n\text{ lẻ}}\sin(2nt)\cos(nx)/n^3$. Có mode tuyến tính nên không tuần hoàn.

Tại $t=0$, hiệu nguyên hàm bằng $0$ và đạo hàm thời gian cho $H'(x)=g(x)$ trong miền dữ liệu; half-sum cho đúng li độ sau khi cộng hàm nâng biên. Mỗi hạng là sóng của $x\pm at$, nên thỏa PDE ngoài các mặt gãy và theo nghĩa phân phối trên toàn miền. Đồ thị hai sóng ban đầu và các lát cắt thời gian được tính trực tiếp từ các nhánh trên.

\textbf{Chi tiết hệ số cos của dữ liệu $x$.} Trên $(0,\pi)$, trung bình là $A_0=\pi^{-1}\int_0^\pi x\,dx=\pi/2$. Với $n\ge1$,
\[\int_0^\pi x\cos nx\,dx=\left[\frac{x\sin nx}{n}\right]_0^\pi-\frac1n\int_0^\pi\sin nx\,dx
=\frac{(-1)^n-1}{n^2}.\]
Nhân $2/\pi$ cho $2((-1)^n-1)/(\pi n^2)$. Nếu dữ liệu này là li độ thì giữ hệ số với $\cos2nt$; nếu là vận tốc thì chia $2n$ và nhân $\sin2nt$. Vì $n$ lẻ có $(-1)^n-1=-2$, hệ số lần lượt là $-4/(\pi n^2)$ và $-2/(\pi n^3)$.

%ID BT_3_2_3_20
\textbf{Bước 1. Đọc dữ liệu và chọn phản xạ.} Đoạn $(0,\pi)$, Dirichlet hai đầu, $a=2$, $f=0$, vận tốc $g$; minh họa $g=\min(x,\pi-x)$.

\textbf{Bước 2. Tính từng nhánh và nối nguyên hàm.} Với $g$ tổng quát, mở rộng lẻ chu kỳ $2\pi$. Với tam giác cụ thể, tích phân lần lượt trên $[0,\pi/2]$ và $[\pi/2,\pi]$, rồi phản xạ nguyên hàm chẵn và tuần hoàn, được $H$ dưới đây.
\[f_e(s)=0,\qquad H(s)=J(r),\quad r=\operatorname{mod}(s,2\pi).\]
\[J(r)=\begin{cases}r^2/2&0\le r\le\pi/2,\\\pi r-r^2/2-\pi^2/4&\pi/2<r\le\pi,\\J(2\pi-r)&\pi<r<2\pi.\end{cases}\]
Ở đây $\operatorname{clip}(s,p,q)=\min(q,\max(p,s))$; $\operatorname{mod}(s,L)\in[0,L)$ kể cả với $s<0$.

\textbf{Bước 3. Xác định hai sóng chạy và ghép nghiệm.} $A=f_e/2-H/(2a)$, $B=f_e/2+H/(2a)$. Thay hai cận:
\[u(x,t)=\frac{f_e(x-2t)+f_e(x+2t)}2+\frac{H(x+2t)-H(x-2t)}{4}.\]

\textbf{Bước 4. Trả lời các yêu cầu và kiểm tra.} Lấy đạo hàm ở biên: $u_x(0,t)=\widetilde g(2t)/2$. Đổi biến $s=2t$ cho $\int_0^{\pi/2}|u_x(0,t)|^2dt=\frac18\int_0^\pi g(s)^2ds$. Tích phân trên thời gian lớn hơn chỉ thêm số không âm, chứng minh yêu cầu với $T\ge\pi/2$. Hai tích phân của tam giác cho $B_n=4\sin(n\pi/2)/(\pi n^2)$, nên chia $2n$ cho chuỗi.

Mô hình biên cố định $u_{tt}=4u_{xx}$ trên $(0,\pi)$, $f=0,u_t=g$. $u_x(0,t)=\widetilde g(2t)/2$. Do đó $\int_0^{\pi/2}|u_x(0,t)|^2dt=\frac18\int_0^\pi g^2$, suy bất đẳng thức với $T\ge\pi/2$. Với $g=\min(x,\pi-x)$, $u=\sum_{n\ge1}2\sin(n\pi/2)\sin(2nt)\sin(nx)/(\pi n^3)$.

Tại $t=0$, hiệu nguyên hàm bằng $0$ và đạo hàm thời gian cho $H'(x)=g(x)$ trong miền dữ liệu; half-sum cho đúng li độ. Mỗi hạng là sóng của $x\pm at$, nên thỏa PDE ngoài các mặt gãy và theo nghĩa phân phối trên toàn miền. Đồ thị hai sóng ban đầu và các lát cắt thời gian được tính trực tiếp từ các nhánh trên.

\textbf{Chi tiết hệ số của tam giác.} Với $f(x)=\min(x,\pi-x)$, chia tại $b=\pi/2$. Đổi $y=\pi-x$ ở nửa sau cho
\[\int_0^\pi f(x)\sin nx\,dx=(1+(-1)^{n+1})\int_0^b x\sin nx\,dx.\]
Tích phân cuối bằng $[-x\cos nx/n+\sin nx/n^2]_0^b=-b\cos(nb)/n+\sin(nb)/n^2$. Nếu $n$ chẵn, hệ số trước tích phân bằng $0$; nếu $n$ lẻ, $\cos(nb)=0$. Nhân $2/\pi$ được $4\sin(n\pi/2)/(\pi n^2)$.

%ID BT_3_2_3_21
\textbf{Bước 1. Đọc dữ liệu và chọn phản xạ.} Đoạn $(0,1)$, Dirichlet tại $0$, Neumann tại $1$, $a=1$, $f=0$, vận tốc $g$; minh họa $g=1$.

\textbf{Bước 2. Tính từng nhánh và nối nguyên hàm.} Mở rộng lẻ tại $0$ và chẵn tại $1$ cho chu kỳ $4$, đổi dấu sau $2$. Với $g=1$, vận tốc mở rộng là $1$ trên $(0,2)$, $-1$ trên $(2,4)$. Nguyên hàm có nhánh $s$ rồi $4-s$; thêm hằng số không đổi nghiệm.
\[f_e(s)=0,\qquad H(s)=1-|\operatorname{mod}(s,4)-2|.\]
Ở đây $\operatorname{clip}(s,p,q)=\min(q,\max(p,s))$; $\operatorname{mod}(s,L)\in[0,L)$ kể cả với $s<0$.

\textbf{Bước 3. Xác định hai sóng chạy và ghép nghiệm.} $A=f_e/2-H/(2a)$, $B=f_e/2+H/(2a)$. Thay hai cận:
\[u(x,t)=\frac{f_e(x-1t)+f_e(x+1t)}2+\frac{H(x+1t)-H(x-1t)}{2}.\]

\textbf{Bước 4. Trả lời các yêu cầu và kiểm tra.} Với $g$ tổng quát, $u_x(0,t)=\widetilde g(t)=g(t)$ khi $0<t<1$. Do đó $\int_0^1g^2=\int_0^1|u_x(0,t)|^2dt\le\int_0^T|u_x(0,t)|^2dt$ với $T\ge1$. Với $g=1$, công thức dưới có chu kỳ $4$ và đổi dấu sau $2$; không có chu kỳ $2$ cho nghiệm khác không.

Tại $t=0$, hiệu nguyên hàm bằng $0$ và đạo hàm thời gian cho $H'(x)=g(x)$ trong miền dữ liệu; half-sum cho đúng li độ. Mỗi hạng là sóng của $x\pm at$, nên thỏa PDE ngoài các mặt gãy và theo nghĩa phân phối trên toàn miền. Đồ thị hai sóng ban đầu và các lát cắt thời gian được tính trực tiếp từ các nhánh trên.

%ID BT_3_2_3_22
\textbf{Bước 1. Đọc dữ liệu và chọn phản xạ.} Đoạn $(0,1)$, Dirichlet hai đầu, $a=1$, $g=0$, $f=4x$ trên $[0,1/4]$, $2-4x$ trên $[1/4,1/2]$, $0$ trên $[1/2,1]$.

\textbf{Bước 2. Tính từng nhánh và nối nguyên hàm.} Mở rộng lẻ tại hai đầu và chu kỳ $2$. Vì vận tốc bằng không, hai sóng chạy đều bằng một nửa $f_e$; không cần tích phân. Công thức này mô tả cả phản xạ tại các thời gian lớn hơn.
\[f_e(s)=\widetilde f\text{ là thác triển lẻ chu kỳ }2\text{ của tam giác trong đề},\qquad H(s)=0.\]
Ở đây $\operatorname{clip}(s,p,q)=\min(q,\max(p,s))$; $\operatorname{mod}(s,L)\in[0,L)$ kể cả với $s<0$.

\textbf{Bước 3. Xác định hai sóng chạy và ghép nghiệm.} $A=f_e/2-H/(2a)$, $B=f_e/2+H/(2a)$. Thay hai cận:
\[u(x,t)=\frac{f_e(x-1t)+f_e(x+1t)}2+\frac{H(x+1t)-H(x-1t)}{2}.\]

\textbf{Bước 4. Trả lời các yêu cầu và kiểm tra.} Với $1/2<x<1$, trước $t=x-1/2$, cả hai cận đều ở vùng dữ liệu không hoặc chưa có sóng phản xạ, nên $u=0$. Ngay sau thời điểm ấy, cận trái vào phần tam giác dương, nên $u\ne0$. Đó là thời điểm đến (giá trị đúng tại thời điểm đến vẫn $0$). Suốt $0\le t\le1/4$, các điểm cân bằng là $[3/4,1]$ cùng đầu cố định $0$. Với tốc độ $c$, thời điểm đến là $(x-1/2)/c$; vùng cân bằng suốt khoảng thời gian này là $[1/2+c/4,1]$ nếu cận không vượt $1$, cộng hai đầu cố định. Nếu $c\ge2$ chỉ còn hai đầu. Nếu đồng thời đổi khoảng thời gian thành $[0,1/(4c)]$, vùng giữ nguyên là $[3/4,1]$ cùng $0$.

Tại $t=0$, hiệu nguyên hàm bằng $0$ và đạo hàm thời gian cho $H'(x)=g(x)$ trong miền dữ liệu; half-sum cho đúng li độ. Mỗi hạng là sóng của $x\pm at$, nên thỏa PDE ngoài các mặt gãy và theo nghĩa phân phối trên toàn miền. Đồ thị hai sóng ban đầu và các lát cắt thời gian được tính trực tiếp từ các nhánh trên.

%ID VD_3_3_2_01
\textbf{Bước 1. Chọn sine theo biên cố định.} Trên $(0,L)$, $X_n=\sin(n\pi x/L)$, $\omega_n=an\pi/L$. Hệ số của li độ là $A_n$, của vận tốc là $B_n$; nghiệm mode là $(A_n\cos\omega_nt+B_n\sin\omega_nt/\omega_n)X_n$.

\textbf{Bước 2. Câu (i).} $a=L=1$, $f=\sin\pi x$, $g=\sin^3\pi x=(3\sin\pi x-\sin3\pi x)/4$. Trực giao cho $A_1=1$, $B_1=3/4$, $B_3=-1/4$, các hệ số khác bằng $0$. Do đó
\[u=\sin\pi x\cos\pi t+\frac3{4\pi}\sin\pi x\sin\pi t-\frac1{12\pi}\sin3\pi x\sin3\pi t.\]
Đạo hàm thời gian tại $0$ khôi phục $g$; từng mode thỏa PDE và biên. Nghiệm $C^\infty$, chu kỳ cơ bản $2$.

\textbf{Bước 3. Câu (ii).} $a=2$, $L=\pi$, dữ liệu trong sách là $f(x)=\sin(\pi x)$, không phải $\sin x$. Dùng $2\sin\pi x\sin nx=\cos((\pi-n)x)-\cos((\pi+n)x)$:
\[A_n=\frac2\pi\int_0^\pi\sin(\pi x)\sin nx\,dx
=\frac{2(-1)^nn\sin(\pi^2)}{\pi(\pi^2-n^2)},\qquad B_n=0.\]
Không có mẫu số bằng không vì $n$ nguyên. Nghiệm là $u=\sum_{n\ge1}A_n\cos2nt\sin nx$, chu kỳ $\pi$.

\textbf{Bước 4. Nêu mức trơn đúng.} $f(\pi)=\sin(\pi^2)\ne0$, trái biên $u(\pi,t)=0$. Do đó câu (ii) không có nghiệm liên tục đến góc $(\pi,0)$ với mọi dữ liệu nguyên văn. Chuỗi xác định nghiệm theo nghĩa $L^2$/ từng miền; bước nhảy của mở rộng lẻ phát sinh mặt sóng từ đầu phải. Câu (i) không có vấn đề tương thích ấy.

\textbf{Chi tiết tích phân câu (ii).} Từ tích–tổng,
\[A_n=\frac1\pi\left[\frac{\sin((\pi-n)x)}{\pi-n}-\frac{\sin((\pi+n)x)}{\pi+n}\right]_0^\pi.\]
Vì $\sin(\pi^2-n\pi)=\sin(\pi^2+n\pi)=(-1)^n\sin\pi^2$,
\[A_n=\frac{(-1)^n\sin\pi^2}{\pi}\left(\frac1{\pi-n}-\frac1{\pi+n}\right)
=\frac{2(-1)^nn\sin\pi^2}{\pi(\pi^2-n^2)}.\]

%ID VD_3_3_2_02
\textbf{Bước 1. Giữ mode hằng của Neumann.} Hệ hàm riêng là $X_0=1$, $X_n=\cos(n\pi x/L)$, $n\ge1$. Mode $0$ bằng $A_0+B_0t$ với hệ số là trung bình của dữ liệu.

\textbf{Bước 2. Câu (i).} $a=L=1$, $f=0$, $g=\cos^2\pi x=(1+\cos2\pi x)/2$. Vì vậy $B_0=1/2$, $B_2=1/2$, còn lại bằng $0$:
\[u=t/2+\cos2\pi x\sin2\pi t/(4\pi).\]
Đạo hàm không gian là bội của $\sin2\pi x$, bằng không ở hai đầu. Đạo hàm thời gian tại $0$ cho đúng $g$. Hạng tuyến tính khiến nghiệm không tuần hoàn dù phần dao động có chu kỳ $1$.

\textbf{Bước 3. Câu (ii).} $a=2$, $L=\pi$, $f=\sin x$, $g=0$. $A_0=2/\pi$. Tích–tổng cho
\[A_n=\frac2\pi\int_0^\pi\sin x\cos nx\,dx
=\frac{2(1+(-1)^n)}{\pi(1-n^2)}\quad(n\ne1),\qquad A_1=0.\]
Tại $n=1$ phải tích phân riêng, không dùng mẫu bằng $0$. Chỉ còn các mode chẵn:
\[u=\frac2\pi+\sum_{m\ge1}\frac4{\pi(1-4m^2)}\cos4mt\cos2mx.\]
\textbf{Bước 4. Chu kỳ và trơn.} Mode $m=1$ khác không cho chu kỳ cơ bản $\pi/2$. Dữ liệu $f'(0)=1$, $f'(\pi)=-1$ không thỏa biên Neumann; mở rộng chẵn có góc, nên không có nghiệm $C^1$ trên toàn miền qua các mặt phản xạ. Chuỗi và công thức phản xạ xác định cùng nghiệm liên tục từng miền.

\textbf{Chi tiết tích–tổng và ngoại lệ $n=1$.} Với $n\ne1$,
\[\begin{aligned}
\int_0^\pi\sin x\cos nx\,dx
&=\frac12\int_0^\pi[\sin((1+n)x)+\sin((1-n)x)]dx\\
&=\frac{1+(-1)^n}{2}\left(\frac1{1+n}+\frac1{1-n}\right)
=\frac{1+(-1)^n}{1-n^2}.
\end{aligned}\]
Nhân $2/\pi$ được $A_n$. Với $n=1$, tích phân riêng $\int_0^\pi\sin x\cos x\,dx=[\sin^2x/2]_0^\pi=0$. Mode hằng dùng chuẩn $\pi$, nên $A_0=\pi^{-1}[-\cos x]_0^\pi=2/\pi$, không dùng chuẩn $\pi/2$ của các mode khác.

%ID VD_3_3_2_03
\textbf{Bước 1. Trị riêng Dirichlet–Neumann.} $X(0)=0$, $X'(L)=0$. Tích phân từng phần cho $\lambda\int X^2=\int(X')^2>0$ vì nghiệm hằng thỏa $X(0)=0$ chỉ là $0$. Với $\lambda=k^2$, $X=C\sin kx$, $\cos kL=0$, nên $k_n=(n+1/2)\pi/L$, $n\ge0$.

\textbf{Bước 2. Trực giao và hệ số.} Với hai trị riêng, trừ $X_m''X_n-X_n''X_m$ và tích phân. Hạng biên $[X_m'X_n-X_n'X_m]_0^L=0$, nên tích vô hướng bằng $0$ nếu khác trị riêng. Chuẩn là $L/2$, suy $A_n=(2/L)\int f\sin k_nx$, $B_n=(2/L)\int g\sin k_nx$.

\textbf{Bước 3. Ghép nghiệm.} $u=\sum_{n\ge0}(A_n\cos ak_nt+B_n\sin ak_nt/(ak_n))\sin k_nx$. Mỗi mode đổi dấu sau $2L/a$ và có chu kỳ $4L/a$. Tính duy nhất theo năng lượng: thông lượng ở $0$ hết vì $u_t=0$, ở $L$ hết vì $u_x=0$.

\textbf{Bước 4. Dữ liệu cụ thể.} $a=1$, $L=\pi$, $f=x$, $g=\sin^3(x/2)$. Với $k_n=n+1/2$, tích phân từng phần cho $\int_0^\pi x\sin k_nx=(-1)^n/k_n^2$, nên $A_n=2(-1)^n/(\pi k_n^2)$. Từ công thức sin lập phương, $B_0=3/4$, $B_1=-1/4$:
\[u=\sum_{n\ge0}\frac{2(-1)^n}{\pi k_n^2}\cos k_nt\sin k_nx
+\frac32\sin(t/2)\sin(x/2)-\frac16\sin(3t/2)\sin(3x/2).\]
Chu kỳ cơ bản $4\pi$. $f'(\pi)=1$ trái biên Neumann, nên nghiệm không $C^1$ qua mặt sóng từ góc phải; đây là nghiệm năng lượng/từng miền, không nghiệm cổ điển toàn bao đóng.

\textbf{Chi tiết điều kiện biên, chuẩn và tích phân dữ liệu.} Từ $X''+k^2X=0$, $X=A\cos kx+B\sin kx$. $X(0)=0$ cho $A=0$; $X'(L)=Bk\cos kL=0$ với $Bk\ne0$ cho $\cos kL=0$. Chuẩn
\[\int_0^L\sin^2k_nx\,dx=\left[x/2-\sin(2k_nx)/(4k_n)\right]_0^L=L/2.\]
Với $L=\pi$,
\[\int_0^\pi x\sin k_nx\,dx=
\left[-x\cos(k_nx)/k_n+\sin(k_nx)/k_n^2\right]_0^\pi
=(-1)^n/k_n^2.\]
Đã dùng $\cos(k_n\pi)=0$, $\sin(k_n\pi)=(-1)^n$. Khai triển $g=3\sin(x/2)/4-\sin(3x/2)/4$ rồi chiếu lên hệ sine trực giao cho $B_0=3/4,B_1=-1/4$.

\textbf{Chi tiết trực giao và dấu trị riêng.} Nhân $-X''=\lambda X$ bởi $X$, tích phân:
\[\lambda\int_0^LX^2dx=-[X'X]_0^L+\int_0^L(X')^2dx=\int_0^L(X')^2dx.\]
Hạng biên hết vì mỗi đầu có $X=0$ hoặc $X'=0$. Nếu $\lambda=0$, $\int(X')^2=0$ cho $X$ hằng; khi có một đầu Dirichlet, hằng đó bằng $0$. Với hai trị riêng,
\[(\lambda_m-\lambda_n)\int_0^LX_mX_n\,dx
=\int_0^L(-X_m''X_n+X_n''X_m)dx
=[-X_m'X_n+X_n'X_m]_0^L=0.\]
Nếu $\lambda_m\ne\lambda_n$, chia hiệu trị riêng cho tích vô hướng bằng $0$.

%ID BT_3_3_3_01
\textbf{Bước 1. Công thức chung.} Hai biên cố định cho $X_n=\sin(n\pi x/L)$, $\omega_n=an\pi/L$ và hệ số $A_n=(2/L)\int fX_n$, $B_n=(2/L)\int gX_n$. Dùng $u=\sum(A_n\cos\omega_nt+B_n\sin\omega_nt/\omega_n)X_n$.

\textbf{Bước 2. Tính hệ số từng dữ liệu.}

(i) $a=L=1$. $A_1=1$, còn lại $0$. Hai lần tích phân từng phần cho
\[\int_0^1x(1-x)\sin(n\pi x)dx=\frac{2(1-(-1)^n)}{n^3\pi^3}.\]
Do đó $B_n=4(1-(-1)^n)/(n^3\pi^3)$, và
\[u=\sin\pi x\cos\pi t+\sum_{n\text{ lẻ}}\frac8{n^4\pi^4}\sin(n\pi t)\sin(n\pi x).\]
Chu kỳ $2$; dữ liệu tương thích cho $C^2$.

(ii) $a=2$, $L=\pi$, $f=\sin x$, $g=1$. $A_1=1$; $B_n=(2/\pi)\int_0^\pi\sin nx=2(1-(-1)^n)/(\pi n)$. Suy $u=\sin x\cos2t+\sum_{n\text{ lẻ}}2\sin2nt\sin nx/(\pi n^2)$, chu kỳ $\pi$. Vận tốc không tương thích ở hai đầu, nên có gãy cấp một.

(iii) $a=1$, $L=\pi$, $f=1$, $g=0$. Cùng tích phân trên cho $A_n=2(1-(-1)^n)/(\pi n)$, nên $u=\sum_{n\text{ lẻ}}4\cos nt\sin nx/(\pi n)$. Chu kỳ $2\pi$; li độ không tương thích biên tạo bước nhảy.

(iv) $a=L=1$, $f=x(1-x)/4$, $g=0$. Chia hệ số của đa thức ở (i) cho $4$: $A_n=(1-(-1)^n)/(n^3\pi^3)$. Vậy $u=\sum_{n\text{ lẻ}}2\cos(n\pi t)\sin(n\pi x)/(n^3\pi^3)$, chu kỳ $2$. $f''(0)=f''(1)=-1/2$ không khớp biên cấp hai, nên không $C^2$ qua mặt phản xạ, dù nghiệm $C^1$.

(v) Cùng $f$, $g=-1$: giữ $A_n$ trên, thêm $B_n=-2(1-(-1)^n)/(n\pi)$. Nghiệm theo công thức chung, chu kỳ $2$, không $C^1$ qua mặt sóng từ góc vì vận tốc biên khác không.

(vi) $f=0$, $g=-1/2$ trên $(0,1/2)$, $1/2$ trên $(1/2,1)$. Tách tích phân tại $1/2$:
\[B_n=-\int_0^{1/2}\sin(n\pi x)dx+\int_{1/2}^1\sin(n\pi x)dx
=\frac{2\cos(n\pi/2)-1-(-1)^n}{n\pi}.\]
Chỉ $n\equiv2\pmod4$ khác không, với $B_n=-4/(n\pi)$. Nghiệm $u=-\sum_{n\equiv2\ (4)}4\sin(n\pi t)\sin(n\pi x)/(n^2\pi^2)$ có chu kỳ cơ bản $1$. Vận tốc bước nhảy cho nghiệm liên tục nhưng gãy cấp một.

\textbf{Bước 3. Kiểm tra.} Mỗi mode thỏa PDE và biên, hai chuỗi tại $t=0$ đúng các hệ số đã tính. Chỉ đạo hàm chuỗi ở nơi hội tụ đủ; với dữ liệu không tương thích dùng phản xạ để xác định các mặt gãy, không khẳng định cổ điển toàn miền.

\textbf{Chi tiết hai lần tích phân từng phần ở (i).} Đặt $k=n\pi$, $P=x(1-x)$, $P'=1-2x$, $P''=-2$.
\[\begin{aligned}
I&=\int_0^1P\sin kx\,dx
=\left[-P\frac{\cos kx}{k}\right]_0^1+\frac1k\int_0^1P'\cos kx\,dx\\
&=\frac1k\left[\frac{P'\sin kx}{k}\right]_0^1-\frac1{k^2}\int_0^1P''\sin kx\,dx\\
&=\frac2{k^2}\left[-\frac{\cos kx}{k}\right]_0^1=\frac{2(1-(-1)^n)}{k^3}.
\end{aligned}\]
Hạng biên lần đầu hết vì $P(0)=P(1)=0$; lần sau hết vì $\sin0=\sin n\pi=0$. Hệ số vận tốc là $B_n=2I$; khi ghép nghiệm còn chia $\omega_n=k$. Với $n$ lẻ, $2I/k=8/k^4$, không bỏ hệ số chia tần số này.

Ở (ii),(iii),(v), dùng $\int_0^L\sin(n\pi x/L)dx=L(1-(-1)^n)/(n\pi)$, rồi nhân $2/L$ và giá trị hằng của dữ liệu. Ở (iv), $\int fX_n=I/4$ nên $A_n=I/2=(1-(-1)^n)/k^3$. Ở (vi),
\[-\int_0^{1/2}\sin kx\,dx=\frac{\cos(k/2)-1}{k},\quad
\int_{1/2}^1\sin kx\,dx=\frac{\cos(k/2)-\cos k}{k}.\]
Cộng rồi xét $n$ lẻ, $n\equiv0\pmod4$, $n\equiv2\pmod4$ cho lần lượt $0,0,-4/k$.

%ID BT_3_3_3_02
\textbf{Bước 1. Chọn hệ sine.} Trên $(0,\pi)$, $f=\min(x,\pi-x)$, $g=0$, $a=0.0002$. Tách tích phân tại $\pi/2$ và đổi biến $y=\pi-x$ ở nửa sau:
\[A_n=\frac2\pi(1+(-1)^{n+1})\int_0^{\pi/2}x\sin nx\,dx.\]
Với $n$ chẵn bằng $0$. Với $n$ lẻ, $\cos(n\pi/2)=0$, tích phân bằng $\sin(n\pi/2)/n^2$, nên chung $A_n=4\sin(n\pi/2)/(\pi n^2)$.

\textbf{Bước 2. Giải ODE.} $T_n''+a^2n^2T_n=0$, $T_n(0)=A_n$, $T_n'(0)=0$, cho
\[u=\sum_{n\ge1}\frac{4\sin(n\pi/2)}{\pi n^2}\cos(0.0002nt)\sin nx.\]
\textbf{Bước 3. Chu kỳ và tính trơn.} Mode $1$ khác không, mọi tần số là bội nguyên của $a$, nên chu kỳ cơ bản $2\pi/a=10000\pi$. Li độ có góc tại $\pi/2$; hai góc chạy theo $x=\pi/2\pm at$ rồi phản xạ ở hai đầu. Nghiệm liên tục/ năng lượng, không $C^1$ trên các nhánh truyền góc. Thế $t=0$ thu lại khai triển tam giác; biên đúng theo từng mode.

\textbf{Chi tiết hai tích phân của tam giác.} Đặt $k=n$, $b=\pi/2$. Đổi $y=\pi-x$ ở nửa sau:
\[\int_b^\pi(\pi-x)\sin nx\,dx=(-1)^{n+1}\int_0^b y\sin ny\,dy.\]
Tích phân từng phần còn lại là
\[\int_0^b x\sin nx\,dx=\left[-x\cos nx/n+\sin nx/n^2\right]_0^b
=-b\cos(nb)/n+\sin(nb)/n^2.\]
Với $n$ chẵn, hai nửa triệt tiêu; với $n$ lẻ, $\cos(nb)=0$ và hệ số trước tích phân là $4/\pi$. Suy $A_n=4\sin(n\pi/2)/(\pi n^2)$ cho mọi $n$.

%ID BT_3_3_3_03
\textbf{Bước 1. Biên Dirichlet–Neumann.} Tách biến cho $k_n=(n+1/2)\pi/L$. Khi dịch $2L/a$, pha tăng $(2n+1)\pi$; cos và sin đều đổi dấu, nên $u(t+2L/a)=-u(t)$. Dịch thêm lần nữa cho chu kỳ $4L/a$.

\textbf{Bước 2. Kiểm tra khẳng định chu kỳ ngắn.} Nếu còn $u(t+2L/a)=u(t)$ thì $u=-u$, suy $u=0$. Lấy $f=\sin(\pi x/(2L))$, $g=0$ cho $u=f\cos(a\pi t/(2L))$ là phản ví dụ rõ ràng với chu kỳ cơ bản $4L/a$. Vì vậy chu kỳ $2L/a$ trong câu (a) không đúng cho nghiệm khác không.

\textbf{Bước 3. Biên Neumann hai đầu.} Tích phân PDE theo $x$ cho trung bình $\bar u''=(a^2/L)[u_x]_0^L=0$, nên $\bar u=\bar f+t\bar g$. Phần còn lại gồm các tần số $an\pi/L$, tất cả có chu kỳ $2L/a$. Nếu $\bar g=0$ thì nghiệm tuần hoàn; nếu $\bar g\ne0$, trung bình tăng tuyến tính và không thể tuần hoàn. Điều kiện cần và đủ là $\int_0^Lg=0$. Nghiệm hằng theo thời gian có mọi chu kỳ dương.

%ID BT_3_3_3_04
\textbf{Bước 1. Hệ cos và mode hằng.} $A_0,B_0$ là trung bình dữ liệu; mode khác không có chuẩn $L/2$. Không được dùng công thức sine của biên cố định.

\textbf{Bước 2. Câu (i).} $a=L=1$, $f=\cos2\pi x$, $g=x-1/2$. $A_2=1$, $B_0=0$. Tích phân từng phần cho
\[B_n=2\int_0^1(x-1/2)\cos(n\pi x)dx=\frac{2((-1)^n-1)}{n^2\pi^2}.\]
Do đó $u=\cos2\pi x\cos2\pi t-\sum_{n\text{ lẻ}}4\sin(n\pi t)\cos(n\pi x)/(n^3\pi^3)$. Mode $1$ khác không cho chu kỳ cơ bản $2$. $g'=1$ ở hai đầu trái đạo hàm thời gian của biên Neumann, nên không $C^2$ qua các mặt sóng góc.

\textbf{Bước 3. Câu (ii).} $a=2$, $L=\pi$, $f=\sin^2x=(1-\cos2x)/2$, $g=\sin x$. $A_0=1/2$, $A_2=-1/2$, $B_0=2/\pi$. Tích–tổng cho $B_n=2(1+(-1)^n)/(\pi(1-n^2))$ nếu $n\ne1$, và $B_1=0$ tính riêng. Suy
\[u=\tfrac12-\tfrac12\cos2x\cos4t+2t/\pi+\sum_{n\ge1}\frac{B_n}{2n}\sin2nt\cos nx.\]
\textbf{Bước 4. Kiểm tra.} Mode hằng đảm bảo vận tốc trung bình đúng; bỏ nó sẽ làm sai $g$. Nó cũng khiến câu (ii) không tuần hoàn. Biên đúng theo mỗi cos; tại $t=0$ hai chuỗi khôi phục dữ liệu. $g'(0),g'(\pi)$ khác không nên vẫn có vấn đề tương thích cấp hai.

\textbf{Chi tiết tích phân câu (i).} Đặt $k=n\pi$:
\[\begin{aligned}
B_n&=2\left[\frac{(x-1/2)\sin kx}{k}\right]_0^1-\frac2k\int_0^1\sin kx\,dx\\
&=-\frac2k\left[-\frac{\cos kx}{k}\right]_0^1
=\frac{2((-1)^n-1)}{k^2}.
\end{aligned}\]
Vì $\sin0=\sin k=0$, hạng biên đầu hết. $B_0=\int_0^1(x-1/2)dx=0$. Ở (ii), trung bình $B_0=\pi^{-1}\int_0^\pi\sin x\,dx=2/\pi$; các mode khác dùng $2\sin x\cos nx=\sin((1+n)x)+\sin((1-n)x)$. Tích phân cho $B_n=2(1+(-1)^n)/[\pi(1-n^2)]$ khi $n\ne1$, còn $B_1=\pi^{-1}\int_0^\pi\sin2x\,dx=0$.

%ID BT_3_3_3_05
\textbf{Bước 1. Trị riêng và trực giao.} $X'(0)=0$, $X(L)=0$. Tích phân từng phần cho $\lambda\int X^2=\int X'^2$. Nếu $\lambda=0$, $X$ hằng rồi $X(L)=0$ nên chỉ có nghiệm không; vì vậy $\lambda>0$. Với $\lambda=k^2$, điều kiện đầu cho $X=C\cos kx$, điều kiện cuối cho $\cos kL=0$. Suy $k_n=(n+1/2)\pi/L$. Trừ hai phương trình của $X_m,X_n$ và tích phân; các hạng biên bằng không, chứng minh trực giao.

\textbf{Bước 2. Chiếu dữ liệu và ghép.} $A_n=(2/L)\int f\cos k_nx$, $B_n=(2/L)\int g\cos k_nx$; $u=\sum(A_n\cos ak_nt+B_n\sin ak_nt/(ak_n))\cos k_nx$. Mọi mode có chu kỳ $4L/a$ và đổi dấu sau $2L/a$.

\textbf{Bước 3. Trường hợp cụ thể.} $a=1$, $L=\pi$, $f=\pi-x$, $g=\cos^3(x/2)$. Tích phân từng phần cho $\int_0^\pi(\pi-x)\cos k_nx=1/k_n^2$, nên $A_n=2/(\pi k_n^2)$. Công thức cos lập phương cho $B_0=3/4$, $B_1=1/4$. Vì vậy
\[u=\sum_{n\ge0}\frac2{\pi k_n^2}\cos k_nt\cos k_nx
+\tfrac32\sin(t/2)\cos(x/2)+\tfrac16\sin(3t/2)\cos(3x/2).\]
\textbf{Bước 4. Kết luận.} Mode đầu khác không nên chu kỳ cơ bản $4\pi$. $f'(0)=-1$ không tương thích Neumann, nên nghiệm không $C^1$ qua mặt phản xạ từ góc trái. Hai chuỗi vẫn đúng dữ liệu trong miền và từng mode thỏa PDE/biên; tính duy nhất trong lớp năng lượng theo năng lượng của hiệu.

\textbf{Chi tiết điều kiện biên và tích phân dữ liệu.} $X=A\cos kx+B\sin kx$, $X'(0)=Bk=0$ cho $B=0$ khi $k>0$; $X(L)=A\cos kL=0$ cho $\cos kL=0$. Chuẩn $\int_0^L\cos^2k_nx\,dx=L/2$ vì $\sin(2k_nL)=0$. Với $L=\pi$,
\[\begin{aligned}
\int_0^\pi(\pi-x)\cos k_nx\,dx
&=\left[\frac{(\pi-x)\sin k_nx}{k_n}\right]_0^\pi
+\frac1{k_n}\int_0^\pi\sin k_nx\,dx\\
&=\frac{1-\cos(k_n\pi)}{k_n^2}=1/k_n^2.
\end{aligned}\]
Chiếu $g=(3\cos(x/2)+\cos(3x/2))/4$ lên hệ cos trực giao cho $B_0=3/4,B_1=1/4$, rồi chia tần số $k_0=1/2,k_1=3/2$ khi ghép nghiệm.

\textbf{Chi tiết trực giao và dấu trị riêng.} Nhân $-X''=\lambda X$ bởi $X$, tích phân:
\[\lambda\int_0^LX^2dx=-[X'X]_0^L+\int_0^L(X')^2dx=\int_0^L(X')^2dx.\]
Hạng biên hết vì mỗi đầu có $X=0$ hoặc $X'=0$. Nếu $\lambda=0$, $\int(X')^2=0$ cho $X$ hằng; khi có một đầu Dirichlet, hằng đó bằng $0$. Với hai trị riêng,
\[(\lambda_m-\lambda_n)\int_0^LX_mX_n\,dx
=\int_0^L(-X_m''X_n+X_n''X_m)dx
=[-X_m'X_n+X_n'X_m]_0^L=0.\]
Nếu $\lambda_m\ne\lambda_n$, chia hiệu trị riêng cho tích vô hướng bằng $0$.

%ID VD_3_4_2_01
\textbf{Bước 1. Hệ sine và đáp ứng nguồn cos.} $a=L=1$, $X_n=\sin(n\pi x)$, $\omega_n=n\pi$. Khai triển $x$ có hệ số $c_n=2\int_0^1x\sin(n\pi x)dx=2(-1)^{n+1}/(n\pi)$. Đặt
\[C_\omega(\alpha,t)=\frac{\cos\alpha t-\cos\omega t}{\omega^2-\alpha^2},\qquad
C_\omega(\omega,t)=\frac{t\sin\omega t}{2\omega}.\]
Đây là nghiệm $C''+\omega^2C=\cos\alpha t$, $C(0)=C'(0)=0$; nhánh cộng hưởng lấy giới hạn hoặc tích phân Duhamel.

\textbf{Bước 2. Câu (i), biên không và nguồn $x\cos\alpha t$.} Dữ liệu đầu không, nên $y_n=c_nC_{\omega_n}$ và $u=\sum_{n\ge1}c_nC_{n\pi}\sin(n\pi x)$. Nếu $\alpha=m\pi$, chỉ mode $m$ dùng nhánh cộng hưởng, không chia $0$.

\textbf{Bước 3. Câu (ii), nguồn không và biên phải $\cos\alpha t$.} Lấy $b=x\cos\alpha t$, $w=u-b$. Khi đó $Q=\alpha^2x\cos\alpha t$, $w(x,0)=-x$, $w_t(x,0)=0$. Suy
\[u=x\cos\alpha t+\sum_{n\ge1}c_n[-\cos(n\pi t)+\alpha^2C_{n\pi}(\alpha,t)]\sin(n\pi x).\]
\textbf{Bước 4. Kiểm tra và góc.} Hàm nâng cho đúng biên, ODE cho nguồn của $w$ rồi triệt tiêu với $b_{tt}$. Tại $t=0$ khai triển $-x$ bù li độ của $b$ trong miền. Nhưng ở $(1,0)$, biên phải bằng $1$ trái li độ đầu $0$, nên câu (ii) không có nghiệm liên tục trên toàn bao đóng. Câu (i) có nguồn không khớp gia tốc biên tại góc phải nếu đòi $C^2$ đến góc; chuỗi vẫn giải trong miền với cách hiểu thích hợp.

\textbf{Chi tiết hệ số nguồn và đáp ứng.} Tích phân từng phần:
\[c_n=2\left[-\frac{x\cos(n\pi x)}{n\pi}+
\frac{\sin(n\pi x)}{n^2\pi^2}\right]_0^1
=-\frac{2(-1)^n}{n\pi}.\]
Với $\alpha^2\ne\omega^2$, thử $C_p=A\cos\alpha t$: $(\omega^2-\alpha^2)A=1$. Tổng $C=A\cos\alpha t+B\cos\omega t+D\sin\omega t$ có dữ liệu không khi $B=-A,D=0$. Với $\alpha=\omega$, thử $C=bt\sin\omega t$; đạo hàm hai lần cho $C''+\omega^2C=2b\omega\cos\omega t$, nên $b=1/(2\omega)$. Trong câu (ii), $b_{xx}=0$, $b_{tt}=-\alpha^2x\cos\alpha t$; vì $u_{tt}-u_{xx}=0$, phần $w$ có nguồn $-b_{tt}=\alpha^2x\cos\alpha t$ và dữ liệu $w_0=-x,w_{t,0}=0$.

%ID VD_3_4_2_02
\textbf{Bước 1. Tính duy nhất.} Với hiệu $v$ của hai nghiệm cùng dữ liệu và biên Neumann, $v_{tt}=a^2v_{xx}$, $v_x(0,t)=v_x(L,t)=0$. Đạo hàm $E=\frac12\int(v_t^2+a^2v_x^2)$ cho $E'=a^2[v_tv_x]_0^L=0$. $E(0)=0$ suy $v_t=v_x=0$; li độ đầu bằng $0$ xác định hằng số là $0$.

\textbf{Bước 2. Nâng biên tổng quát.} Với $u_x(0,t)=\mu(t)$, $u_x(L,t)=\nu(t)$, chọn $b=\mu x+(\nu-\mu)x^2/(2L)$, vì $b_x(0)=\mu$, $b_x(L)=\nu$. Đặt $w=u-b$; nguồn mới là $Q=F+a^2(\nu-\mu)/L-b_{tt}$, hai dữ liệu mới là $f-b_0$, $g-b_{t,0}$.

\textbf{Bước 3. Tách mode.} Với $X_0=1$, $X_n=\cos(n\pi x/L)$, lấy hệ số $A_0,B_0,Q_0$ là trung bình và các hệ số $n\ge1$ là $2/L$ lần tích phân với $X_n$. Mode hằng giải $y_0''=Q_0$, nên $y_0=A_0+B_0t+\int_0^t(t-s)Q_0(s)ds$. Mode còn lại giải
\[y_n=A_n\cos\omega_nt+\frac{B_n}{\omega_n}\sin\omega_nt+\frac1{\omega_n}\int_0^t\sin(\omega_n(t-s))Q_n(s)ds.\]
\textbf{Bước 4. Trở lại nghiệm.} $u=b+y_0+\sum_{n\ge1}y_nX_n$. Kiểm tra $t=0$ sau khi cộng hàm nâng; kiểm tra từng đạo hàm biên. Công thức mode hằng khác các mode dao động, nên không áp dụng mẫu $1/\omega_0$. Nếu dữ liệu không tương thích $f'(0)=\mu(0)$, $g'(0)=\mu'(0)$ (và tương tự đầu phải), nêu rõ mức trơn và xét nghiệm yếu thay vì khẳng định cổ điển đến góc.

\textbf{Chi tiết tích phân từng phần trong đạo hàm năng lượng.} Đặt $E=\frac12\int_0^L(u_t^2+a^2u_x^2+\kappa u^2)dx$. Với $u_{tt}-a^2u_{xx}+\gamma u_t+\kappa u=F$,
\[\begin{aligned}
E'&=\int_0^L(u_tu_{tt}+a^2u_xu_{xt}+\kappa uu_t)dx\\
&=\int_0^Lu_t(a^2u_{xx}-\gamma u_t-\kappa u+F)dx
+a^2\int_0^Lu_xu_{xt}dx+\kappa\int_0^Luu_tdx\\
&=a^2[u_tu_x]_0^L-a^2\int_0^Lu_{tx}u_xdx
+a^2\int_0^Lu_xu_{xt}dx-\gamma\int_0^Lu_t^2dx+\int_0^LFu_tdx\\
&=a^2[u_tu_x]_0^L-\gamma\int_0^Lu_t^2dx+\int_0^LFu_tdx.
\end{aligned}\]
Hai hạng $\kappa uu_t$ triệt tiêu ngay khi thế PDE; hai tích phân $u_xu_{xt}$ triệt tiêu sau tích phân từng phần. Với biên Dirichlet thuần nhất, đạo hàm vết theo $t$ cho $u_t=0$; với biên Neumann thuần nhất, $u_x=0$. Nếu định nghĩa năng lượng không có $1/2$, nhân cả đẳng thức đạo hàm với $2$.

\textbf{Chi tiết đạo hàm đáp ứng mode nguồn.} Đặt $D_n(t)=\omega_n^{-1}\int_0^t\sin(\omega_n(t-s))Q_n(s)ds$. Vì $\sin0=0$,
\[D_n'=\int_0^t\cos(\omega_n(t-s))Q_n(s)ds,\quad
D_n''=Q_n(t)-\omega_n\int_0^t\sin(\omega_n(t-s))Q_n(s)ds=Q_n-\omega_n^2D_n.\]
Tại $0$, $D_n=D_n'=0$. Vì thế tích phân vừa cho đúng nguồn, vừa giữ nguyên cả hai dữ liệu ban đầu của các hạng thuần nhất. Với mode hằng, $D_0=\int_0^t(t-s)Q_0(s)ds$, $D_0'=\int_0^tQ_0(s)ds$, $D_0''=Q_0(t)$, không chia cho $\omega_0=0$.

%ID BT_3_4_3_01
\textbf{Bước 1. Tính hệ số và nhân đáp ứng.} Trên $(0,\pi)$, khai triển $1$ theo sine có $c_n=2(1-(-1)^n)/(\pi n)$. Đặt
\[R_\omega(\alpha,t)=\frac{\sin\alpha t-(\alpha/\omega)\sin\omega t}{\omega^2-\alpha^2}.\]
Nó thỏa $R''+\omega^2R=\sin\alpha t$, $R(0)=R'(0)=0$. Khi $\alpha=\omega$, dùng $R=\sin\omega t/(2\omega^2)-t\cos\omega t/(2\omega)$.

\textbf{Bước 2. Câu (i).} $a=1$, biên và dữ liệu đầu không, nguồn $\sin\alpha t$. Chiếu nguồn cho $y_n''+n^2y_n=c_n\sin\alpha t$, nên $u=\sum c_nR_n\sin nx$. Tại cộng hưởng $\alpha=m$, thay đúng mode $m$ bằng nhánh giới hạn; nếu $c_m=0$, mode ấy không được kích thích.

\textbf{Bước 3. Câu (ii).} $a=2$, nguồn không, cả hai biên bằng $\sin\alpha t$. Lấy $b=\sin\alpha t$, $w=u-b$, suy $Q=\alpha^2\sin\alpha t$, $w_0=0$, $w_{t,0}=-\alpha$. Hệ số vận tốc là $-\alpha c_n$, vì vậy
\[u=\sin\alpha t+\sum_{n\ge1}c_n\left[-\frac\alpha{2n}\sin2nt+\alpha^2R_{2n}(\alpha,t)\right]\sin nx.\]
Đạo hàm ban đầu của hàm nâng và chuỗi triệt tiêu trong miền. Tại góc, vận tốc biên $\alpha$ trái $g=0$, nên không có nghiệm $C^1$ trên toàn bao đóng khi $\alpha\ne0$.

\textbf{Bước 4. Câu (iii).} $a=1$, nguồn $1$, biên trái $0$, biên phải $-\pi^2/2$, dữ liệu đầu không. Chọn $P=-x^2/2$: $P_{tt}-P_{xx}=1$ và đúng biên. $w=u-P$ thuần nhất, li độ $x^2/2$, vận tốc $0$. Tích phân từng phần hai lần cho
\[A_n=\frac1\pi\int_0^\pi x^2\sin nx\,dx
=-\frac{\pi(-1)^n}{n}+\frac{2((-1)^n-1)}{\pi n^3}.\]
Suy $u=-x^2/2+\sum A_n\cos nt\sin nx$. Li độ đầu và biên phải mâu thuẫn ở góc; công thức là nghiệm từng miền/ yếu. Mỗi câu đã kiểm tra nguồn qua ODE và dữ liệu qua hệ số, không chỉ ghi chuỗi hình thức.

\textbf{Chi tiết đáp ứng sine và tích phân đa thức.} Không cộng hưởng, thử $R_p=A\sin\alpha t$, có $A=1/(\omega^2-\alpha^2)$. Viết $R=R_p+B\cos\omega t+D\sin\omega t$; dữ liệu $R(0)=0$ cho $B=0$, $R'(0)=0$ cho $D=-\alpha A/\omega$. Ở cộng hưởng, $-t\cos\omega t/(2\omega)$ có nguồn $\sin\omega t$ nhưng vận tốc đầu $-1/(2\omega)$; cộng $\sin\omega t/(2\omega^2)$ để bù.

Trong (iii), đặt $I=\int_0^\pi x^2\sin nx\,dx$. Hai lần tích phân:
\[\begin{aligned}
I&=\left[-x^2\cos nx/n\right]_0^\pi+\frac2n\int_0^\pi x\cos nx\,dx,\\
\int_0^\pi x\cos nx\,dx
&=\left[x\sin nx/n+\cos nx/n^2\right]_0^\pi
=\frac{(-1)^n-1}{n^2}.
\end{aligned}\]
Vậy $I=-\pi^2(-1)^n/n+2((-1)^n-1)/n^3$, rồi chia $\pi$ cho $A_n$.

%ID BT_3_4_3_02
\textbf{Bước 1. Giữ mode hằng Neumann.} Hệ hàm riêng trên $(0,\pi)$ là $1,\cos nx$. Mode hằng dùng $y_0''=Q_0$; các mode $n\ge1$ dùng ODE dao động và $R_\omega$ như Bài 1.

\textbf{Bước 2. Câu (i).} Biên Neumann không, $a=1$, nguồn $\sin\alpha t$ không phụ thuộc $x$. Chỉ mode $0$ có nguồn khác không. Với dữ liệu không,
\[u=\int_0^t(t-s)\sin\alpha s\,ds=t/\alpha-\sin\alpha t/\alpha^2.\]
Thế lại cho $u_{tt}=\sin\alpha t$, $u_{xx}=0$, dữ liệu không. Nếu $\alpha=0$, lấy giới hạn $u=0$.

\textbf{Bước 3. Câu (ii).} $a=2$, $\mu=\nu=\sin\alpha t$, nguồn không. Lấy $b=x\sin\alpha t$, $w=u-b$, suy $Q=\alpha^2x\sin\alpha t$, $w_0=0$, $w_{t,0}=-\alpha x$. Khai triển $x$ theo cos có $d_0=\pi/2$, $d_n=2((-1)^n-1)/(\pi n^2)$, do tích phân từng phần.

\textbf{Bước 4. Giải từng ODE và ghép.} Mode hằng là $y_0=-d_0\sin\alpha t$ (đạo hàm đầu $-\alpha d_0$ và $y_0''=\alpha^2d_0\sin\alpha t$). Các mode khác là $y_n=d_n[-\alpha\sin2nt/(2n)+\alpha^2R_{2n}]$. Do đó
\[u=(x-\pi/2)\sin\alpha t+\sum_{n\ge1}d_n[-\alpha\sin2nt/(2n)+\alpha^2R_{2n}(\alpha,t)]\cos nx.\]
Hàm nâng cho đúng đạo hàm biên; các ODE cho đúng PDE. Tại góc, $g'=0$ trái $\mu'(0)=\nu'(0)=\alpha$, nên không $C^2$ đến góc nếu $\alpha\ne0$. Tại $\alpha=2n$ phải dùng nhánh cộng hưởng, không chia mẫu bằng không.

\textbf{Chi tiết tích phân nguồn không gian hằng và hệ số cos.} Với $\alpha\ne0$,
\[\int_0^t(t-s)\sin\alpha s\,ds
=\left[-(t-s)\cos\alpha s/\alpha\right]_0^t-\frac1\alpha\int_0^t\cos\alpha s\,ds
=t/\alpha-\sin\alpha t/\alpha^2.\]
Với dữ liệu $x$, $d_0=\pi^{-1}[x^2/2]_0^\pi=\pi/2$ và
\[d_n=\frac2\pi\left[x\sin nx/n+\cos nx/n^2\right]_0^\pi
=\frac{2((-1)^n-1)}{\pi n^2}.\]
Mode hằng giải $y_0''=\alpha^2d_0\sin\alpha t$ với $y_0'(0)=-\alpha d_0$: tích phân cho $y_0'=-\alpha d_0\cos\alpha t$, $y_0=-d_0\sin\alpha t$. Đây là phép bù tạo ra $(x-\pi/2)\sin\alpha t$.

%ID BT_3_4_3_03
\textbf{Bước 1. Chứng minh duy nhất.} Hiệu $v$ có nguồn và dữ liệu không, $v(0,t)=v_x(L,t)=0$. Đạo hàm năng lượng $E'=a^2[v_tv_x]_0^L$: ở đầu trái $v_t=0$ vì vết $v=0$, đầu phải $v_x=0$. Suy $E=E(0)=0$, $v$ hằng cả không gian và thời gian, rồi dữ liệu đầu buộc $v=0$.

\textbf{Bước 2. Nâng biên.} Lấy $b=\mu(t)+x\nu(t)$ để $b(0,t)=\mu$, $b_x(L,t)=\nu$. $w=u-b$ có biên hỗn hợp không, nguồn $Q=F-\mu''-x\nu''$, li độ $f-\mu(0)-x\nu(0)$, vận tốc $g-\mu'(0)-x\nu'(0)$.

\textbf{Bước 3. Chiếu nguồn cùng dữ liệu.} $k_n=(n+1/2)\pi/L$, $X_n=\sin k_nx$, chuẩn $L/2$. Tính $A_n=(2/L)\int w_0X_n$, $B_n=(2/L)\int w_{t,0}X_n$, $Q_n(s)=(2/L)\int Q(x,s)X_n(x)dx$.

\textbf{Bước 4. Giải và trả về biến gốc.}
\[u=b+\sum_{n\ge0}\left[A_n\cos\omega_nt+\frac{B_n}{\omega_n}\sin\omega_nt+\int_0^t\frac{\sin(\omega_n(t-s))}{\omega_n}Q_n(s)ds\right]X_n,\quad\omega_n=ak_n.\]
Với $\mu=\nu=0$, $b=0$ và dữ liệu/nguồn chiếu chính là dữ liệu gốc. Đạo hàm hai lần của tích phân mode cho $Q_n$, còn hai dữ liệu ban đầu của nó bằng không. Cộng hàm nâng xác nhận mọi biên. Cần điều kiện tương thích ở góc nếu muốn nghiệm cổ điển, chẳng hạn $f(0)=\mu(0)$, $g(0)=\mu'(0)$, $f'(L)=\nu(0)$, $g'(L)=\nu'(0)$.

\textbf{Chi tiết tích phân từng phần trong đạo hàm năng lượng.} Đặt $E=\frac12\int_0^L(u_t^2+a^2u_x^2+\kappa u^2)dx$. Với $u_{tt}-a^2u_{xx}+\gamma u_t+\kappa u=F$,
\[\begin{aligned}
E'&=\int_0^L(u_tu_{tt}+a^2u_xu_{xt}+\kappa uu_t)dx\\
&=\int_0^Lu_t(a^2u_{xx}-\gamma u_t-\kappa u+F)dx
+a^2\int_0^Lu_xu_{xt}dx+\kappa\int_0^Luu_tdx\\
&=a^2[u_tu_x]_0^L-a^2\int_0^Lu_{tx}u_xdx
+a^2\int_0^Lu_xu_{xt}dx-\gamma\int_0^Lu_t^2dx+\int_0^LFu_tdx\\
&=a^2[u_tu_x]_0^L-\gamma\int_0^Lu_t^2dx+\int_0^LFu_tdx.
\end{aligned}\]
Hai hạng $\kappa uu_t$ triệt tiêu ngay khi thế PDE; hai tích phân $u_xu_{xt}$ triệt tiêu sau tích phân từng phần. Với biên Dirichlet thuần nhất, đạo hàm vết theo $t$ cho $u_t=0$; với biên Neumann thuần nhất, $u_x=0$. Nếu định nghĩa năng lượng không có $1/2$, nhân cả đẳng thức đạo hàm với $2$.

\textbf{Chi tiết đạo hàm đáp ứng mode nguồn.} Đặt $D_n(t)=\omega_n^{-1}\int_0^t\sin(\omega_n(t-s))Q_n(s)ds$. Vì $\sin0=0$,
\[D_n'=\int_0^t\cos(\omega_n(t-s))Q_n(s)ds,\quad
D_n''=Q_n(t)-\omega_n\int_0^t\sin(\omega_n(t-s))Q_n(s)ds=Q_n-\omega_n^2D_n.\]
Tại $0$, $D_n=D_n'=0$. Vì thế tích phân vừa cho đúng nguồn, vừa giữ nguyên cả hai dữ liệu ban đầu của các hạng thuần nhất. Với mode hằng, $D_0=\int_0^t(t-s)Q_0(s)ds$, $D_0'=\int_0^tQ_0(s)ds$, $D_0''=Q_0(t)$, không chia cho $\omega_0=0$.

%ID BT_3_4_3_04
\textbf{Bước 1. Nhận dạng một mode.} Biên cố định trên $(0,1)$ và nguồn $\sin\pi x\sin t$ gợi $u=\sin\pi x\,T(t)$. Đạo hàm $u_{xx}=-\pi^2\sin\pi x\,T$; vì tốc độ bằng $1$, PDE trở thành $T''+\pi^2T=\sin t$, dữ liệu $T(0)=T'(0)=0$.

\textbf{Bước 2. Nghiệm riêng và thuần nhất.} Nguồn tần số $1$ khác $\pi$, nên thử $T_p=C\sin t$, được $C=1/(\pi^2-1)$. Nghiệm tổng $T=C\sin t+A\cos\pi t+B\sin\pi t$. Dữ liệu đầu cho $A=0$, $C+\pi B=0$, nên $B=-1/(\pi(\pi^2-1))$.

\textbf{Bước 3. Nghiệm.}
\[u(x,t)=\frac{\sin t-\sin\pi t/\pi}{\pi^2-1}\sin\pi x.\]
\textbf{Bước 4. Kiểm tra.} Hai biên bằng không vì $\sin\pi x=0$; $T(0)=0$ và $T'(0)=(1-1)/(\pi^2-1)=0$. Thế $T''+\pi^2T$ cho $\sin t$, nên đúng nguồn. Tính duy nhất theo năng lượng của hiệu hai nghiệm biên cố định.

%ID VD_3_5_2_01
\textbf{Bước 1. Năng lượng của sóng có cản.} PDE là $u_{tt}+u_t=4u_{xx}$, biên cố định trên $(0,\pi)$. Với $E=\frac12\int(u_t^2+4u_x^2)$, tích phân từng phần cho $E'=4[u_tu_x]_0^\pi-\int u_t^2=-\int u_t^2\le0$. Do vậy năng lượng giảm; hiệu hai nghiệm cùng dữ liệu có năng lượng đầu $0$, chứng minh duy nhất.

\textbf{Bước 2. Tách biến và bỏ cản bằng đổi ẩn.} Mode $\sin nx$ có $T_n''+T_n'+4n^2T_n=0$. Đặt $T_n=e^{-t/2}Y_n$, được $Y_n''+(4n^2-1/4)Y_n=0$, $Y_n(0)=A_n$, $Y_n'(0)=B_n+A_n/2$. Đặt $\omega_n=\sqrt{4n^2-1/4}$:
\[T_n=e^{-t/2}[A_n\cos\omega_nt+(B_n+A_n/2)\sin\omega_nt/\omega_n].\]
Hạng $A_n/2$ cần để bù đạo hàm của $e^{-t/2}$.

\textbf{Bước 3. Dữ liệu cụ thể.} $f=\sin^3x$ cho $A_1=3/4$, $A_3=-1/4$, các hệ số khác $0$. $g=1$ cho $B_n=2(1-(-1)^n)/(\pi n)$. Nghiệm $u=\sum_{n\ge1}T_n\sin nx$ với các hệ số này.

\textbf{Bước 4. Kiểm tra.} Tại $t=0$, $T_n=A_n$, $T_n'=B_n$, đúng cả hai dữ liệu. Biên đúng theo sine và ODE chứng minh PDE. Tuy nhiên $g(0)=g(\pi)=1$ trái biên cố định, nên công thức là nghiệm năng lượng, có mặt gãy từ góc; không có nghiệm $C^1$ trên toàn bao đóng. Đẳng thức năng lượng vẫn đúng trong lớp năng lượng bằng xấp xỉ dữ liệu.

\textbf{Chi tiết tích phân của dữ liệu hằng.} Với $k_n=n\pi/L$,
\[\int_0^L\sin k_nx\,dx=\left[-\frac{\cos k_nx}{k_n}\right]_0^L=\frac{1-(-1)^n}{k_n}.\]
Nhân $2/L$ cho hệ số của hằng $1$ là $2(1-(-1)^n)/(n\pi)$; hằng $-1$ đổi dấu. Khi chiếu vận tốc, còn chia $\omega_n$ trong hệ số $\sin\omega_nt$; khi chiếu li độ, hệ số $\cos\omega_nt$ giữ nguyên.

\textbf{Chi tiết tích phân từng phần trong đạo hàm năng lượng.} Đặt $E=\frac12\int_0^L(u_t^2+a^2u_x^2+\kappa u^2)dx$. Với $u_{tt}-a^2u_{xx}+\gamma u_t+\kappa u=F$,
\[\begin{aligned}
E'&=\int_0^L(u_tu_{tt}+a^2u_xu_{xt}+\kappa uu_t)dx\\
&=\int_0^Lu_t(a^2u_{xx}-\gamma u_t-\kappa u+F)dx
+a^2\int_0^Lu_xu_{xt}dx+\kappa\int_0^Luu_tdx\\
&=a^2[u_tu_x]_0^L-a^2\int_0^Lu_{tx}u_xdx
+a^2\int_0^Lu_xu_{xt}dx-\gamma\int_0^Lu_t^2dx+\int_0^LFu_tdx\\
&=a^2[u_tu_x]_0^L-\gamma\int_0^Lu_t^2dx+\int_0^LFu_tdx.
\end{aligned}\]
Hai hạng $\kappa uu_t$ triệt tiêu ngay khi thế PDE; hai tích phân $u_xu_{xt}$ triệt tiêu sau tích phân từng phần. Với biên Dirichlet thuần nhất, đạo hàm vết theo $t$ cho $u_t=0$; với biên Neumann thuần nhất, $u_x=0$. Nếu định nghĩa năng lượng không có $1/2$, nhân cả đẳng thức đạo hàm với $2$.

\textbf{Chi tiết đổi ẩn trong ODE có cản.} Với $T=e^{-t/2}Y$,
\[T'=e^{-t/2}(Y'-Y/2),\quad T''=e^{-t/2}(Y''-Y'+Y/4).\]
Do đó $T''+T'+\Omega^2T=e^{-t/2}[Y''+(\Omega^2-1/4)Y]$. Dữ liệu $T(0)=A,T'(0)=B$ trở thành $Y(0)=A,Y'(0)=B+A/2$. Nếu $\Omega=0$, giải trực tiếp $T''+T'=0$: đặt $Z=T'$, $Z=Be^{-t}$, rồi $T=A+B(1-e^{-t})$.

%ID VD_3_5_2_02
\textbf{Bước 1. Năng lượng đúng của dầm.} $u_{tt}=-a^2u_{xxxx}$, hai đầu ngàm $u=u_x=0$. Nhân với $u_t$ và tích phân hai lần:
\[E=\tfrac12\int_0^L(u_t^2+a^2u_{xx}^2)dx,\qquad E'=a^2[u_{tx}u_{xx}-u_tu_{xxx}]_0^L=0.\]
Do $u_t=u_{tx}=0$ ở hai đầu. Năng lượng có $u_x^2$ trong đề không phù hợp toán tử cấp bốn; cần sửa thành $u_{xx}^2$. Với hiệu hai nghiệm, $E=0$ suy $u_t=u_{xx}=0$; hàm affine còn lại bị biên ngàm và dữ liệu đầu buộc bằng $0$, nên duy nhất.

\textbf{Bước 2. Trị riêng dương và trực giao.} $X''''=\lambda X$, $X=X'=0$ ở hai đầu. Tích phân hai lần cho $\lambda\int X^2=\int(X'')^2$. Nếu $\lambda=0$, $X$ affine rồi biên loại hết; vậy $\lambda>0$. Trừ tích phân hai phương trình cho $(\lambda_m-\lambda_n)\int X_mX_n=0$ vì hạng biên $[X_m'''X_n-X_m''X_n'-X_n'''X_m+X_n''X_m']$ bằng $0$.

\textbf{Bước 3. Giải đủ bốn điều kiện.} Đặt $\lambda=k^4$, $\beta=kL$. Hai điều kiện trái cho $X=A(\cosh kx-\cos kx)+B(\sinh kx-\sin kx)$. Hai điều kiện phải tạo định thức $2(1-\cosh\beta\cos\beta)=0$, nên $\cosh\beta\cos\beta=1$. Chỉ lấy nghiệm dương; $\beta=0$ là nghiệm giả của cơ sở suy biến. Nghiệm đầu $\beta_1\approx4.730040745$. Chọn
\[X_n=\cosh k_nx-\cos k_nx-\sigma_n(\sinh k_nx-\sin k_nx),\quad
\sigma_n=\frac{\cosh\beta_n-\cos\beta_n}{\sinh\beta_n-\sin\beta_n}.\]
Phương trình đặc trưng bảo đảm cả $X_n(L)=X_n'(L)=0$.

\textbf{Bước 4. Chiếu dữ liệu và ghép.} $A_n=\int fX_n/\int X_n^2$, $B_n=\int gX_n/\int X_n^2$, tần số $\omega_n=ak_n^2$. Nghiệm $u=\sum(A_n\cos\omega_nt+B_n\sin\omega_nt/\omega_n)X_n$. Kiểm tra bốn biên bằng hàm riêng, hai dữ liệu bằng hệ số; không dùng tần số $ak_n$ của phương trình sóng cấp hai. Với dữ liệu đủ trơn và tương thích thu nghiệm cổ điển; với dữ liệu năng lượng hiểu chuỗi trong không gian thích hợp.

\textbf{Chi tiết hai tích phân năng lượng và định thức.} Nhân PDE với $u_t$:
\[\begin{aligned}
\int_0^Lu_tu_{xxxx}dx
&=[u_tu_{xxx}]_0^L-\int_0^Lu_{tx}u_{xxx}dx\\
&=[u_tu_{xxx}-u_{tx}u_{xx}]_0^L+\int_0^Lu_{txx}u_{xx}dx.
\end{aligned}\]
Hạng tích phân cuối bằng $\frac12(d/dt)\int u_{xx}^2$; chuyển sang cùng động năng cho $E'$ đã viết.

Đặt $C=\cosh\beta,c=\cos\beta,S=\sinh\beta,s=\sin\beta$. Hai điều kiện phải, sau chia đạo hàm cho $k$, là
\[\begin{pmatrix}C-c&S-s\\S+s&C-c\end{pmatrix}\binom AB=\binom00.\]
Định thức $(C-c)^2-(S-s)(S+s)=C^2-2Cc+c^2-S^2+s^2=2-2Cc$, vì $C^2-S^2=c^2+s^2=1$. Dòng đầu cho $B/A=-(C-c)/(S-s)$; dòng hai đúng khi định thức bằng $0$. Đây là nguồn gốc của $\sigma_n$.

%ID BT_3_5_3_01
\textbf{Bước 1. Sửa khẳng định bảo toàn và chứng minh duy nhất.} Phương trình (3.5.46) là $u_{tt}+u_t=a^2u_{xx}$. Với $I=\frac12\int(u_t^2+a^2u_x^2)$, $I'=a^2[u_tu_x]_0^L-\int u_t^2$. Ở biên Dirichlet, $u_t=0$; ở biên Neumann, $u_x=0$. Cả ba kiểu biên đều cho $I'=-\int u_t^2$, không bảo toàn nói chung. Hiệu hai nghiệm có $I(0)=0$, suy $u_t=u_x=0$, rồi dữ liệu đầu buộc hằng số bằng $0$.

\textbf{Bước 2. Ba bài trị riêng.} Tách $X''+\lambda X=0$. Tích phân cho $\lambda\int X^2=\int(X')^2$; hai kiểu hỗn hợp loại mode hằng, còn Neumann hai đầu giữ nó. Trừ hai phương trình chứng minh trực giao như mục 3.3. Cụ thể:

(a) $X(0)=X'(L)=0$: $X_n=\sin k_nx$, $k_n=(n+1/2)\pi/L$, $n\ge0$.

(b) $X'(0)=X(L)=0$: cùng $k_n$, $X_n=\cos k_nx$.

(c) $X'(0)=X'(L)=0$: $X_0=1$, $X_n=\cos(n\pi x/L)$, $n\ge1$. Chuẩn là $L$ cho mode $0$, $L/2$ cho mode khác không.

\textbf{Bước 3. Giải ODE và chiếu cả hai dữ liệu.} Đặt $A_n=\int fX_n/\int X_n^2$, $B_n=\int gX_n/\int X_n^2$. $T_n''+T_n'+a^2\lambda_nT_n=0$. Đổi $T_n=e^{-t/2}Y_n$ cho $Y_n''+(a^2\lambda_n-1/4)Y_n=0$. Nếu $\omega_n^2=a^2\lambda_n-1/4>0$,
\[T_n=e^{-t/2}[A_n\cos\omega_nt+(B_n+A_n/2)\sin\omega_nt/\omega_n].\]
Nếu $\omega_n=0$, dùng $e^{-t/2}[A_n+(B_n+A_n/2)t]$. Nếu $\omega_n^2<0$, thay cos/sin bằng cosh/sinh với $\sqrt{-\omega_n^2}$. Mode hằng là $A_0+B_0(1-e^{-t})$, giải trực tiếp $T_0''+T_0'=0$.

\textbf{Bước 4. Nghiệm và ví dụ cho cả ba kiểu.} $u=\sum T_nX_n$. Với $a=L=1$, lấy lần lượt $f=\sin(\pi x/2)$, $f=\cos(\pi x/2)$, $f=\cos\pi x$, tất cả $g=0$. Chỉ có một mode, $A=1$, $B=0$, nên nghiệm là $X(x)e^{-t/2}(\cos\omega t+\sin\omega t/(2\omega))$, với $\omega=\sqrt{\pi^2/4-1/4}$ cho (a),(b), $\sqrt{\pi^2-1/4}$ cho (c). Đồ thị thể hiện giảm biên độ; ví dụ mode hằng $f=1,g=0$ ở (c) cho nghiệm hằng $1$, phù hợp vì năng lượng đạo hàm bằng $0$.

\textbf{Chi tiết tích phân từng phần trong đạo hàm năng lượng.} Đặt $E=\frac12\int_0^L(u_t^2+a^2u_x^2+\kappa u^2)dx$. Với $u_{tt}-a^2u_{xx}+\gamma u_t+\kappa u=F$,
\[\begin{aligned}
E'&=\int_0^L(u_tu_{tt}+a^2u_xu_{xt}+\kappa uu_t)dx\\
&=\int_0^Lu_t(a^2u_{xx}-\gamma u_t-\kappa u+F)dx
+a^2\int_0^Lu_xu_{xt}dx+\kappa\int_0^Luu_tdx\\
&=a^2[u_tu_x]_0^L-a^2\int_0^Lu_{tx}u_xdx
+a^2\int_0^Lu_xu_{xt}dx-\gamma\int_0^Lu_t^2dx+\int_0^LFu_tdx\\
&=a^2[u_tu_x]_0^L-\gamma\int_0^Lu_t^2dx+\int_0^LFu_tdx.
\end{aligned}\]
Hai hạng $\kappa uu_t$ triệt tiêu ngay khi thế PDE; hai tích phân $u_xu_{xt}$ triệt tiêu sau tích phân từng phần. Với biên Dirichlet thuần nhất, đạo hàm vết theo $t$ cho $u_t=0$; với biên Neumann thuần nhất, $u_x=0$. Nếu định nghĩa năng lượng không có $1/2$, nhân cả đẳng thức đạo hàm với $2$.

\textbf{Chi tiết trực giao và dấu trị riêng.} Nhân $-X''=\lambda X$ bởi $X$, tích phân:
\[\lambda\int_0^LX^2dx=-[X'X]_0^L+\int_0^L(X')^2dx=\int_0^L(X')^2dx.\]
Hạng biên hết vì mỗi đầu có $X=0$ hoặc $X'=0$. Nếu $\lambda=0$, $\int(X')^2=0$ cho $X$ hằng; khi có một đầu Dirichlet, hằng đó bằng $0$. Với hai trị riêng,
\[(\lambda_m-\lambda_n)\int_0^LX_mX_n\,dx
=\int_0^L(-X_m''X_n+X_n''X_m)dx
=[-X_m'X_n+X_n'X_m]_0^L=0.\]
Nếu $\lambda_m\ne\lambda_n$, chia hiệu trị riêng cho tích vô hướng bằng $0$.

\textbf{Chi tiết đổi ẩn trong ODE có cản.} Với $T=e^{-t/2}Y$,
\[T'=e^{-t/2}(Y'-Y/2),\quad T''=e^{-t/2}(Y''-Y'+Y/4).\]
Do đó $T''+T'+\Omega^2T=e^{-t/2}[Y''+(\Omega^2-1/4)Y]$. Dữ liệu $T(0)=A,T'(0)=B$ trở thành $Y(0)=A,Y'(0)=B+A/2$. Nếu $\Omega=0$, giải trực tiếp $T''+T'=0$: đặt $Z=T'$, $Z=Be^{-t}$, rồi $T=A+B(1-e^{-t})$.

%ID BT_3_5_3_02
\textbf{Bước 1. Năng lượng và tính duy nhất.} Dầm $u_{tt}=-a^2u_{xxxx}$ dùng $E=\frac12\int(u_t^2+a^2u_{xx}^2)$, không phải năng lượng chứa $u_x^2$. Tích phân hai lần cho $E'=a^2[u_{tx}u_{xx}-u_tu_{xxx}]_0^L$. Mỗi cặp điều kiện trong đề làm cả hai tích biên bằng $0$: ngàm có $u_t=u_{tx}=0$; tựa có $u_t=u_{xx}=0$; đầu trượt có $u_{tx}=u_{xxx}=0$. Do đó $E$ bảo toàn. Hiệu hai nghiệm có $E=0$, $u_t=u_{xx}=0$, rồi dữ liệu đầu loại hàm affine còn lại, chứng minh duy nhất.

\textbf{Bước 2. Đọc đúng năm kiểu biên của bản scan.} (a) ngàm trái, tựa phải: $X(0)=X'(0)=X(L)=X''(L)=0$. (b) ngàm trái, trượt phải: $X(0)=X'(0)=X'(L)=X'''(L)=0$. Hai kiểu này khác hai đầu ngàm/ngàm–tự do của lời giải cũ.

\textbf{Bước 3. Giải phương trình đặc trưng.} Đặt $X''''=k^4X$, $\beta=kL$. Với (a),(b), hai điều kiện trái cho $X=A(\cosh kx-\cos kx)+B(\sinh kx-\sin kx)$.

(a) Từ $X(L)=X''(L)=0$, cộng/trừ suy $A\cosh\beta+B\sinh\beta=0$, $A\cos\beta+B\sin\beta=0$. Định thức cho $\cosh\beta\sin\beta-\sinh\beta\cos\beta=0$, hay $\tan\beta=\tanh\beta$. Chọn
\[X=\cosh kx-\cos kx-\coth\beta(\sinh kx-\sin kx).\]
Nghiệm dương đầu $\beta_1\approx3.926602312$.

(b) Từ $X'(L)=X'''(L)=0$, suy $A\sinh\beta+B\cosh\beta=0$, $A\sin\beta-B\cos\beta=0$. Định thức cho $\sinh\beta\cos\beta+\cosh\beta\sin\beta=0$, hay $\tan\beta=-\tanh\beta$. Chọn
\[X=\cosh kx-\cos kx-\tanh\beta(\sinh kx-\sin kx),\]
nghiệm dương đầu $\beta_1\approx2.365020372$.

(c) Tựa hai đầu $X=X''=0$: điều kiện trái loại cos/cosh, điều kiện phải loại sinh, nên $k_n=n\pi/L$, $X_n=\sin k_nx$, $n\ge1$.

(d) Tựa trái, trượt phải $X(0)=X''(0)=X'(L)=X'''(L)=0$: còn sine và sinh; cộng/trừ hai điều kiện phải loại sinh rồi cho $\cos kL=0$. $k_n=(n+1/2)\pi/L$, $X_n=\sin k_nx$, $n\ge0$.

(e) $X'=X'''=0$ hai đầu: $X_n=\cos(n\pi x/L)$, $n\ge1$, cộng mode $X_0=1$, $\lambda_0=0$.

\textbf{Bước 4. Dương, trực giao và nghiệm.} Tích phân hai lần cho $\lambda\int X^2=\int(X'')^2\ge0$ dưới cả năm kiểu biên. Nếu $\lambda=0$, $X$ affine; chỉ (e) giữ hằng. Các nghiệm $\beta=0$ trong (a),(b) phải loại bằng kiểm tra đa thức, không coi là mode. Trừ hai phương trình tích phân chứng minh trực giao với các trị riêng khác nhau. Lấy $A_n=\int fX_n/\int X_n^2$, $B_n=\int gX_n/\int X_n^2$:
\[u=\sum_n[A_n\cos(ak_n^2t)+B_n\sin(ak_n^2t)/(ak_n^2)]X_n.\]
Ở (e) thêm $\bar f+t\bar g$. Với $a=L=1$, mỗi kiểu lấy $f$ là hàm riêng dương đầu vừa tìm, $g=0$, cho ví dụ $u=X_1\cos(k_1^2t)$; đồ thị kèm đủ năm kiểu biên và không dùng nhầm các phương trình $\cos\beta\cosh\beta=\pm1$.

\textbf{Chi tiết giải đủ các điều kiện của (a),(b).} Viết $C=\cosh\beta,c=\cos\beta,S=\sinh\beta,s=\sin\beta$.
Ở (a), $X(L)=A(C-c)+B(S-s)=0$, $X''(L)/k^2=A(C+c)+B(S+s)=0$. Cộng hai dòng rồi chia $2$ cho $AC+BS=0$; lấy dòng sau trừ dòng trước rồi chia $2$ cho $Ac+Bs=0$. Định thức của hệ này là $Cs-Sc=0$, và $B/A=-C/S=-\coth\beta$.

Ở (b), $X'(L)/k=A(S+s)+B(C-c)=0$, $X'''(L)/k^3=A(S-s)+B(C+c)=0$. Cộng cho $AS+BC=0$; lấy dòng đầu trừ dòng sau cho $As-Bc=0$. Định thức là $-Sc-Cs=0$, và $B/A=-S/C=-\tanh\beta$. Với $\beta>0$, $\sinh\beta,\cosh\beta$ khác $0$; $\cos\beta=0$ không thỏa các định thức, nên các phép chia để viết tan hợp lệ tại nghiệm.

\textbf{Chi tiết ba kiểu còn lại.} Từ $X=A\cosh kx+B\sinh kx+C\cos kx+D\sin kx$, ở (c),(d), $X(0)=A+C=0$, $X''(0)/k^2=A-C=0$ cho $A=C=0$. (c) Hai dòng phải $B\sinh\beta+D\sin\beta=0$, $B\sinh\beta-D\sin\beta=0$ cho $B=0$, $\sin\beta=0$. (d) Hai dòng phải $B\cosh\beta+D\cos\beta=0$, $B\cosh\beta-D\cos\beta=0$ cho $B=0$, $\cos\beta=0$.

Ở (e), $X'(0)/k=B+D=0$, $X'''(0)/k^3=B-D=0$ cho $B=D=0$. Hai dòng phải $A\sinh\beta-C\sin\beta=0$, $A\sinh\beta+C\sin\beta=0$ cho $A=0$, $\sin\beta=0$. Riêng $\lambda=0$, $X''''=0$ cho $X=A+Bx+Cx^2+Dx^3$; $X'''=0$ cho $D=0$, $X'(0)=0$ cho $B=0$, $X'(L)=0$ cho $C=0$, còn hằng $A$.

\textbf{Chi tiết tích phân chứng minh trực giao dầm.} Tích phân từng phần hai lần:
\[\int_0^L(X_m''''X_n-X_n''''X_m)dx=
[X_m'''X_n-X_m''X_n'-X_n'''X_m+X_n''X_m']_0^L.\]
Hai tích phân còn lại là $\int X_m''X_n''$ và $\int X_n''X_m''$, triệt tiêu khi trừ. Tại đầu ngàm $X=X'=0$, đầu tựa $X=X''=0$, đầu trượt $X'=X'''=0$, mỗi tích biên mất một thừa số. Thay phương trình hàm riêng cho $(\lambda_m-\lambda_n)\int X_mX_n=0$, rồi chia hiệu trị riêng.

\textbf{Chi tiết hai lần tích phân từng phần của dầm.}
\[\int_0^L u_tu_{xxxx}dx
=[u_tu_{xxx}]_0^L-\int_0^Lu_{tx}u_{xxx}dx
=[u_tu_{xxx}-u_{tx}u_{xx}]_0^L+\int_0^Lu_{txx}u_{xx}dx.\]
Thế $u_{tt}=-a^2u_{xxxx}$ vào đạo hàm động năng. Tích phân cuối cùng triệt tiêu với đạo hàm $\frac{a^2}{2}\int u_{xx}^2$, còn $E'=a^2[u_{tx}u_{xx}-u_tu_{xxx}]_0^L$.

%ID BT_3_5_3_03
\textbf{Bước 1. Hệ sine và dữ liệu tam giác.} PDE $u_{tt}=u_{xx}-u$ trên $(0,\pi)$, biên không. $f=\min(x,\pi-x)$, $g=0$ có $A_n=4\sin(n\pi/2)/(\pi n^2)$, tính bằng tách tích phân tại $\pi/2$ như Bài 3.3.3.2.

\textbf{Bước 2. Giải đúng tần số mới.} $T_n''+(n^2+1)T_n=0$, $T_n(0)=A_n$, $T_n'(0)=0$. Do đó
\[u=\sum_{n\ge1}\frac{4\sin(n\pi/2)}{\pi n^2}\cos(\sqrt{n^2+1}\,t)\sin nx.\]
Hạng $+1$ dưới căn đến từ $-u$ trong PDE, không thể dùng tần số $n$.

\textbf{Bước 3. Không có chu kỳ chung.} Hệ số mode $1,3$ đều khác không. Nếu chu kỳ $T>0$ tồn tại, chiếu lên hai mode cho $\sqrt2T$ và $\sqrt{10}T$ đều là bội của $2\pi$, buộc tỷ số $\sqrt5$ hữu tỷ, mâu thuẫn. Vậy nghiệm này không tuần hoàn.

\textbf{Bước 4. Kiểm tra và mức trơn.} Mỗi mode thỏa PDE/biên và đúng dữ liệu. Tam giác có góc nên nghiệm thuộc lớp năng lượng, không cổ điển qua mặt truyền góc. Năng lượng bảo toàn phải chứa $u^2$: $\frac12\int(u_t^2+u_x^2+u^2)$.

\textbf{Chi tiết hệ số của tam giác.} Với $f(x)=\min(x,\pi-x)$, chia tại $b=\pi/2$. Đổi $y=\pi-x$ ở nửa sau cho
\[\int_0^\pi f(x)\sin nx\,dx=(1+(-1)^{n+1})\int_0^b x\sin nx\,dx.\]
Tích phân cuối bằng $[-x\cos nx/n+\sin nx/n^2]_0^b=-b\cos(nb)/n+\sin(nb)/n^2$. Nếu $n$ chẵn, hệ số trước tích phân bằng $0$; nếu $n$ lẻ, $\cos(nb)=0$. Nhân $2/\pi$ được $4\sin(n\pi/2)/(\pi n^2)$.

%ID BT_3_5_3_04
\textbf{Bước 1. Tách phương trình điện báo.} Với $u_{tt}+u_t+u=u_{xx}$, hai biên cố định trên $(0,1)$, đặt $u=\sum T_n\sin(n\pi x)$. ODE là $T_n''+T_n'+(n^2\pi^2+1)T_n=0$.

\textbf{Bước 2. Đổi ẩn và dữ liệu.} $A_n=2\int_0^1f(x)\sin(n\pi x)dx$, vận tốc đầu bằng $0$. Đặt $T_n=e^{-t/2}Y_n$ cho $Y_n''+(n^2\pi^2+3/4)Y_n=0$, $Y_n(0)=A_n$, $Y_n'(0)=A_n/2$.

\textbf{Bước 3. Nghiệm.}
\[u=e^{-t/2}\sum_{n\ge1}A_n\left[\cos\omega_nt+\frac{\sin\omega_nt}{2\omega_n}\right]\sin(n\pi x),\qquad\omega_n=\sqrt{n^2\pi^2+3/4}.\]
\textbf{Bước 4. Kiểm tra.} Tại $0$, đạo hàm của ngoặc là $A_n/2$, triệt tiêu đạo hàm $-A_n/2$ của hệ số mũ; vì thế vận tốc đầu đúng bằng $0$. Biên theo sine và PDE theo ODE. Năng lượng có thêm $u^2$, đạo hàm bằng $-\int u_t^2$, chứng minh duy nhất trong lớp thích hợp. Cần $f$ tương thích biên nếu đòi nghiệm trơn đến góc; công thức hệ số áp dụng cho dữ liệu năng lượng.

\textbf{Chi tiết tích phân từng phần trong đạo hàm năng lượng.} Đặt $E=\frac12\int_0^L(u_t^2+a^2u_x^2+\kappa u^2)dx$. Với $u_{tt}-a^2u_{xx}+\gamma u_t+\kappa u=F$,
\[\begin{aligned}
E'&=\int_0^L(u_tu_{tt}+a^2u_xu_{xt}+\kappa uu_t)dx\\
&=\int_0^Lu_t(a^2u_{xx}-\gamma u_t-\kappa u+F)dx
+a^2\int_0^Lu_xu_{xt}dx+\kappa\int_0^Luu_tdx\\
&=a^2[u_tu_x]_0^L-a^2\int_0^Lu_{tx}u_xdx
+a^2\int_0^Lu_xu_{xt}dx-\gamma\int_0^Lu_t^2dx+\int_0^LFu_tdx\\
&=a^2[u_tu_x]_0^L-\gamma\int_0^Lu_t^2dx+\int_0^LFu_tdx.
\end{aligned}\]
Hai hạng $\kappa uu_t$ triệt tiêu ngay khi thế PDE; hai tích phân $u_xu_{xt}$ triệt tiêu sau tích phân từng phần. Với biên Dirichlet thuần nhất, đạo hàm vết theo $t$ cho $u_t=0$; với biên Neumann thuần nhất, $u_x=0$. Nếu định nghĩa năng lượng không có $1/2$, nhân cả đẳng thức đạo hàm với $2$.

\textbf{Chi tiết đổi ẩn trong ODE có cản.} Với $T=e^{-t/2}Y$,
\[T'=e^{-t/2}(Y'-Y/2),\quad T''=e^{-t/2}(Y''-Y'+Y/4).\]
Do đó $T''+T'+\Omega^2T=e^{-t/2}[Y''+(\Omega^2-1/4)Y]$. Dữ liệu $T(0)=A,T'(0)=B$ trở thành $Y(0)=A,Y'(0)=B+A/2$. Nếu $\Omega=0$, giải trực tiếp $T''+T'=0$: đặt $Z=T'$, $Z=Be^{-t}$, rồi $T=A+B(1-e^{-t})$.

%ID BT_3_5_3_05
\textbf{Câu (i): tìm một hàm riêng của toán tử có đạo hàm cấp một.} Đặt $f=e^{-x}(\pi\cos\pi x+\sin\pi x)$. Đạo hàm trực tiếp cho $f'=-(\pi^2+1)e^{-x}\sin\pi x$, nên biên $f'(0)=f'(1)=0$. Đạo hàm thêm cho $f''+2f'=-(\pi^2+1)f$. Vì vậy $u=fT$ dẫn tới $T''+(\pi^2+1)T=0$, $T(0)=1,T'(0)=0$:
\[u=f(x)\cos(\sqrt{\pi^2+1}\,t).\]
Các tính toán trên kiểm tra đồng thời PDE, biên Neumann và dữ liệu đầu.

\textbf{Câu (ii), bước 1: đổi tọa độ và khử đạo hàm cấp một.} $z=\log(1+x)$, $\ell=\log2$. Vì $\partial_x=e^{-z}\partial_z$, $\partial_x((1+x)^2u_x)=u_{zz}+u_z$. Đặt $u=e^{-z/2}v$; tính $u_{zz}+u_z=e^{-z/2}(v_{zz}-v/4)$. Nguồn $t/\sqrt{1+x}=te^{-z/2}$, nên $v_{tt}=v_{zz}-v/4+t$, biên Dirichlet ở $z=0,\ell$.

\textbf{Câu (ii), bước 2: chiếu và giải nguồn tuyến tính.} $v_0=e^{z/2}(e^z-1)(2-e^z)$, $v_{t,0}=0$. Với $k_n=n\pi/\ell$, $\omega_n=\sqrt{k_n^2+1/4}$,
\[A_n=\frac2\ell\int_0^\ell e^{z/2}(e^z-1)(2-e^z)\sin k_nz\,dz,\quad c_n=\frac{2(1-(-1)^n)}{n\pi}.\]
ODE $T_n''+\omega_n^2T_n=c_nt$ cho nghiệm riêng dữ liệu không $c_n(t/\omega_n^2-\sin\omega_nt/\omega_n^3)$. Vì vậy
\[u=e^{-z/2}\sum_{n\ge1}\left[A_n\cos\omega_nt+c_n\left(\frac t{\omega_n^2}-\frac{\sin\omega_nt}{\omega_n^3}\right)\right]\sin k_nz.\]
Tích phân của $A_n$ là tích phân sơ cấp: khai triển $v_0=-e^{5z/2}+3e^{3z/2}-2e^{z/2}$ rồi dùng
\[\int_0^\ell e^{pz}\sin k_nz\,dz=\frac{k_n(1-(-1)^ne^{p\ell})}{p^2+k_n^2}.\]
Thay $p=5/2,3/2,1/2$ cùng các hệ số $-1,3,-2$ cho hệ số tường minh.

\textbf{Câu (iii), bước 1: chiếu nguồn và vận tốc.} PDE $u_{tt}+2u_t=u_{xx}+\sin x$, biên không trên $(0,1)$, $f=0,g=1$. Với $k_n=n\pi$, $\omega_n=\sqrt{k_n^2-1}$, vận tốc có $c_n=2(1-(-1)^n)/(n\pi)$. Tích–tổng cho nguồn
\[F_n=2\int_0^1\sin x\sin k_nx\,dx=\frac{2(-1)^{n+1}k_n\sin1}{k_n^2-1}.\]
\textbf{Câu (iii), bước 2: giải ODE có cản.} $T_n''+2T_n'+k_n^2T_n=F_n$, $T_n(0)=0$, $T_n'(0)=c_n$. Nghiệm riêng tĩnh $F_n/k_n^2$ được bù hai dữ liệu đầu bởi phần thuần nhất:
\[u=\sum_{n\ge1}\left[c_ne^{-t}\frac{\sin\omega_nt}{\omega_n}+\frac{F_n}{k_n^2}\left\{1-e^{-t}\left(\cos\omega_nt+\frac{\sin\omega_nt}{\omega_n}\right)\right\}\right]\sin k_nx.\]
Thế $t=0$ kiểm tra hai dữ liệu; từng ODE kiểm tra PDE. $g=1$ không tương thích biên cố định ở góc, nên câu (iii) không có nghiệm cổ điển toàn bao đóng; nguồn cũng có vết không phù hợp nếu đòi cấp hai đến biên.

\textbf{Chi tiết đạo hàm ở (i).} Với $f=e^{-x}(\pi\cos\pi x+\sin\pi x)$, quy tắc tích cho
\[f'=e^{-x}[-\pi\cos\pi x-\sin\pi x-\pi^2\sin\pi x+\pi\cos\pi x]
=-(\pi^2+1)e^{-x}\sin\pi x.\]
Đạo hàm tiếp: $f''=(\pi^2+1)e^{-x}(\sin\pi x-\pi\cos\pi x)$. Vì $2f'=-2(\pi^2+1)e^{-x}\sin\pi x$, cộng cho $f''+2f'=-(\pi^2+1)e^{-x}(\sin\pi x+\pi\cos\pi x)=-(\pi^2+1)f$.

\textbf{Chi tiết đổi ẩn và tính hệ số ở (ii).} $u=e^{-z/2}v$ cho $u_z=e^{-z/2}(v_z-v/2)$, $u_{zz}=e^{-z/2}(v_{zz}-v_z+v/4)$. Cộng $u_z$ được $e^{-z/2}(v_{zz}-v/4)$. Vì $(e^z-1)(2-e^z)=-e^{2z}+3e^z-2$, dữ liệu mới là $-e^{5z/2}+3e^{3z/2}-2e^{z/2}$. Nguyên hàm
\[\int e^{pz}\sin k_nz\,dz=
\frac{e^{pz}(p\sin k_nz-k_n\cos k_nz)}{p^2+k_n^2}\]
được kiểm tra bằng đạo hàm tích. Tại $0,\ell$, $\sin k_nz=0$, $\cos k_n\ell=(-1)^n$, nên
\[A_n=\frac{2k_n}{\ell}\left[
-\frac{1-(-1)^n2^{5/2}}{k_n^2+25/4}
+3\frac{1-(-1)^n2^{3/2}}{k_n^2+9/4}
-2\frac{1-(-1)^n2^{1/2}}{k_n^2+1/4}\right].\]
Với nguồn $c_nt$, nghiệm riêng $c_nt/\omega_n^2$ có vận tốc đầu $c_n/\omega_n^2$; trừ $c_n\sin\omega_nt/\omega_n^3$ để bù, cho đáp ứng dữ liệu không.

\textbf{Chi tiết nguồn và hằng số của (iii).} Tích–tổng cho
\[F_n=\left[\frac{\sin((1-k_n)x)}{1-k_n}
-\frac{\sin((1+k_n)x)}{1+k_n}\right]_0^1
=(-1)^n\sin1\left(\frac1{1-k_n}-\frac1{1+k_n}\right).\]
Rút gọn cho $F_n=2(-1)^{n+1}k_n\sin1/(k_n^2-1)$. Đặt $H_n=F_n/k_n^2$; tổng nghiệm ODE là $T_n=H_n+e^{-t}(A\cos\omega_nt+B\sin\omega_nt)$. $T_n(0)=0$ cho $A=-H_n$; $T_n'(0)=c_n$ cho $-A+\omega_nB=c_n$, nên $B=(c_n-H_n)/\omega_n$. Thế lại và gom hai nhóm cho công thức ở bước 2.

%ID BT_3_5_3_06
\textbf{Bước 1. Đạo hàm năng lượng Klein–Gordon.} Với $u_{tt}-u_{xx}+u=0$, đặt $E=\int_0^L(u^2+u_t^2+u_x^2)dx$. Tích phân từng phần:
\[E'=2\int_0^Lu_t(u_{tt}-u_{xx}+u)dx+2[u_tu_x]_0^L=0.\]
Ở biên cố định, $u_t=0$, nên thông lượng hết.

\textbf{Bước 2. Thế dữ liệu đầu và sửa câu (i).}
\[E(t)=E(0)=\int_0^L(f^2+g^2+(f')^2)dx.\]
Công thức in thiếu $g^2$. Nó chỉ đúng nếu thêm $g=0$. Lấy $f=0$, $g=\sin(\pi x/L)$: tại $t=0$ vế trái là $L/2$, còn vế phải in bằng $0$, nên yêu cầu nguyên văn sai.

\textbf{Bước 3. Câu (ii).} Khai triển và tích phân hạng chéo:
\[\int_0^L|f+f'|^2dx=\int_0^L(f^2+(f')^2)dx+[f^2]_0^L.\]
Nếu $f(0)=f(L)$ thì hạng biên hết; để bằng năng lượng còn cần $g=0$. Dưới biên cố định tương thích, hai giá trị $f$ đều $0$, nhưng điều đó không xóa năng lượng vận tốc. Công thức đúng tổng quát là $E=\int|f+f'|^2+\int g^2$ khi hai giá trị đầu bằng nhau. Đoạn không gian ghi $(0,1)$ cùng cận $L$ cần thống nhất; chứng minh đúng trên mọi $(0,L)$.

\textbf{Chi tiết tích phân từng phần trong đạo hàm năng lượng.} Đặt $E=\frac12\int_0^L(u_t^2+a^2u_x^2+\kappa u^2)dx$. Với $u_{tt}-a^2u_{xx}+\gamma u_t+\kappa u=F$,
\[\begin{aligned}
E'&=\int_0^L(u_tu_{tt}+a^2u_xu_{xt}+\kappa uu_t)dx\\
&=\int_0^Lu_t(a^2u_{xx}-\gamma u_t-\kappa u+F)dx
+a^2\int_0^Lu_xu_{xt}dx+\kappa\int_0^Luu_tdx\\
&=a^2[u_tu_x]_0^L-a^2\int_0^Lu_{tx}u_xdx
+a^2\int_0^Lu_xu_{xt}dx-\gamma\int_0^Lu_t^2dx+\int_0^LFu_tdx\\
&=a^2[u_tu_x]_0^L-\gamma\int_0^Lu_t^2dx+\int_0^LFu_tdx.
\end{aligned}\]
Hai hạng $\kappa uu_t$ triệt tiêu ngay khi thế PDE; hai tích phân $u_xu_{xt}$ triệt tiêu sau tích phân từng phần. Với biên Dirichlet thuần nhất, đạo hàm vết theo $t$ cho $u_t=0$; với biên Neumann thuần nhất, $u_x=0$. Nếu định nghĩa năng lượng không có $1/2$, nhân cả đẳng thức đạo hàm với $2$.

%ID BT_3_5_3_07
\textbf{Bước 1. Đạo hàm năng lượng.} Với $u_{tt}+\xi u_t=u_{xx}$, $\xi\ge0$, biên cố định, $E=\int_0^L(u_t^2+u_x^2)dx$. Nhân PDE với $2u_t$ và tích phân từng phần:
\[E'=2[u_tu_x]_0^L-2\xi\int_0^Lu_t^2dx=-2\xi\int_0^Lu_t^2dx.\]
Hạng biên hết vì đạo hàm thời gian của vết Dirichlet bằng $0$.

\textbf{Bước 2. Ước lượng trên.} Tích phân bình phương không âm, nên $E'\le0$ và $E(t)\le E(0)$.

\textbf{Bước 3. Ước lượng dưới.} $\int u_t^2\le E$, nên $E'\ge-2\xi E$. Do đó $(e^{2\xi t}E)'=e^{2\xi t}(E'+2\xi E)\ge0$. So với giá trị đầu:
\[e^{-2\xi t}E(0)\le E(t)\le E(0).\]
\textbf{Bước 4. Kiểm tra hệ số.} Vì định nghĩa $E$ không có $1/2$, hệ số đạo hàm là $2\xi$. Với $\xi=0$, hai cận bằng nhau và thu bảo toàn năng lượng. Với $E(0)=0$ suy nghiệm duy nhất bằng $0$ khi li độ đầu cũng bằng $0$. Các công thức đúng trên đoạn $(0,L)$ sau khi thống nhất ký hiệu cận phải của đề.

%ID VD_3_6_2_01
\textbf{Bước 1. Đưa Kirchhoff về diện tích chỏm.} Tốc độ $c=2$, li độ đầu không, vận tốc là chỉ thị $z>0$. Tại $P=(x,y,z)$, điểm trên mặt cầu là $P+2t\omega$ và thuộc miền vận tốc khi $z+2t\omega_z>0$. Đặt $\mu=\omega_z$; diện tích hướng có $\mu$ trong một khoảng bằng $2\pi$ nhân độ dài khoảng ấy.

\textbf{Bước 2. Cắt ngưỡng vào $[-1,1]$.} Điều kiện là $\mu>-z/(2t)$. Nếu $z\ge2t$, cả mặt cầu nằm phía trên; nếu $z\le-2t$, không có phần nào; nếu $|z|<2t$, độ dài khoảng là $1+z/(2t)$. Vì $u=t|E|/2$,
\[u(P,t)=\begin{cases}t&z\ge2t,\\t/2+z/4&|z|<2t,\\0&z\le-2t.\end{cases}\]
\textbf{Bước 3. Kiểm tra.} Ngoài các mặt $z=\pm2t$, mỗi nhánh affine nên $u_{tt}=4\Delta u=0$. $u(P,0)=0$; với $z\ne0$, $u_t(P,0)$ là $1$ hoặc $0$ đúng dữ liệu. Nghiệm liên tục nhưng gãy cấp một trên các mặt sóng. Tại $z=0$ dữ liệu chỉ thị gián đoạn nên không có vết cổ điển thống nhất; công thức cho $u=t/2$, là cách hiểu từ trung bình mặt cầu.

%ID VD_3_6_2_02
\textbf{Bước 1. Chọn trục theo điểm quan sát.} $c=2$, $\psi=\chi_{B_R(0)}$, $r=|P|$, $\rho=2t$. Nếu $r>0$, lấy trục của $\mu$ hướng về tâm từ $P$. Khi đó $|P+\rho\omega|^2=r^2+\rho^2-2r\rho\mu$.

\textbf{Bước 2. Bất đẳng thức quả cầu.} Điều kiện nằm trong $B_R$ tương đương $\mu>(r^2+\rho^2-R^2)/(2r\rho)$. Tỷ lệ diện tích chỏm là ngưỡng đã cắt:
\[K_R(r,t)=t\left[\frac{R^2-(r-2t)^2}{8rt}\right]_0^1,\qquad [s]_0^1=\min(1,\max(0,s)).\]
\textbf{Bước 3. Giải thích các nhánh.} Nếu $r<R$ và $2t<R-r$, cả mặt cầu ở trong, $u=t$. Nếu $|R-r|<2t<R+r$, chỉ một chỏm giao, $u=[R^2-(r-2t)^2]/(8r)$. Nếu không giao, $u=0$. Với $r=0$, tính riêng: $u=t$ khi $2t<R$, $0$ khi $2t>R$; giá trị tại $2t=R$ phụ thuộc quy ước chỉ thị trên toàn mặt cầu, với quả cầu mở là $0$.

\textbf{Bước 4. Kiểm tra.} Khi điểm ban đầu nằm trong quả cầu, nhánh đầu cho vận tốc $1$; ở ngoài cho $0$. Thời gian đến/đi là khoảng cách tới mặt cầu chia tốc độ. Dữ liệu gián đoạn cho nghiệm không cổ điển trên các mặt sóng; trường hợp tâm quan sát có thể có bước nhảy khi toàn mặt cầu đồng thời đi qua biên.

%ID VD_3_6_2_03
\textbf{Bước 1. Dùng phương pháp hạ chiều.} Dữ liệu hai chiều $u_0=0$, $u_{t,0}=f(x)$ không phụ thuộc $y$. Mở rộng sang ba chiều cũng không phụ thuộc $y,z$. Kirchhoff cho $u=t(4\pi)^{-1}\int_{S^2}f(x+ct\omega_x)d\omega$.

\textbf{Bước 2. Tích phân góc.} Chọn trục cực là trục $x$, $\mu=\omega_x$, $d\omega=d\phi\,d\mu$. Tích phân theo $\phi$ cho $2\pi$, nên $u=(t/2)\int_{-1}^1f(x+ct\mu)d\mu$. Đổi $s=x+ct\mu$, $ds=ct\,d\mu$:
\[u(x,y,t)=\frac1{2c}\int_{x-ct}^{x+ct}f(s)ds.\]
\textbf{Bước 3. Kiểm tra bằng d'Alembert.} Nghiệm không phụ thuộc $y$, nên $\Delta u=u_{xx}$. Công thức vừa tìm đúng d'Alembert một chiều với li độ không và vận tốc $f$. Tại $0$ tích phân có độ dài $0$, còn đạo hàm thời gian cho $f(x)$. Nếu $f$ không trơn, công thức vẫn có nghĩa và mức trơn được xét tại các cận chạm điểm gãy của $f$.

%ID VD_3_6_2_04
\textbf{Bước 1. Tách hai nguồn vận tốc.} $c=2$, $u_0=0$, $\psi=\chi_{[0,1]}(x)-\chi_{[0,1]}(y)$. Tuyến tính cho hiệu của một nghiệm chỉ phụ thuộc $x$ và một nghiệm chỉ phụ thuộc $y$.

\textbf{Bước 2. Tính độ dài giao.} Đặt $h(s)=[s]_0^1$, nguyên hàm của $\chi_{[0,1]}$. Theo Ví dụ 3,
\[u=\frac14\{h(x+2t)-h(x-2t)-h(y+2t)+h(y-2t)\}.\]
Mỗi hiệu $h$ bằng độ dài giao $[q-2t,q+2t]\cap[0,1]$ với $q=x$ hoặc $y$, nên có thể đọc mọi nhánh từ min/max.

\textbf{Bước 3. Kiểm tra và đối xứng.} Tại $t=0$, $u=0$; đạo hàm thời gian cho $\chi_{[0,1]}(x)-\chi_{[0,1]}(y)$. Hai phần thỏa PDE một chiều, nên tổng thỏa PDE hai chiều. Đổi $x,y$ đổi dấu, do đó $u(x,x,t)=0$. Khi cả hai khoảng ảnh hưởng đã bao trọn $[0,1]$, hai phần cùng bằng $1/4$ và triệt tiêu. Các mặt gãy là $x\pm2t=0,1$ và $y\pm2t=0,1$.

%ID VD_3_6_2_05
\textbf{Bước 1. Duhamel cho nguồn bật trong một đơn vị thời gian.} $c=2$, dữ liệu đầu không, $F(P,s)=\chi_{[0,1]}(s)\chi_{B_2(0)}(P)$. Xung tại thời điểm $s$ có nghiệm $K_2(r,t-s)$, $r=|P|$, nên
\[u=\int_0^{\min(t,1)}K_2(r,t-s)ds=J_2(r,t)-J_2(r,(t-1)_+).\]
\textbf{Bước 2. Tính nguyên hàm chỏm với tốc độ tổng quát.} Với $R,c>0$, đặt $b=|R-r|/c$, $d=(R+r)/c$, $q=[t]_b^d$ và
\[A(s)=\frac{(R^2-r^2)s+rcs^2-c^2s^3/3}{4rc}.\]
Ở phần giao chỏm, $K_R=[R^2-(r-cs)^2]/(4rc)=A'(s)$. Nếu $r<R$, trước $b$ có $K_R=s$; nếu $r\ge R$, trước $b$ bằng $0$. Do đó, với $r>0$,
\[J_R(r,t)=\tfrac12\min(t,(R-r)_+/c)^2+A(q)-A(b).\]
Với $r=0$, $J_R(0,t)=\frac12\min(t,R/c)^2$.

\textbf{Bước 3. Áp dụng $R=c=2$ và kiểm tra.} Thay hai bộ thời gian $t,(t-1)_+$ vào công thức; mọi nhánh đều là đa thức và min/max. Duhamel đảm bảo hai dữ liệu đầu không và PDE có đúng nguồn. Khi $t\ge1+(r+2)/2$, cả hai $J$ đã đạt cùng hằng số, nên $u=0$ tại điểm ấy: xung nguồn hữu hạn đi qua hoàn toàn trong ba chiều.

%ID BT_3_6_3_01
\textbf{Bước 1. Quy các miền về cầu, nửa không gian và lớp.} Với vận tốc chỉ thị, Kirchhoff bằng $t$ nhân tỷ lệ diện tích mặt cầu. Các quả cầu bị bỏ ở (a),(b),(c) nằm trong miền lớn và đôi cầu ở (a) không chồng thể tích, nên được lấy hiệu các chỉ thị. Các quả cầu ở (d) rời nhau, nên cộng ba đóng góp. Đặt $\mu$ hướng về tâm cầu; điều kiện $r^2+c^2t^2-2rct\mu<R^2$ cho ngưỡng chỏm.

\textbf{Bước 2. Tính các khoảng cách đúng điểm quan sát.} (a) $P=(0,100,0)$ cách tâm lớn $100$, cách hai tâm $(\pm1,0,0)$ cùng $\sqrt{10001}$. (b) $P_z=100$, cầu bị bỏ có tâm $z=-2$, cách $102$; nửa không gian $z<0$ cho $\mu<-100/(2t)$, tỷ lệ $(1-50/t)_+/2$. (c) lớp $0<z<4$ cho $-100/(3t)<\mu<-96/(3t)$, còn cầu tâm $z=2$ cách $98$. (d) khoảng cách là $98$ tới $(0,2,0)$ và $\sqrt{10004}$ tới hai tâm $(\pm2,0,0)$. (e) hiệu hai góc phần tư là $(\operatorname{sgn}x+\operatorname{sgn}y)/2$; mỗi tọa độ của hướng cầu đều phân bố đều trên $[-1,1]$.

\textbf{Bước 3. Thay tỷ lệ diện tích vào Kirchhoff.}

Kirchhoff với li độ đầu 0: $u(P,t)=t(4\pi)^{-1}\int_{S^2}\psi(P+ct\omega)d\omega$. Đối với một quả cầu bán kính $R$ có tâm cách $P$ một khoảng $d>0$, diện tích chỏm cho
\[K_R(d,t)=t\left[\frac{R^2-(d-ct)^2}{4dct}\right]_0^1,\quad [s]_0^1=\max(0,\min(1,s)).\]
Nếu $d=0$, $K_R(0,t)=t\chi_{ct<R}$. Giá trị trên mặt sóng hiểu theo giới hạn; dữ liệu bước nhảy không cho nghiệm cổ điển trên mặt ấy.

(a) $u=K_2(100,t)-2K_1(\sqrt{10001},t)$, $c=2$.
(b) $u=(t-50)_+/2-K_1(102,t)$, $c=2$; hạng đầu là tỷ lệ mặt cầu ở nửa không gian $z<0$.
(c) Với $c=3$, tỷ lệ của lớp $0<z<4$ cho
\[u=\frac t2\left[\left[-\frac{96}{3t}\right]_{-1}^1-\left[-\frac{100}{3t}\right]_{-1}^1\right]-K_1(98,t).\]
Ở đây $[s]_{-1}^1=\max(-1,\min(1,s))$.
(d) $u=K_1(98,t)+2K_1(\sqrt{10004},t)$, $c=3$.
(e) $\psi=(\operatorname{sgn}x+\operatorname{sgn}y)/2$. Tọa độ mỗi hướng cầu phân bố đều trên $[-1,1]$; ở $P=(100,100,0)$ trung bình sign là $\min(1,100/(3t))$. Do đó $u=\min(t,100/3)$.

\textbf{Bước 4. Đối chiếu và kết luận.} Ở (a),(c),(d), sau khi mặt cầu ra ngoài toàn bộ các miền bị chặn, đóng góp bằng $0$; ở (b),(e), miền vận tốc không bị chặn nên không suy như vậy. Tại các điểm dữ liệu liên tục, đạo hàm đầu đúng bằng chỉ thị; chuẩn hóa $\psi=1$ cho $u=t$. Các mặt tiếp xúc với biên miền là nơi cần xét vết một phía.

%ID BT_3_6_3_02
\textbf{Bước 1. Tham số hóa mặt cầu và giữ dấu của nón.} Tốc độ ký hiệu $c$, $\rho=ct$. Trên trục $z$, điểm cầu có chiều cao $z+\rho\mu$, bán kính trụ $\rho\sqrt{1-\mu^2}$. Miền $D$ là giao của $z'>\alpha r'$ với quả cầu $|Q|<R$; điều kiện $z'>0$ không được bỏ khi bình phương. Tích phân phương vị cho $2\pi$, nên $u=t|E|/2$.

\textbf{Bước 2. Giải từng ngưỡng; ở mặt phẳng $z=0$ khử góc phương vị.} Nhờ đối xứng chỉ cần $P=(r,0,0)$. Đặt $\omega=(\sqrt{1-\mu^2}\cos\phi,\sqrt{1-\mu^2}\sin\phi,\mu)$. Nón đòi $\mu>0$ và $r'^2<\rho^2\mu^2/\alpha^2$; cầu đòi $r'^2+\rho^2\mu^2<R^2$. Viết $r'^2=r^2+\rho^2(1-\mu^2)+2r\rho\sqrt{1-\mu^2}\cos\phi$, được hai cận trên cho $\cos\phi$. Độ dài góc đúng là $A(\min(B,C))$, không nhân hai tỷ lệ chỏm độc lập.

\textbf{Bước 3. Công thức tính $u$ và cực trị.}

(a) Đặt $\rho=ct$, $\mu=\cos\theta$. Tại $P=(0,0,z)$, miền hướng hợp lệ là
\[E=\{\mu\in[-1,1]:z+\rho\mu>\alpha\rho\sqrt{1-\mu^2},\quad z^2+\rho^2+2z\rho\mu<R^2\}.\]
Kirchhoff cho kết quả chính xác $u(0,0,z,t)=t|E|/2$. Để tính $|E|$ bằng hữu hạn phép toán: chia $[-1,1]$ ở các nghiệm thực của
\[(1+\alpha^2)\rho^2\mu^2+2z\rho\mu+z^2-\alpha^2\rho^2=0,\]
và điểm $(R^2-z^2-\rho^2)/(2z\rho)$ nếu $z\ne0$, cùng điểm $-z/\rho$. Trên từng khoảng kiểm tra hai bất đẳng thức gốc tại trung điểm rồi cộng độ dài các khoảng đúng. Kiểm tra dấu $z+\rho\mu$ loại nghiệm ngoại lai do bình phương. Khi $z=0$, $u=t(1-\alpha/\sqrt{1+\alpha^2})/2$ nếu $ct<R$, và 0 nếu $ct>R$.

(b) Do đối xứng quay, chỉ cần $r=\sqrt{x^2+y^2}$. Với $r>0$, đặt $D=2r\rho\sqrt{1-\mu^2}$,
\[B=\frac{R^2-r^2-\rho^2}{D},\qquad C=\frac{\rho^2\mu^2/\alpha^2-r^2-\rho^2(1-\mu^2)}D.\]
Gọi $A(v)=0$ nếu $v\le-1$, $A(v)=2\pi-2\arccos v$ nếu $|v|<1$, $A(v)=2\pi$ nếu $v\ge1$. Khi đó
\[U_r(t)=\frac t{4\pi}\int_0^1A(\min(B,C))d\mu.\]
Đây là tích phân một chiều đã loại hoàn toàn tích phân mặt. $U=0$ khi $ct>r+R$; để tìm giá trị lớn nhất, tính $U$ ở các điểm chuyển nhánh và các nghiệm $U'(t)=0$ trong $0<t<(r+R)/c$, lấy giá trị lớn nhất. Kèm mã tính tích phân và tối ưu cho $r,R,\alpha,c$ bất kỳ trong file Python chung. Không có một con số duy nhất vì đề để các tham số tùy ý. Với $r=0$, supremum là $R(1-\alpha/\sqrt{1+\alpha^2})/(2c)$; dữ liệu chỉ thị gián đoạn ở mặt cầu có thể khiến supremum chỉ đạt theo giới hạn trái.

\textbf{Bước 4. Đối chiếu và kết luận.} Với $r>0$, công thức tích phân và bài toán cực trị một biến là biểu diễn chính xác theo tham số; các con số do mã tích phân/tối ưu trả về là xấp xỉ. Mã chia lưới thời gian, tinh chỉnh từng cực đại và cho phép tăng độ phân giải để kiểm tra hội tụ, nên không dùng một lần tối ưu cục bộ làm bằng chứng cực đại toàn cục. Ở $r=0$ công thức supremum nêu trên là tường minh, suy bằng lấy $t\uparrow R/c$. Điều kiện nón strict và biên cầu phải được giữ khi xét đạt cực đại.

\textbf{Chi tiết nghiệm biên nón và độ dài phương vị.} Phương trình bậc hai ở (a) có biệt thức
\[4z^2\rho^2-4(1+\alpha^2)\rho^2(z^2-\alpha^2\rho^2)
=4\alpha^2\rho^2[(1+\alpha^2)\rho^2-z^2].\]
Nếu biểu thức dưới căn không âm và $\rho>0$, hai nghiệm là
\[\mu_\pm=\frac{-z\pm\alpha\sqrt{(1+\alpha^2)\rho^2-z^2}}{(1+\alpha^2)\rho}.\]
Giữ những nghiệm thuộc $[-1,1]$, thêm các điểm cắt cầu và dấu đã nêu, rồi kiểm tra bất đẳng thức gốc trên từng khoảng. Ở (b), $\cos\phi<v$ với $-1<v<1$ cho $\arccos v<\phi<2\pi-\arccos v$; độ dài là $2\pi-2\arccos v$. Hai điều kiện cùng là cận trên nên lấy $\min(B,C)$. Khi $r=0$, nón cho $\mu>\alpha/\sqrt{1+\alpha^2}$ và cầu cho $\rho<R$, nên tỷ lệ diện tích bằng $(1-\alpha/\sqrt{1+\alpha^2})/2$.

%ID BT_3_6_3_03
\textbf{Bước 1. Đưa điều kiện xuyến về một biến góc.} Trên trục $z$, bán kính trụ của điểm cầu là $\rho\sqrt{1-\mu^2}$, chiều cao $z+\rho\mu$. Thay vào $(r'-R)^2+z'^2<a^2$, khai triển để được tập $E$ dưới đây. Bình phương phương trình biên có thể sinh nghiệm ngoại lai, nên chia các khoảng và kiểm tra bất đẳng thức gốc.

\textbf{Bước 2. Ở mặt phẳng $z=0$, tính góc phương vị.} Chọn $P=(r,0,0)$, đặt $h=\sqrt{a^2-\rho^2\mu^2}$. Điều kiện xuyến tương đương $R-h<r'<R+h$, với $|\rho\mu|<a$. Vì $R>a$, cận dưới dương nên có thể bình phương hai vế. Thay $r'^2=r^2+\rho^2(1-\mu^2)+2r\rho\sqrt{1-\mu^2}\cos\phi$ cho $C_-<\cos\phi<C_+$. Độ dài góc là hiệu $A(C_+)-A(C_-)$.

\textbf{Bước 3. Tích phân một chiều và cực trị.}

(a) Đặt $\rho=ct$. Tại $(0,0,z)$, điều kiện nằm trong xuyến là
\[E=\{\mu\in[-1,1]:\rho^2+R^2+z^2+2z\rho\mu-2R\rho\sqrt{1-\mu^2}<a^2\}.\]
Do đó $u=t|E|/2$. Tính các đầu mút bằng nghiệm của
\[(\rho^2+R^2+z^2-a^2+2z\rho\mu)^2=4R^2\rho^2(1-\mu^2),\]
kiểm tra bất đẳng thức gốc ở từng khoảng để loại nghiệm ngoại lai. Đây là phương trình bậc hai theo $\mu$, nên kết quả hữu hạn phép toán. Với $z=0$, miền đúng là $|\mu|<\sqrt{1-h^2}$ khi $h=(\rho^2+R^2-a^2)/(2R\rho)<1$, nên $u=t\sqrt{1-h^2}$; ngoài $R-a<\rho<R+a$, $u=0$.

(b) Với $r=\sqrt{x^2+y^2}>0$, $|\mu|<\min(1,a/\rho)$, đặt $h=\sqrt{a^2-\rho^2\mu^2}$, $D=2r\rho\sqrt{1-\mu^2}$,
\[C_\pm=\frac{(R\pm h)^2-r^2-\rho^2(1-\mu^2)}D.\]
Độ dài góc phương vị là $A(C_+)-A(C_-)$ với $A$ định nghĩa ở bài nón. Vậy
\[U_r(t)=\frac t{4\pi}\int_{-\min(1,a/\rho)}^{\min(1,a/\rho)}[A(C_+)-A(C_-)]d\mu.\]
Chỉ cần xét $0<t<(r+R+a)/c$; tính các cực trị $U'=0$ và điểm chuyển nhánh, lấy lớn nhất. Mã Python kèm theo thực hiện bước tích phân/tối ưu với tham số tùy chọn. Riêng $r=0$, từ câu (a), $U^2=[4R^2\rho^2-(\rho^2+R^2-a^2)^2]/(4R^2c^2)$; cực đại tại $\rho=\sqrt{R^2+a^2}$, cho $\max U=a/c$.

\textbf{Bước 4. Đối chiếu và kết luận.} Ở $r=0$, $U^2$ là một đa thức bậc hai theo $\rho^2$ với hệ số đầu âm. Đạo hàm bằng $0$ tại $\rho^2=R^2+a^2$, nằm giữa $(R-a)^2,(R+a)^2$, nên đây là cực đại toàn cục $a/c$. Ở $r>0$, biểu thức tích phân và cực trị một biến là chính xác theo tham số; kết quả từ mã là xấp xỉ, cần tăng cả độ chính xác tích phân và số khoảng thời gian để kiểm tra ổn định.

\textbf{Chi tiết nghiệm bậc hai tại điểm $(0,0,z)$.} Đặt $A=\rho^2+R^2+z^2-a^2$. Phương trình biên sau bình phương là
\[4\rho^2(z^2+R^2)\mu^2+4Az\rho\mu+A^2-4R^2\rho^2=0.\]
Với $\rho>0$, hai nghiệm thực nếu biểu thức dưới căn không âm:
\[\mu_\pm=\frac{-Az\pm R\sqrt{4\rho^2(z^2+R^2)-A^2}}{2\rho(z^2+R^2)}.\]
Giữ các nghiệm trong $[-1,1]$ và kiểm tra $A+2z\rho\mu=2R\rho\sqrt{1-\mu^2}$ để loại nghiệm ngoại lai. Chia khoảng ở các nghiệm còn lại rồi kiểm tra bất đẳng thức gốc trên từng khoảng để cộng độ dài; tại $t=0$ dùng dữ liệu $u=0$.

\textbf{Chi tiết khai triển biên và cực đại trên trục.} Ở (a), $(\rho\sqrt{1-\mu^2}-R)^2+(z+\rho\mu)^2$ khai triển thành $\rho^2+R^2+z^2+2z\rho\mu-2R\rho\sqrt{1-\mu^2}$: hai hạng $\rho^2\mu^2$ triệt tiêu. Với $z=r=0$, đặt $q=\rho^2$,
\[4R^2c^2U^2=4R^2q-(q+R^2-a^2)^2.\]
Đạo hàm theo $q$ bằng $4R^2-2(q+R^2-a^2)=2(R^2+a^2-q)$, đạo hàm bậc hai bằng $-2$. Cực đại duy nhất là $q=R^2+a^2$; thay vào được $4R^2c^2U^2=4R^2a^2$, tức $U=a/c$. $q-(R-a)^2=2Ra>0$ và $(R+a)^2-q=2Ra>0$, nên điểm cực đại nằm trong miền có sóng. Ở (b), $C_-<\cos\phi<C_+$ cho độ dài phương vị bằng hiệu hai độ dài $A(C_+)-A(C_-)$ sau khi cắt cận.

%ID BT_3_6_3_04
\textbf{Bước 1. Kiểm tra hai dữ liệu trước khi dùng công thức in.} Tại $t=0$, công thức $\psi+t\varphi$ trong sách cho li độ $\psi$ và vận tốc $\varphi$, đảo với dữ liệu đã ghi. Vì thế cần đổi lại $\varphi+t\psi$. Đặt $I(t)=\int_0^t(t-s)g(s)ds$, ta có $I(0)=I'(0)=0$, $I''=g$.

\textbf{Bước 2. Trường hợp điều hòa.} Mọi Laplace của $f,\varphi,\psi$ đều bằng không, nên nghiệm không có đóng góp $c^2\Delta u$; chỉ cần giải ODE thời gian tại mỗi điểm. Trường hợp $k$-điều hòa dùng chuỗi Taylor chẵn/lẻ theo thời gian và chuỗi Duhamel, nhưng kết thúc vì $\Delta^k$ của từng dữ liệu bằng không.

\textbf{Bước 3. Nghiệm sửa đúng và chuỗi hữu hạn.}

(a) Với dữ liệu đầu $u(0)=\varphi,u_t(0)=\psi$, nghiệm đúng là
\[u=\varphi+t\psi+f\int_0^t(t-s)g(s)ds.\]
Do các hàm không gian điều hòa, $\Delta u=0$ và $u_{tt}=fg$. Dòng công thức trong đề đảo $\varphi,\psi$ nên không thỏa dữ liệu đã ghi; cần sửa như trên.
(b) Với $\Delta^kf=\Delta^k\varphi=\Delta^k\psi=0$, chuỗi hữu hạn:
\begin{align*}
u=\sum_{j=0}^{k-1}c^{2j}\bigg[&\frac{t^{2j}}{(2j)!}\Delta^j\varphi+\frac{t^{2j+1}}{(2j+1)!}\Delta^j\psi\\
&+\Delta^jf\int_0^t\frac{(t-s)^{2j+1}}{(2j+1)!}g(s)ds\bigg].
\end{align*}
Đạo hàm hai lần làm dịch chỉ số; hạng cuối hết vì $\Delta^k=0$, chứng minh phương trình và dữ liệu.

\textbf{Bước 4. Đối chiếu và kết luận.} Để kiểm tra chuỗi ở (b), đạo hàm hai lần hạng li độ chuyển $t^{2j}/(2j)!$ thành $t^{2j-2}/(2j-2)!$; Laplace hạng trước sinh cùng $\Delta^j\varphi$. Tương tự cho vận tốc và tích phân nguồn. Chỉ hạng nguồn $j=0$ sinh $fg$, các hạng cuối mất vì $\Delta^k=0$. Tại $0$ chỉ hai hạng $j=0$ có thể cho li độ/vận tốc, nên đúng $\varphi,\psi$.

%ID BT_3_6_3_05
\textbf{Bước 1. Duhamel của nguồn cầu không đổi theo thời gian.} Li độ/vận tốc đầu bằng không, tốc độ $1$. Tại mỗi thời điểm nguồn $s$, vận tốc xung là $\chi_{B_1}$, nên đóng góp ở thời gian trễ $\tau=t-s$ là $K_1(r,\tau)$. Tích phân $s$ cho $u=J_1(r,t)$.

\textbf{Bước 2. Tính nguyên hàm ở từng giai đoạn.} Với cầu bán kính $R$, nếu $r<R$, trước $R-r$ có $K_R(s)=s$; nếu $r\ge R$, trước $r-R$ có $K_R=0$. Trong $|R-r|<s<R+r$, $K_R=(R^2-(r-s)^2)/(4r)$ là đa thức bậc hai. Sau $R+r$ đóng góp xung bằng không, nên $J_R$ giữ hằng số. Tích phân ba nhánh cho công thức dưới đây.

\textbf{Bước 3. Nghiệm và trạng thái dừng.}

Đặt $r=|P|$, $J_R(r,t)=\int_0^tK_R(r,s)ds$ với $c=1$. Duhamel cho $u=J_1(r,t)$. Khi $r>0$, đặt $b=|R-r|$, $q=\min(\max(t,b),R+r)$,
\[A(s)=\frac{(R^2-r^2)s+rs^2-s^3/3}{4r},\quad J_R=\frac12\min(t,(R-r)_+)^2+A(q)-A(b).\]
Với $r=0$, $J_R(0,t)=\min(t,R)^2/2$. Khi $t\ge r+R$, nghiệm bằng thế Newton của quả cầu: $R^3/(3r)$ nếu $r\ge R$, $(3R^2-r^2)/6$ nếu $r\le R$. Kiểm tra $\Delta u=-\chi_{r<R}$ cho trạng thái dừng.

\textbf{Bước 4. Đối chiếu và kết luận.} Với $R=1$, trạng thái dừng là $(3-r^2)/6$ bên trong và $1/(3r)$ bên ngoài. Công thức Laplace hướng tâm $v_{rr}+2v_r/r$ cho $-1$ bên trong và $0$ bên ngoài; giá trị cùng đạo hàm nối tại $r=1$, nên không sinh nguồn mặt cầu. Hai dữ liệu đầu của $J$ bằng $0$ vì $J(0)=0$ và $J_t=K$, $K(r,0)=0$.

\textbf{Chi tiết nguyên hàm và trạng thái dừng.} Trong nhánh chỏm,
\[K_R(r,s)=\frac{R^2-r^2+2rs-s^2}{4r},\quad
\int K_Rds=\frac{(R^2-r^2)s+rs^2-s^3/3}{4r}=A(s).\]
Nếu $r\ge R$, trước $b=r-R$ tích phân bằng $0$, và
\[J_\infty=A(r+R)-A(r-R)
=\frac{2R(R^2-r^2)+4r^2R-(2r^2R+2R^3/3)}{4r}
=\frac{R^3}{3r}.\]
Nếu $r<R$, trước $b=R-r$ tích phân bằng $(R-r)^2/2$. Tính hiệu:
\[A(R+r)-A(R-r)=Rr-\frac23r^2,\quad
J_\infty=\frac{(R-r)^2}{2}+Rr-\frac23r^2=\frac{3R^2-r^2}{6}.\]
Tại $r=R$, hai giá trị bằng $R^2/3$; đạo hàm trong $-r/3$ và ngoài $-R^3/(3r^2)$ cùng bằng $-R/3$.

%ID BT_3_6_3_06
\textbf{Bước 1. Chồng nghiệm vận tốc và nguồn.} Hai cầu vận tốc có dấu trái nhau cho $K_1(d_+,t)-K_1(d_-,t)$. Nguồn chỉ bật khi $0<s<T$, nên phần nguồn là $\int_0^{\min(t,T)}K_1(r,t-s)ds$. Đổi $\tau=t-s$ cho hiệu hai nguyên hàm $J$.

\textbf{Bước 2. Dùng đối xứng và thời gian yêu cầu.} Ở $x=0$, khoảng cách tới hai tâm $(\pm2,0,0)$ bằng nhau, nên phần vận tốc triệt tiêu hoàn toàn; câu (a) chỉ còn nguồn với $T=1$. Câu (b) giữ hai khoảng cách khác nhau và thay $t=r$ vào biểu thức, không thay $d_\pm$ bằng $r$.

\textbf{Bước 3. Công thức áp dụng.}

Gọi $d_\pm=|P-(\pm2,0,0)|$, $r=|P|$. Chồng Kirchhoff và Duhamel cho
\[u=K_1(d_+,t)-K_1(d_-,t)+J_1(r,t)-J_1(r,(t-T)_+),\]
với $c=1$, $K,J$ có biểu thức chỏm và nguyên hàm ở hai bài trước.
(a) $P=(0,y,z)$ cho $d_+=d_-$, nên $u=J_1(\sqrt{y^2+z^2},t)-J_1(\sqrt{y^2+z^2},(t-1)_+)$.
(b) Thay $t=r$: $u=K_1(d_+,r)-K_1(d_-,r)+J_1(r,r)-J_1(r,(r-T)_+)$. Đây là biểu thức tường minh bằng min/max và đa thức, áp dụng mọi $P$ kể cả $r=0$ theo giới hạn.

\textbf{Bước 4. Đối chiếu và kết luận.} Tại $t=0$, cả $J$ và $K$ bằng $0$; đạo hàm của phần $K$ cho đúng hiệu hai chỉ thị cầu và phần $J$ có vận tốc đầu không. Với $t=r=0$, kết quả là $0$ và dùng nhánh tâm của $J$ để tránh chia $r$. Các min/max của $K,J$ nêu trong phần kiến thức và Ví dụ 3.6.2.5 làm biểu thức câu (b) hoàn toàn tường minh.

\textbf{Chi tiết các hàm dùng trong biểu thức tường minh.} Với tốc độ $1$, bán kính $1$, $d>0,t>0$,
\[K_1(d,t)=t\left[\frac{1-(d-t)^2}{4dt}\right]_0^1.\]
Với $r>0$, đặt $b=|1-r|$, $q=[t]_b^{r+1}$ và $A(s)=((1-r^2)s+rs^2-s^3/3)/(4r)$:
\[J_1(r,t)=\tfrac12\min(t,(1-r)_+)^2+A(q)-A(b).\]
Tại tâm, $K_1(0,t)=t$ khi $t<1$, bằng $0$ khi $t>1$, $J_1(0,t)=\min(t,1)^2/2$; tại $t=0$, $K_1=J_1=0$. Khi đổi $\tau=t-s$, $s$ chạy $0$ đến $\min(t,T)$ nên $\tau$ chạy ngược từ $t$ đến $(t-T)_+$:
\[\int_0^{\min(t,T)}K_1(r,t-s)ds=\int_{(t-T)_+}^{t}K_1(r,\tau)d\tau=J_1(r,t)-J_1(r,(t-T)_+).\]

%ID BT_3_6_3_07
\textbf{Bước 1. Tích phân Poisson theo tung độ trước.} Với $c=2$, tại hoành độ $X$, đĩa ảnh hưởng có tung độ $P_y\pm h(X)$. Một dải $l<Y<v$ góp nguyên hàm $\arcsin((Y-P_y)/h)$, sau khi cắt hai cận vào $[-1,1]$. Các miền trong đề có cận tung độ tường minh theo $|X|$, nên chỉ còn một tích phân theo $X$.

\textbf{Bước 2. Hai đối xứng cần dùng.} Câu (a): trước $2t<98/\sqrt5$, cả đĩa quanh $(0,100)$ nằm trong miền $D_1$, vận tốc bằng $2$, nên $u=2t$. Câu (b): $\psi(Y,X)=-\psi(X,Y)$ và điểm quan sát $(100,100)$ bất biến khi đổi tọa độ. Đĩa và nhân Poisson cũng bất biến, nên tích phân đổi thành đối của chính nó: $u(100,100,t)=0$ với mọi $t$. Công thức tích phân dưới đây áp dụng cả cho các điểm khác khi vẽ đồ thị, nhưng tại điểm của câu (b) kết quả rút gọn là $0$.

\textbf{Bước 3. Tính các dải và dựng đồ thị.}

Poisson hai chiều: $u(P,t)=(2\pi c)^{-1}\int_{|Q-P|<ct}\psi(Q)/\sqrt{c^2t^2-|Q-P|^2}\,dQ$. Tích phân theo tung độ được thực hiện chính xác. Đặt $h(X)=\sqrt{c^2t^2-(X-P_x)^2}$,
\[B_X(l,v)=\arcsin\left[\frac{v-P_y}{h(X)}\right]_{-1}^1-\arcsin\left[\frac{l-P_y}{h(X)}\right]_{-1}^1.\]
(a) Với $P=(0,100),c=2$,
\[u=\frac1{4\pi}\int_{-2t}^{2t}\{2B_X(2+2|X|,\infty)-B_X(-\infty,\min(2-2|X|,0))+B_X(\max(2-2|X|,0),2+2|X|)\}dX.\]
(b) Với $P=(100,100),c=2$, tích phân $X\in[100-2t,100+2t]$: nếu $|X|>1$, tích phân tung độ là $B_X(-1,1)$; nếu $|X|<1$, là $-B_X(-\infty,-1)-B_X(1,\infty)$. Nhân $1/(4\pi)$. Đây là tích phân một chiều của hàm sơ cấp; không có kỳ dị căn ở tung độ nữa. Ký hiệu $D_1$ thứ hai trong đề là $D_2$.
(c) Đồ thị kèm theo được tính từ tích phân Poisson; mã cho phép đổi điểm/thời gian. Dữ liệu chỉ thị nên hiểu nghiệm yếu và sự gãy ở biên vùng ảnh hưởng.

\textbf{Bước 4. Đối chiếu và kết luận.} Với vận tốc hằng $1$, tích phân hai chiều của nhân bằng $2\pi ct$, cho $u=t$, kiểm tra hệ số $1/(4\pi)$ khi $c=2$. Đồ thị bề mặt dùng tích phân số, phải đọc tham số lưới trong mã; đường $x=y$ của câu (b) luôn bằng $0$ theo đối xứng, là một kiểm tra độc lập cho đồ thị.

\textbf{Chi tiết loại tích phân tung độ và kiểm tra hệ số.} Tại $X$ cố định, đặt $h^2=c^2t^2-(X-P_x)^2$ và $z=(Y-P_y)/h$. Khi $h>0$,
\[\int\frac{dY}{\sqrt{h^2-(Y-P_y)^2}}
=\int\frac{dz}{\sqrt{1-z^2}}=\arcsin z.\]
Chỉ lấy giao của dải với $P_y-h<Y<P_y+h$, tương đương cắt $z$ vào $[-1,1]$, cho $B_X(l,v)$. Với vận tốc hằng $1$, dùng tọa độ cực quanh $P$:
\[\int_{|Q-P|<ct}\frac{dQ}{\sqrt{c^2t^2-|Q-P|^2}}
=2\pi\int_0^{ct}\frac{r\,dr}{\sqrt{c^2t^2-r^2}}
=2\pi ct.\]
Nhân $1/(2\pi c)$ cho $u=t$. Ở (a), khoảng cách từ $(0,100)$ đến hai đường $Y=2\pm2X$ là $|100-2|/\sqrt{1+4}=98/\sqrt5$, nên trước $2t<98/\sqrt5$ cả đĩa có vận tốc $2$.

%ID BT_3_6_3_08
\textbf{Bước 1. Hạ về hai bài một chiều có nguồn.} Vì $F=\chi_{[0,1]}(x)+\chi_{[0,1]}(y)$, đặt $u=U(x,t)+U(y,t)$. Mỗi $U$ có dữ liệu không, giải $U_{tt}-U_{xx}=\chi_{[0,1]}$.

\textbf{Bước 2. Tích phân hai lần hàm chỉ thị.} $H(s)=s_+-(s-1)_+$ là nguyên hàm liên tục của chỉ thị, và $Q(s)=(s_+^2-(s-1)_+^2)/2$ thỏa $Q''=\chi_{(0,1)}$. Duhamel cho $U=\frac12\int_0^t[H(x+\tau)-H(x-\tau)]d\tau=\frac12[Q(x+t)+Q(x-t)-2Q(x)]$. Khai triển $Q$ cho công thức dưới đây.

\textbf{Bước 3. Nghiệm và kiểm tra.}

Nguồn tách thành hai hàm chỉ phụ thuộc một tọa độ nên $u=U(x,t)+U(y,t)$. Duhamel một chiều với nguồn $\chi_{[0,1]}$ cho
\[U(x,t)=\frac14\{(x+t)_+^2-2x_+^2+(x-t)_+^2-(x-1+t)_+^2+2(x-1)_+^2-(x-1-t)_+^2\}.\]
Đạo hàm hai lần theo $t,x$ cho $U_{tt}-U_{xx}=\chi_{(0,1)}(x)$; $U(0)=U_t(0)=0$ theo thời gian. Do đó tổng đúng phương trình hai chiều.

\textbf{Bước 4. Đối chiếu và kết luận.} Từ $U=\frac12(Q(x+t)+Q(x-t)-2Q(x))$, có $U_{tt}-U_{xx}=Q''(x)=\chi_{(0,1)}(x)$. Tại $t=0$ các hạng triệt tiêu và $U_t=0$. Vì tổng không có đạo hàm chéo giữa $x,y$, PDE hai chiều có đúng hai nguồn. Công thức đúng theo nghĩa phân phối cả tại các đường cận chạm $0,1$.

%ID BT_3_6_3_09
\textbf{Bước 1. Đạo hàm tích và Euler.} $q=x^2+y^2-t^2$, $\Box=\Delta-\partial_{tt}$, nên $\Box q=6$. Với $H$ thuần nhất bậc $m$, hạng đạo hàm chéo bằng $4j q^{j-1}(xH_x+yH_y+tH_t)=4jm q^{j-1}H$. Cộng $\Box q^j=2j(2j+1)q^{j-1}$ cho công thức tích dưới đây.

\textbf{Bước 2. Lập hệ số thay vì chép dấu in.} Đặt $u=\sum a_kq^{k+1}\Box^k\Phi$. Hệ số hạng hạ bậc là $C_k=2(k+1)(2d-2k-1)$. Vì PDE là $-\Box u=\Phi$, cần $a_0=-1/C_0$; các hạng sau triệt tiêu khi $a_k=-a_{k-1}/C_k$. Suy biểu thức giai thừa kép, có dấu âm phía ngoài và số mũ $q^{k+1}$, không phải $q^{2k+2}$.

\textbf{Bước 3. Công thức đúng, chứng minh duy nhất và nguồn cụ thể.}

Đặt $q=x^2+y^2-t^2$, $\Box=\Delta_{x,y}-\partial_{tt}$; phương trình là $-\Box u=\Phi$. Với $H$ thuần nhất bậc $m$,
\[\Box(q^jH)=q^j\Box H+2j(2m+2j+1)q^{j-1}H.\]
Lấy $H=\Box^k\Phi$ bậc $d-2-2k$, $j=k+1$: hệ số là $2(k+1)(2d-2k-1)$. Quy nạp triệt tiêu các hạng cho
\[u=-\sum_{k\ge0}\frac{(-1)^k(2d-2k-3)!!}{(2k+2)!!(2d-1)!!}\,q^{k+1}\Box^k\Phi.\]
Tổng hữu hạn vì $\Phi$ đa thức. Đề có dấu dương và chỉ số giai thừa kép không nhất quán; công thức trên thỏa $u_{tt}-\Delta u=\Phi$ và $u=0$ trên nón.

(ii-a,b) Với hiệu hai nghiệm nguồn và vết nón 0, đặt năng lượng trên $B_{-t}$, $t<0$. Đạo hàm miền di động cho $E'=\int_{r=-t}[u_tu_r-(u_t^2+|\nabla u|^2)/2]$. Trên nón, vết 0 buộc đạo hàm tiếp tuyến $u_t=u_r$ và đạo hàm góc 0, nên thông lượng bằng 0. Nếu nghiệm đủ trơn đến đỉnh, $E(t)\to0$ khi $t\uparrow0$, vì diện tích miền co về 0. Vậy $E=0$, $u_t=\nabla u=0$ trên mọi mặt $t=-h$, suy $u=0$ và duy nhất trong nón. Điều kiện chính quy ở đỉnh cần thiết cho lập luận.

(ii-c) $\Phi=xyt$ có bậc 3 nên $d=5$, $\Box\Phi=0$:
\[u=-\frac{(x^2+y^2-t^2)xyt}{18}.\]
Kiểm tra $u_{tt}-u_{xx}-u_{yy}=xyt$ và vết nón 0, nên đây là nghiệm duy nhất trong lớp chính quy đã nêu.

\textbf{Bước 4. Đối chiếu và kết luận.} Đạo hàm biên chuyển động trong năng lượng ở (ii) có tốc độ bán kính $-1$, nên hạng biên mang dấu âm. Vết $u(r,-r)=0$ cho $u_r-u_t=0$; thế vào thông lượng triệt tiêu bình phương. Với $\Phi=xyt$, tính trực tiếp $u=-[x^3yt+xy^3t-xyt^3]/18$: $u_{tt}=xyt/3$, $u_{xx}=-xyt/3$, $u_{yy}=-xyt/3$, suy hiệu bằng $xyt$. Hệ số $-1/18$ và dấu đã được kiểm tra độc lập.

\textbf{Chi tiết đạo hàm tích và truy hồi hệ số.} $\nabla_{x,y}q=(2x,2y)$, $q_t=-2t$, nên $|\nabla q|^2-q_t^2=4q$ và $\Box q=6$. Quy tắc dây chuyền:
\[\Box q^j=j q^{j-1}\Box q+j(j-1)q^{j-2}(|\nabla q|^2-q_t^2)
=[6j+4j(j-1)]q^{j-1}=2j(2j+1)q^{j-1}.\]
Hạng chéo trong $\Box(q^jH)$ là
\[2(\nabla q^j\cdot\nabla H-(q^j)_tH_t)
=4jq^{j-1}(xH_x+yH_y+tH_t)=4jmq^{j-1}H.\]
Gọi $H_k=\Box^k\Phi$, $C_k=2(k+1)(2d-2k-1)$. Thế chuỗi:
\[\Box u=a_0C_0\Phi+\sum_{k\ge1}(a_{k-1}+C_ka_k)q^kH_k,\]
vì hạng cuối $H_{k+1}$ hết khi bậc đa thức âm. Muốn $\Box u=-\Phi$ cần $a_0C_0=-1$, $a_{k-1}+C_ka_k=0$. Do đó
\[a_k=\frac{(-1)^{k+1}}{C_0\cdots C_k}
=\frac{(-1)^{k+1}(2d-2k-3)!!}{(2k+2)!!(2d-1)!!},\]
vì tích các thừa số chẵn là $(2k+2)!!$, tích lẻ là $(2d-1)!!/(2d-2k-3)!!$.
